% Le Cameroun indépendant — Histoire Tle Tech — Cours signé OnBuch+
% Programme officiel MINESEC. Compiler avec : tectonic N07-cameroun-independant.tex
\newcommand{\DOCMATIERE}{Histoire}
\newcommand{\DOCNIVEAU}{Tle Tech}
\newcommand{\DOCMODULE}{Histoire – classe de Terminale technique}
\newcommand{\DOCLECON}{N07}
\newcommand{\DOCTITRE}{Le Cameroun indépendant}
% OnBuch+ — préambule « cours signés » ENRICHI (compilé avec tectonic / XeLaTeX)
% Système visuel « Pop » d'OnBuch+ : fond crème (jamais blanc), bordures encre
% épaisses, ombres « dures » décalées, titres Archivo Black en capitales.
% Version MATHÉMATIQUES : graphes (pgfplots/tikz), boîtes pédagogiques
% (définition, propriété, méthode, attention, le savais-tu, exercice résolu,
% exercice, TP, bilan), page de garde et signature OnBuch+ ancrée en bas.
%
% Macros attendues dans main.tex AVANT \input{preamble} :
%   \DOCMATIERE  \DOCNIVEAU  \DOCMODULE  \DOCLECON (numéro)  \DOCTITRE
\documentclass[11pt]{article}

\usepackage{fontspec}
\usepackage[french]{babel}
\usepackage[a4paper,margin=2cm,top=3.4cm,bottom=2.8cm,headheight=1.4cm,footskip=1.3cm]{geometry}
\usepackage{amsmath,amssymb,amsfonts}
\usepackage{mathtools} % \xmapsto, \coloneqq... souvent utilisés par les modèles
\usepackage{cancel} % simplifications barrées (bilans de forces)
% \abs{z} (module, valeur absolue) et \norm{u} (norme), utilisés par les modèles
\providecommand{\abs}{}\let\abs\relax\DeclarePairedDelimiter{\abs}{\lvert}{\rvert}
\providecommand{\norm}{}\let\norm\relax\DeclarePairedDelimiter{\norm}{\lVert}{\rVert}
\usepackage[table]{xcolor} % [table] : fournit aussi \rowcolors (alternance de lignes)
\usepackage{pagecolor}
\usepackage{tikz}
\usetikzlibrary{arrows.meta,positioning,calc,shapes.geometric,shapes.misc,shapes.arrows,decorations.pathmorphing,decorations.pathreplacing,decorations.markings,patterns,shadows,mindmap,trees,fit,backgrounds,matrix,intersections,angles,quotes}
\usepackage{pgfplots}
\usepackage[version=4]{mhchem} % équations nucléaires \ce{^{235}_{92}U}
\usepackage[european,siunitx]{circuitikz} % schémas électriques
\pgfplotsset{compat=1.18}
\usepgfplotslibrary{fillbetween,groupplots} % aires entre courbes (calcul intégral)
\usepackage{enumitem}
\usepackage{fancyhdr}
% [most] charge toutes les bibliothèques tcolorbox (listings, minted, raster,
% documentation...) et gonfle inutilement la mémoire TeX sur un document long
% (~300 boîtes) : on ne charge que ce qui est réellement utilisé.
\usepackage[skins,breakable]{tcolorbox}
\usepackage{graphicx}
\usepackage{colortbl}
\usepackage{booktabs}
\usepackage{tabularx}
\usepackage{multirow}
\usepackage{diagbox} % case d'en-tête diagonale des tableaux à double entrée
% Colonnes extensibles centrée / gauche / droite (C, L, R), souvent utilisées par les modèles
\newcolumntype{C}{>{\centering\arraybackslash}X}
\newcolumntype{L}{>{\raggedright\arraybackslash}X}
\newcolumntype{R}{>{\raggedleft\arraybackslash}X}
\usepackage{array}
\usepackage{multicol}
\usepackage{siunitx}
% parse-numbers=false : accepte aussi \SI{2,5 \times 10^{-3}}{...}
\sisetup{locale=FR,output-decimal-marker={,},group-minimum-digits=4,parse-numbers=false}
\DeclareSIUnit\ohm{\ensuremath{\symup{\Omega}}} % U+2126 absent de la police
\DeclareSIUnit\division{div} % réglages d'oscilloscope (ms/div, V/div)
\DeclareSIUnit\tr{tr} % tours (vitesse de rotation, ex. \SI{1500}{\tr\per\minute})
\DeclareSIUnit\molar{M} % concentration molaire (ex. \SI{3}{\molar}, \SI{1}{\micro\molar})
\DeclareSIUnit\an{an} % année (le modèle confond parfois avec \year, un compteur TeX) — ex. \SI{5}{\kilo\meter\per\an}
\DeclareSIUnit\FCFA{FCFA} % franc CFA (ex. \SI{500}{\FCFA\per\kilo\gram})
\DeclareSIUnit\degreeBrix{\textdegree Brix} % teneur en sucre d'un sirop/jus (réfractométrie)
\DeclareSIUnit\month{mois} % le modèle utilise parfois \month, un compteur TeX, comme unité de durée
\DeclareSIUnit\microliter{\micro\liter} % le modèle écrit parfois \microliter en un seul token
\DeclareSIUnit\million{millions} % dénombrement (ex. \SI{15}{\million\per\mL} spermatozoïdes)
\usepackage{qrcode}
% \degree : commande fréquemment utilisée par les modèles pour un angle ou une
% température en dehors de \SI{}{} (non fournie par les paquets chargés).
\providecommand{\degree}{^{\circ}}
% \qed (fin de démonstration), utilisé par les modèles sans amsthm
\providecommand{\qed}{\hfill$\square$}
% \barre : double barre des valeurs interdites dans un tableau de variations (tkz-tab)
\providecommand{\barre}{\ensuremath{\|}}
% environnement proof (amsthm non chargé) utilisé par les modèles
\ifdefined\proof\else\newenvironment{proof}[1][Démonstration]{\par\noindent\textit{#1.}\ }{\qed\par}\fi
\usepackage{lastpage}
\usepackage{hyperref}
\hypersetup{hidelinks}

% ============================================================
% Palette « Pop » OnBuch+ — mêmes hex que la classe P (plus_ui.dart)
% ============================================================
\definecolor{poporange}{HTML}{F59321}
\definecolor{popdark}{HTML}{E07A0C}
\definecolor{popcream}{HTML}{FFF6E8}
\definecolor{popink}{HTML}{1C1714}
\definecolor{poppurple}{HTML}{7A5AE0}
\definecolor{popgreen}{HTML}{2E9E6A}
\definecolor{poppink}{HTML}{E0457B}
\definecolor{popblue}{HTML}{2F7DE1}
\definecolor{popgold}{HTML}{C98A12}
\definecolor{popmuted}{HTML}{8A7F76}
\definecolor{popline}{HTML}{EADFD2}
\definecolor{popgray}{HTML}{8A7F76}
\colorlet{popgrey}{popgray}
% Alias tolérants (le modèle nomme parfois les couleurs différemment)
\colorlet{popwhite}{white}
\colorlet{popblack}{popink}
\colorlet{popred}{poppink}
\colorlet{popyellow}{popgold}
% Teintes claires (fonds de tableaux/encadrés) — le modèle nomme parfois
% les couleurs avec un suffixe L (ex. popblueL) sans les avoir définies.
\colorlet{poporangeL}{poporange!20}
\colorlet{popdarkL}{popdark!20}
\colorlet{poppurpleL}{poppurple!20}
\colorlet{popgreenL}{popgreen!20}
\colorlet{poppinkL}{poppink!20}
\colorlet{popblueL}{popblue!20}
\colorlet{popgoldL}{popgold!20}
\colorlet{popredL}{popred!20}
\colorlet{popgrayL}{popgray!20}
% Teintes claires pour fonds de tableaux / zones
\colorlet{popblueL}{popblue!10!white}
\colorlet{poporangeL}{poporange!14!white}
\colorlet{popgreenL}{popgreen!12!white}
\colorlet{poppurpleL}{poppurple!12!white}
\colorlet{poppinkL}{poppink!12!white}
\colorlet{popgoldL}{popgold!14!white}

\pagecolor{popcream}

% ============================================================
% Typographie
% ============================================================
\defaultfontfeatures{Ligatures=TeX}
\setmainfont{PlusJakartaSans-Medium.ttf}[Path=../../../fonts/,BoldFont=PlusJakartaSans-Bold.ttf]
% Maths en sans-serif (Fira Math) avec chiffres et lettres droites de Plus
% Jakarta, pour que les formules se fondent dans le texte.
\usepackage[math-style=ISO,bold-style=ISO]{unicode-math}
\setmathfont{FiraMath-Regular.otf}[Path=../../../fonts/]
\setmathfont{PlusJakartaSans-Medium.ttf}[Path=../../../fonts/,range={up/num,up/{Latin,latin},\mathup}]
\usepackage{icomma} % virgule décimale sans espace en mode math
% Caractères mathématiques Unicode absents des polices de texte (Plus Jakarta,
% Archivo Black) : rendus par leur équivalent mathématique, en texte comme en
% formule — sinon ils s'affichent en carré vide (ex. « Arithmétique dans ℤ »).
\usepackage{newunicodechar}
\newunicodechar{ℝ}{\ensuremath{\mathbb{R}}}
\newunicodechar{ℤ}{\ensuremath{\mathbb{Z}}}
\newunicodechar{ℂ}{\ensuremath{\mathbb{C}}}
\newunicodechar{ℕ}{\ensuremath{\mathbb{N}}}
\newunicodechar{ℚ}{\ensuremath{\mathbb{Q}}}
\newunicodechar{𝔻}{\ensuremath{\mathbb{D}}}
\newunicodechar{α}{\ensuremath{\alpha}}
\newunicodechar{β}{\ensuremath{\beta}}
\newunicodechar{γ}{\ensuremath{\gamma}}
\newunicodechar{δ}{\ensuremath{\delta}}
\newunicodechar{ε}{\ensuremath{\varepsilon}}
\newunicodechar{θ}{\ensuremath{\theta}}
\newunicodechar{λ}{\ensuremath{\lambda}}
\newunicodechar{μ}{\ensuremath{\mu}}
\newunicodechar{σ}{\ensuremath{\sigma}}
\newunicodechar{τ}{\ensuremath{\tau}}
\newunicodechar{φ}{\ensuremath{\varphi}}
\newunicodechar{ψ}{\ensuremath{\psi}}
\newunicodechar{ω}{\ensuremath{\omega}}
\newunicodechar{Σ}{\ensuremath{\Sigma}}
% majuscules produites par \MakeUppercase dans les titres (α-aminés -> Α)
\newunicodechar{Α}{\ensuremath{\alpha}}
\newunicodechar{Β}{\ensuremath{\beta}}
\newunicodechar{Γ}{\ensuremath{\Gamma}}
\newunicodechar{Δ}{\ensuremath{\Delta}}
\newunicodechar{Ε}{\ensuremath{\varepsilon}}
\newunicodechar{Θ}{\ensuremath{\Theta}}
\newunicodechar{Λ}{\ensuremath{\Lambda}}
\newunicodechar{Μ}{\ensuremath{\mu}}
\newunicodechar{Τ}{\ensuremath{\tau}}
\newunicodechar{Φ}{\ensuremath{\Phi}}
\newunicodechar{Ψ}{\ensuremath{\Psi}}
\newunicodechar{Ω}{\ensuremath{\Omega}}
\newunicodechar{↦}{\ensuremath{\mapsto}}
\newunicodechar{∈}{\ensuremath{\in}}
\newunicodechar{∉}{\ensuremath{\notin}}
\newunicodechar{∧}{\ensuremath{\wedge}}
\newunicodechar{⇔}{\ensuremath{\Leftrightarrow}}
\newunicodechar{⟺}{\ensuremath{\Longleftrightarrow}}
\newunicodechar{⇒}{\ensuremath{\Rightarrow}}
\newunicodechar{⟹}{\ensuremath{\Longrightarrow}}
\newunicodechar{∘}{\ensuremath{\circ}}
\newunicodechar{∪}{\ensuremath{\cup}}
\newunicodechar{∩}{\ensuremath{\cap}}
\newunicodechar{≡}{\ensuremath{\equiv}}
\newunicodechar{⊂}{\ensuremath{\subset}}
\newunicodechar{⊥}{\ensuremath{\perp}}
\newunicodechar{∃}{\ensuremath{\exists}}
\newunicodechar{∀}{\ensuremath{\forall}}
\newunicodechar{∛}{\ensuremath{\sqrt[3]{\,}}}
\newunicodechar{∜}{\ensuremath{\sqrt[4]{\,}}}
\newunicodechar{⃗}{\ensuremath{{}^{\to}}}
\newunicodechar{ⁿ}{\ifmmode^{n}\else\textsuperscript{\ensuremath{n}}\fi}
\newunicodechar{ˣ}{\ifmmode^{x}\else\textsuperscript{\ensuremath{x}}\fi}
\newunicodechar{ᵇ}{\ifmmode^{b}\else\textsuperscript{\ensuremath{b}}\fi}
\newunicodechar{ʳ}{\ifmmode^{r}\else\textsuperscript{\ensuremath{r}}\fi}
\newunicodechar{ᵘ}{\ifmmode^{u}\else\textsuperscript{\ensuremath{u}}\fi}
\newunicodechar{ʸ}{\ifmmode^{y}\else\textsuperscript{\ensuremath{y}}\fi}
\newunicodechar{ᵃ}{\ifmmode^{a}\else\textsuperscript{\ensuremath{a}}\fi}
\newunicodechar{ᵉ}{\ifmmode^{e}\else\textsuperscript{\ensuremath{e}}\fi}
\newunicodechar{ᵏ}{\ifmmode^{k}\else\textsuperscript{\ensuremath{k}}\fi}
\newunicodechar{ᵖ}{\ifmmode^{p}\else\textsuperscript{\ensuremath{p}}\fi}
\newunicodechar{ᴺ}{\ifmmode^{N}\else\textsuperscript{\ensuremath{N}}\fi}
\newunicodechar{ᶜ}{\ifmmode^{c}\else\textsuperscript{\ensuremath{c}}\fi}
\newunicodechar{ʰ}{\ifmmode^{h}\else\textsuperscript{\ensuremath{h}}\fi}
\newunicodechar{ᵠ}{\ifmmode^{q}\else\textsuperscript{\ensuremath{q}}\fi}
\newunicodechar{ₙ}{\ifmmode_{n}\else\textsubscript{\ensuremath{n}}\fi}
\newunicodechar{ₐ}{\ifmmode_{a}\else\textsubscript{\ensuremath{a}}\fi}
\newunicodechar{ᵢ}{\ifmmode_{i}\else\textsubscript{\ensuremath{i}}\fi}
\newunicodechar{ᵣ}{\ifmmode_{r}\else\textsubscript{\ensuremath{r}}\fi}
\newunicodechar{ₖ}{\ifmmode_{k}\else\textsubscript{\ensuremath{k}}\fi}
\newunicodechar{ᵦ}{\ifmmode_{\beta}\else\textsubscript{\ensuremath{\beta}}\fi}
\newunicodechar{ₓ}{\ifmmode_{x}\else\textsubscript{\ensuremath{x}}\fi}
\newfontfamily\popdisplay{ArchivoBlack-Regular.ttf}[Path=../../../fonts/]
\newfontfamily\poplogo{SpaceGrotesk-Bold.ttf}[Path=../../../fonts/]
\newfontfamily\popsemi{PlusJakartaSans-SemiBold.ttf}[Path=../../../fonts/]
\newfontfamily\popxbold{PlusJakartaSans-ExtraBold.ttf}[Path=../../../fonts/]
\newfontfamily\popxboldls{PlusJakartaSans-ExtraBold.ttf}[Path=../../../fonts/,LetterSpace=3.0]

\setlist[enumerate]{leftmargin=1.7em,itemsep=5pt,topsep=4pt}
\setlist[itemize]{leftmargin=1.7em,itemsep=4pt,topsep=4pt,label={\color{poporange}\textbullet}}
\setlength{\parindent}{0pt}
\setlength{\parskip}{5pt}
\renewcommand{\arraystretch}{1.25}

% pgfplots : style Pop par défaut
\pgfplotsset{
  every axis/.append style={axis line style={popink,line width=1pt},tick style={popink},
    grid style={popline,line width=0.5pt},label style={font=\small},tick label style={font=\footnotesize}},
  popcurve/.style={poporange,line width=1.8pt,no markers,smooth},
}
% Même style disponible dans un \draw TikZ simple (hors pgfplots)
\tikzset{popcurve/.style={poporange,line width=1.8pt}}
\tikzset{popgrid/.style={popline,very thin}}
\colorlet{popcurve}{poporange} % le nom du style est parfois utilisé comme couleur

% ============================================================
% Pill / squiggle
% ============================================================
\newcommand{\poppill}[3]{%
  \tikz[baseline=-0.55ex]\node[draw=popink,line width=1.2pt,rounded corners=2.6pt,%
    fill=#2,inner xsep=6.5pt,inner ysep=3pt]{%
    \popxboldls\fontsize{7}{7}\selectfont\color{#3}\MakeUppercase{#1}};%
}
\newcommand{\popsquiggle}[1]{%
  \tikz[baseline=-0.4ex]{
    \draw[#1,line width=2.6pt,line cap=round]
      plot [smooth] coordinates {(0,0.09) (0.34,0.01) (0.68,0.21) (1.02,0.01) (1.36,0.10)};
  }%
}
% Mot-clé surligné (marqueur orange sous le texte)
\newcommand{\cle}[1]{{\popxbold\color{popdark}#1}}

% ============================================================
% En-tête / pied
% ============================================================
\pagestyle{fancy}
\fancyhf{}
\renewcommand{\headrulewidth}{0pt}
\renewcommand{\footrulewidth}{0pt}
\lhead{\raisebox{-0.15ex}{\poplogo\fontsize{13}{13}\selectfont\color{popink}OnBuch\color{poporange}+}}
\rhead{\poppill{\DOCMATIERE\ \textbullet\ \DOCNIVEAU\ \textbullet\ Leçon \DOCLECON}{white}{popink}}
\lfoot{\popxbold\fontsize{7.6}{7.6}\selectfont\color{popmuted}\MakeUppercase{Cours signé OnBuch+}\ \textbullet\ {\popsemi\DOCTITRE}}
\rfoot{\tikz[baseline=-0.6ex]\node[rounded corners=8.5pt,fill=popink,minimum height=17pt,minimum width=17pt,inner xsep=5pt,inner sep=0]{\popxbold\fontsize{8}{8}\selectfont\color{popcream}\thepage\,/\,\pageref*{LastPage}};}

% ============================================================
% Titres
% ============================================================
\newcommand{\chaptitre}[2]{%
  \begin{tcolorbox}[enhanced,colback=poporange,colframe=popink,boxrule=2.6pt,arc=11pt,
      drop shadow=popink,left=15pt,right=15pt,top=12pt,bottom=11pt,before skip=3pt,after skip=18pt]
    \poppill{#2}{popink}{popcream}\par\vspace{9pt}
    {\raggedright\hyphenpenalty=10000\exhyphenpenalty=10000\popdisplay\fontsize{22}{24}\selectfont\color{popcream}\MakeUppercase{#1}\par}
  \end{tcolorbox}
}
% Section numérotée : pastille orange avec le numéro + titre Archivo
\newcounter{popsec}
\newcommand{\coursec}[1]{%
  \stepcounter{popsec}\setcounter{popsub}{0}%
  \par\vspace{16pt}\noindent
  \begin{minipage}{\linewidth}
  \tikz[baseline=-0.7ex]\node[draw=popink,line width=1.6pt,fill=poporange,rounded corners=4pt,minimum size=22pt,inner sep=0]{\popdisplay\fontsize{12}{12}\selectfont\color{popcream}\thepopsec};%
  \hspace{8pt}{\popdisplay\fontsize{15}{17}\selectfont\color{popink}\MakeUppercase{#1}}\par
  \vspace{4pt}\noindent{\color{poporange}\rule{36pt}{3.6pt}}
  \end{minipage}\par\nopagebreak\vspace{8pt}
}
\newcounter{popsub}
\newcommand{\courssub}[1]{%
  \stepcounter{popsub}%
  \par\vspace{9pt}\noindent{\popxbold\fontsize{11.5}{13}\selectfont\color{popink}{\color{poporange}\thepopsec.\thepopsub}\hspace{6pt}#1}\par\nopagebreak\vspace{4pt}
}

% ============================================================
% Boîtes pédagogiques — même recette : fond blanc, bordure épaisse colorée,
% ombre dure encre, badge plein. Titre optionnel [#1].
% ============================================================
\tcbset{popbase/.style={breakable,enhanced,colback=white,boxrule=2.3pt,arc=9pt,
  shadow={3pt}{-3pt}{0pt}{popink},left=14pt,right=14pt,top=11pt,bottom=11pt,before skip=11pt,after skip=15pt,
  fontupper=\popsemi\fontsize{10.6}{14.5}\selectfont\color{popink},
  fontlower=\popsemi\fontsize{10.6}{14.5}\selectfont\color{popink}}}
% Coche dessinée (le glyphe ✓ n'existe pas dans les polices du document)
\newcommand{\popcheck}[1]{\tikz[baseline=-0.5ex]\draw[#1,line width=1.6pt,line cap=round,line join=round](-0.13,0.0)--(-0.04,-0.09)--(0.13,0.1);}
\providecommand{\coche}{\popcheck{popgreen}}
\providecommand{\resultat}[1]{\par\smallskip\noindent\fcolorbox{popgreen}{popgreenL}{\parbox{\dimexpr\linewidth-2\fboxsep-2\fboxrule\relax}{\textbf{Résultat.} #1}}\par\smallskip}
\newcommand{\popboxhead}[3]{\poppill{#1}{#2}{white}\ifx\relax#3\relax\else\hspace{6pt}{\popxbold\fontsize{10.6}{12}\selectfont\color{popink}#3}\fi\par\vspace{7pt}}

\newtcolorbox{aretenir}[1][]{popbase,colframe=popgreen,before upper={\popboxhead{À retenir}{popgreen}{#1}}}
\newtcolorbox{exemplebox}[1][]{popbase,colframe=poporange,before upper={\popboxhead{Exemple}{poporange}{#1}}}
\newtcolorbox{definition}[1][]{popbase,colframe=popblue,before upper={\popboxhead{Définition}{popblue}{#1}}}
\newtcolorbox{propriete}[1][]{popbase,colframe=poppurple,before upper={\popboxhead{Propriété}{poppurple}{#1}}}
\newtcolorbox{methode}[1][]{popbase,colframe=popgold,before upper={\popboxhead{Méthode}{popgold}{#1}}}
\newtcolorbox{attention}[1][]{popbase,colframe=poppink,before upper={\popboxhead{Attention}{poppink}{#1}}}
\newtcolorbox{experience}[1][]{popbase,colframe=popblue,colback=popblueL,before upper={\popboxhead{Expérience}{popblue}{#1}}}
% « Le savais-tu ? » : carte violette pleine (contexte, culture, Cameroun)
\newtcolorbox{savaistu}[1][]{popbase,colback=poppurple,colframe=popink,
  fontupper=\popsemi\fontsize{10.4}{14.2}\selectfont\color{white},
  before upper={\poppill{Le savais-tu\,?}{white}{poppurple}\ifx\relax#1\relax\else\hspace{6pt}{\popxbold\color{white}#1}\fi\par\vspace{7pt}}}
% Exercice résolu : énoncé en haut, \tcblower puis la solution sur fond crème
\newcounter{popexo}\newcounter{popexr}
\newtcolorbox{exoresolu}[1][]{popbase,colframe=poporange,colbacklower=popcream,
  segmentation style={popink,line width=1.2pt,dashed},
  before upper={\stepcounter{popexr}\popboxhead{Exercice résolu \thepopexr}{poporange}{#1}},
  before lower={\poppill{Solution}{popink}{popcream}\par\vspace{6pt}}}
% Exercice (non résolu en place) — la correction est dans la partie Corrigés
\newtcolorbox{exercice}[1][]{popbase,colframe=popink,
  before upper={\stepcounter{popexo}\popboxhead{Exercice \thepopexo}{popink}{#1}}}
% Corrigé
\newtcolorbox{corrige}[1][]{popbase,colframe=popgreen,colback=popgreenL,
  before upper={\popboxhead{Corrigé}{popgreen}{#1}}}
% Objectifs / prérequis / bilan
\newtcolorbox{objectifs}{popbase,colframe=popink,colback=poporangeL,
  before upper={\popboxhead{Objectifs de la leçon}{popink}{}}}
\newtcolorbox{prerequis}{popbase,colframe=popink,colback=popblueL,
  before upper={\popboxhead{Prérequis}{popblue}{}}}
\newtcolorbox{bilan}{popbase,colframe=popink,colback=white,boxrule=2.8pt,
  before upper={\popboxhead{Fiche bilan}{poporange}{L'essentiel à connaître pour l'examen}}}
% Figure encadrée avec légende
\newtcolorbox{popfigure}[1][]{enhanced,breakable=false,colback=white,colframe=popline,boxrule=1.4pt,arc=8pt,
  left=10pt,right=10pt,top=10pt,bottom=8pt,before skip=10pt,after skip=12pt,halign=center,
  lower separated=false,halign lower=center,
  fontlower=\popsemi\fontsize{9}{11.5}\selectfont\color{popmuted},
  #1}
\newcommand{\legende}[1]{\par\vspace{6pt}{\popsemi\fontsize{9}{11.5}\selectfont\color{popmuted}#1}}

% ============================================================
% Signature OnBuch+
% ============================================================
\newcommand{\poplogoinline}[1]{{\poplogo\fontsize{#1}{#1}\selectfont\color{popink}OnBuch\color{poporange}+}}
\newcommand{\signatureonbuch}{%
  \begin{tcolorbox}[enhanced,colback=popink,colframe=popink,boxrule=2.2pt,arc=10pt,
      drop shadow=poporange,left=14pt,right=14pt,top=10pt,bottom=10pt,before skip=0pt,after skip=0pt]
    \tikz[baseline=-0.7ex]\node[circle,fill=popgreen,draw=popcream,line width=1.3pt,minimum size=22pt,inner sep=0]{\popcheck{white}};%
    \hspace{9pt}\begin{minipage}[c]{0.62\linewidth}
      {\popxbold\fontsize{12}{13}\selectfont\color{popcream}Signé\ \textbullet\ {\poplogo OnBuch\color{poporange}+}}\par\vspace{2pt}
      {\raggedright\popsemi\fontsize{8}{10.5}\selectfont\color{popline}\MakeUppercase{Équipe pédagogique OnBuch+}\\\MakeUppercase{Conforme au programme officiel MINESEC}\par}
    \end{minipage}\hfill
    {\poplogo\fontsize{22}{22}\selectfont\color{popcream}OnBuch\color{poporange}+}
  \end{tcolorbox}%
}

% ============================================================
% Page de garde
% #1 titre · #2 matière · #3 niveau · #4 module · #5 n° leçon · #6 durée ·
% #7 description / objectifs (texte ou itemize)
% ============================================================
\newcommand{\pagegarde}[7]{%
  \thispagestyle{empty}
  \newgeometry{a4paper,margin=1.7cm,top=1.6cm,bottom=1.4cm}%
  \begin{tikzpicture}[remember picture,overlay]
    \fill[poporange!18] ([xshift=-3.2cm,yshift=-3.6cm]current page.north east) circle (4.6cm);
    \fill[poppurple!14] ([xshift=2.4cm,yshift=7.6cm]current page.south west) circle (3.8cm);
  \end{tikzpicture}%
  {\poplogo\fontsize{30}{30}\selectfont\color{popink}OnBuch\color{poporange}+}\hfill
  \raisebox{0.6ex}{\poppill{Cours signé \textbullet\ Premium}{popink}{popcream}}\par
  \vspace{1pt}\popsquiggle{poppurple}\par
  \vspace{8pt}
  \begin{tcolorbox}[enhanced,colback=poporange,colframe=popink,boxrule=2.8pt,arc=13pt,
      drop shadow=popink,left=18pt,right=18pt,top=12pt,bottom=13pt,after skip=10pt]
    \poppill{#2 \textbullet\ #3}{popink}{popcream}\hspace{5pt}\poppill{#4}{white}{popink}\par\vspace{8pt}
    {\popdisplay\fontsize{14}{14}\selectfont\color{popink}LEÇON #5}\par\vspace{4pt}
    {\raggedright\hyphenpenalty=10000\exhyphenpenalty=10000\popdisplay\fontsize{25}{27}\selectfont\color{popcream}\MakeUppercase{#1}\par}
    \vspace{8pt}
    {\popxbold\fontsize{9.5}{9.5}\selectfont\color{popink}\MakeUppercase{Durée conseillée : #6}}
  \end{tcolorbox}
  \begin{tcolorbox}[enhanced,colback=white,colframe=popink,boxrule=2.2pt,arc=10pt,
      drop shadow=popink,left=16pt,right=16pt,top=10pt,bottom=10pt,after skip=8pt]
    \poppill{Dans cette leçon}{poppurple}{white}\par\vspace{5pt}
    \popsemi\fontsize{9.3}{12.6}\selectfont\color{popink}#7
  \end{tcolorbox}
  \noindent\begin{minipage}[t]{0.487\linewidth}
    \begin{tcolorbox}[enhanced,colback=popgreen,colframe=popink,boxrule=2.2pt,arc=10pt,
        drop shadow=popink,left=11pt,right=11pt,top=10pt,bottom=10pt,height=3.1cm,valign=center]
      \tcbox[colback=white,colframe=popink,boxrule=1.3pt,arc=5pt,boxsep=0pt,nobeforeafter,
        left=3pt,right=3pt,top=3pt,bottom=3pt,valign=center]{\qrcode[height=1.9cm]{https://play.google.com/store/apps/details?id=com.luvvix.onbuch&hl=fr}}%
      \hspace{8pt}\begin{minipage}[c]{0.5\linewidth}
        {\popdisplay\fontsize{11}{12.5}\selectfont\color{white}\MakeUppercase{Télécharge OnBuch}}\par\vspace{4pt}
        {\raggedright\popsemi\fontsize{8.4}{11}\selectfont\color{white}Gratuit \textbullet\ Google Play\\Tuteur IA, annales, concours, résultats\par}
      \end{minipage}
    \end{tcolorbox}
  \end{minipage}\hfill
  \begin{minipage}[t]{0.487\linewidth}
    \begin{tcolorbox}[enhanced,colback=poppurple,colframe=popink,boxrule=2.2pt,arc=10pt,
        drop shadow=popink,left=11pt,right=11pt,top=10pt,bottom=10pt,height=3.1cm,valign=center]
      \tikz[baseline=-0.7ex]\node[circle,fill=white,draw=popink,line width=1.3pt,minimum size=1.6cm,inner sep=0]{\popdisplay\fontsize{24}{24}\selectfont\color{poppurple}+};%
      \hspace{8pt}\begin{minipage}[c]{0.6\linewidth}
        {\popdisplay\fontsize{11}{12.5}\selectfont\color{white}\MakeUppercase{Passe à OnBuch+}}\par\vspace{4pt}
        {\raggedright\popsemi\fontsize{8.4}{11}\selectfont\color{white}Tous les cours signés, Tuteur IA illimité, fiches PDF — onglet OnBuch+ de l'app\par}
      \end{minipage}
    \end{tcolorbox}
  \end{minipage}
  \vspace{8pt}
  \popcontenu
  \vfill
  \signatureonbuch
  \restoregeometry
  \newpage
}

% Rangée « Ce cours contient » (page de garde)
\newcommand{\poptile}[3]{% #1 couleur #2 gros texte #3 légende
  \begin{tcolorbox}[enhanced,colback=#1,colframe=popink,boxrule=2pt,arc=9pt,shadow={2.5pt}{-2.5pt}{0pt}{popink},
      left=6pt,right=6pt,top=9pt,bottom=9pt,halign=center,valign=center,height=1.75cm,width=\linewidth,nobeforeafter]
    {\popdisplay\fontsize{12.5}{13}\selectfont\color{white}\MakeUppercase{#2}}\par\vspace{3pt}
    {\centering\popsemi\fontsize{7.8}{9.5}\selectfont\color{white}#3\par}
  \end{tcolorbox}}
\newcommand{\popcontenu}{%
  \par\noindent{\popxboldls\fontsize{7.5}{7.5}\selectfont\color{popmuted}CE COURS SIGNÉ CONTIENT}\par\vspace{7pt}
  \noindent\begin{minipage}[t]{0.235\linewidth}\poptile{popblue}{Cours}{complet, illustré, pas à pas}\end{minipage}\hfill
  \begin{minipage}[t]{0.235\linewidth}\poptile{poporange}{Méthodes}{et exercices résolus}\end{minipage}\hfill
  \begin{minipage}[t]{0.235\linewidth}\poptile{poppink}{Exercices}{type examen + corrigés}\end{minipage}\hfill
  \begin{minipage}[t]{0.235\linewidth}\poptile{popgreen}{Fiche bilan}{l'essentiel à retenir}\end{minipage}\par}

% Sommaire
\newtcolorbox{sommaire}{enhanced,colback=white,colframe=popink,boxrule=2.2pt,arc=10pt,
  drop shadow=popink,left=18pt,right=18pt,top=13pt,bottom=13pt,before skip=6pt,after skip=20pt,
  before upper={\poppill{Sommaire}{poppurple}{white}\par\vspace{9pt}\popsemi\fontsize{10.6}{14.5}\selectfont\color{popink}}}

% Signature de fin de cours — poussée tout en bas de la dernière page
\newcommand{\findecours}{%
  \par\vfill\nopagebreak
  \begin{center}\popsquiggle{poporange}\end{center}\vspace{4pt}
  \signatureonbuch
}
\AtBeginDocument{\def\llbracket{\mathopen{[\mkern-3.5mu[}}\def\rrbracket{\mathclose{]\mkern-3.5mu]}}} % intervalles d'entiers (glyphe ⟦ absent de la police)
% Glyphes absents de Fira Math : redéfinis à partir de glyphes disponibles
\AtBeginDocument{%
  \def\setminus{\mathbin{\backslash}}\let\smallsetminus\setminus
  \def\vdots{\mathord{\vbox{\baselineskip3.2pt\lineskiplimit0pt\kern4pt\hbox{.}\hbox{.}\hbox{.}}}}%
  \def\ddots{\mathinner{\mkern1mu\raise6.5pt\hbox{.}\mkern2mu\raise3.6pt\hbox{.}\mkern2mu\raise0.7pt\hbox{.}\mkern1mu}}%
  \def\triangle{\mathord{\scalebox{0.9}{$\symup{\Delta}$}}}%
  \def\nabla{\mathord{\raisebox{\depth}{\scalebox{1}[-1]{$\symup{\Delta}$}}}}%
  \def\checkmark{\popcheck{popdark}}%
  \def\star{\mathbin{\ast}}\def\bigstar{\mathord{\ast}}%
  \def\diamond{\mathbin{\scalebox{0.6}{$\Diamond$}}}%
  \def\bigcup{\mathop{\mathchoice{\raisebox{-0.3em}{\scalebox{1.7}{$\cup$}}}{\scalebox{1.25}{$\cup$}}{\cup}{\cup}}\displaylimits}%
  \def\bigcap{\mathop{\mathchoice{\raisebox{-0.3em}{\scalebox{1.7}{$\cap$}}}{\scalebox{1.25}{$\cap$}}{\cap}{\cap}}\displaylimits}%
  \def\lesseqqgtr{\mathrel{\lessgtr}}\def\bigoplus{\mathop{\oplus}\displaylimits}%
}
\providecommand{\ssi}{si et seulement si} % abréviation fréquente
\AtBeginDocument{\providecommand{\wideparen}{\overparen}} % arc de cercle

%% FIGFIX-BEGIN v4 — correctifs globaux des figures (docs/onbuch-generation/tools/figfix)
\usepackage{adjustbox}\usepackage{etoolbox}
% toute figure TikZ/pgfplots plus large que la ligne est réduite à la largeur disponible
\BeforeBeginEnvironment{tikzpicture}{\begin{adjustbox}{max width=\linewidth}}
\AfterEndEnvironment{tikzpicture}{\end{adjustbox}}
% légendes pgfplots lisibles par-dessus les courbes
\pgfplotsset{every axis/.append style={legend style={fill=white,fill opacity=0.92,text opacity=1,draw=black!15,inner sep=3pt}}}
% étiquettes d'arêtes (syntaxe edge["..."]) avec fond blanc
\tikzset{every edge quotes/.append style={fill=white,inner sep=1.2pt}}
%% FIGFIX-END
\begin{document}

\pagegarde{Le Cameroun indépendant}{Histoire}{Tle Tech}{Histoire – classe de Terminale technique}{N07}{2 séances (fiche de progression harmonisée)}{Cette leçon retrace l'histoire politique, économique et sociale du Cameroun depuis les indépendances de 1960-1961 jusqu'à nos jours. Elle analyse les deux grandes périodes que sont la présidence Ahidjo (1960-1982) et la présidence Biya (1982 à nos jours), à travers documents, cartes et données chiffrées.
\vspace{4pt}
\begin{multicols}{2}\small
\begin{itemize}
\item À la fin de cette leçon, je sais situer chronologiquement les grandes étapes de l'indépendance et de la réunification du Cameroun (1960-1972)
\item Je sais décrire l'organisation politique du Cameroun sous Ahidjo (parti unique, État unitaire)
\item Je sais analyser les réalisations économiques et sociales du Cameroun de 1960 à 1982 à partir de données chiffrées
\item Je sais expliquer les causes et les circonstances de la succession de 1982
\item Je sais caractériser la gouvernance du Cameroun sous Paul Biya (libéralisation, multipartisme, décentralisation)
\item Je sais identifier les défis contemporains du Cameroun (crise économique, sécurité, cohésion nationale)
\end{itemize}\end{multicols}}

\begin{sommaire}
\begin{multicols}{2}
\begin{enumerate}
\item Les indépendances et la réunification (1960-1961)
\item Le Cameroun sous Ahidjo : la construction de l'État unitaire (1961-1972)
\item Le bilan économique et social de l'ère Ahidjo (1960-1982)
\item La succession de 1982 et l'avènement de Paul Biya
\item Le Cameroun de 1982 à nos jours : mutations politiques
\item Les défis du Cameroun contemporain
\item Activité d'intégration
\item Méthodes et exercices résolus
\item Exercices
\item Corrigés des exercices
\item Fiche bilan
\end{enumerate}
\end{multicols}
\end{sommaire}

\begin{objectifs}
\begin{itemize}
\item À la fin de cette leçon, je sais situer chronologiquement les grandes étapes de l'indépendance et de la réunification du Cameroun (1960-1972)
\item Je sais décrire l'organisation politique du Cameroun sous Ahidjo (parti unique, État unitaire)
\item Je sais analyser les réalisations économiques et sociales du Cameroun de 1960 à 1982 à partir de données chiffrées
\item Je sais expliquer les causes et les circonstances de la succession de 1982
\item Je sais caractériser la gouvernance du Cameroun sous Paul Biya (libéralisation, multipartisme, décentralisation)
\item Je sais identifier les défis contemporains du Cameroun (crise économique, sécurité, cohésion nationale)
\item Je sais construire une frise chronologique et un croquis cartographique de l'évolution territoriale du Cameroun
\item Je sais rédiger et justifier une conclusion argumentée à partir de documents historiques
\item Je sais appliquer les notions de la leçon à des situations concrètes de la vie citoyenne camerounaise
\end{itemize}
\end{objectifs}

\begin{prerequis}
\begin{itemize}
\item La colonisation et la mise sous tutelle du Cameroun (classe de Première)
\item Les indépendances africaines et le contexte de la Guerre froide
\item La notion d'État, de régime politique et de parti unique
\item Les grandes lignes de la géographie administrative du Cameroun (régions, capitale)
\item La lecture de documents historiques (textes, cartes, statistiques)
\end{itemize}
\end{prerequis}

\newpage

\coursec{Les indépendances et la réunification (1960-1961)}

\noindent En janvier 1960, un jeune Camerounais de Douala voit hisser pour la première fois le drapeau vert-rouge-jaune sur le palais des gouverneurs : son pays devient indépendant. Vingt-deux ans plus tard, le 4 novembre 1982, son fils assiste, médusé, à la démission télévisée du président Ahmadou Ahidjo et à l'arrivée au pouvoir de Paul Biya. Aujourd'hui, son petit-fils, élève en Terminale, s'interroge : comment le Cameroun est-il passé de la tutelle franco-britannique à un État unifié et stable ? Quels défis politiques, économiques et sociaux a-t-il traversés depuis 1960 ? C'est à ces questions que cette leçon répond.

\courssub{L'indépendance du Cameroun oriental (1er janvier 1960)}

\noindent Après la Seconde Guerre mondiale, le Cameroun, ancienne colonie allemande partagée entre la France et le Royaume-Uni sous mandat de la Société des Nations (SDN) puis tutelle de l'ONU, connaît une montée des nationalismes. L'\cle{Union des Populations du Cameroun (UPC)}, fondée en 1948, réclame l'indépendance immédiate et la réunification. Face à la répression coloniale, l'UPC bascule dans la clandestinité et la lutte armée dès 1955 (insurrection du 25 mai 1955 à Douala, puis maquis dans le Sanaga-Maritime et le Bamiléké).

\begin{savaistu}[La tutelle, un cadre international contraignant]
Contrairement aux colonies de peuplement, les territoires sous tutelle (ex-mandats SDN) font l'objet de rapports annuels à l'ONU. La France doit rendre des comptes sur l'évolution politique, économique et sociale du Cameroun. Cette pression internationale accélère le processus d'autonomie interne (loi-cadre Defferre, 1956) puis d'indépendance.
\end{savaistu}

\noindent La France, confrontée à la guerre d'Algérie et à la pression de l'ONU, organise un référendum d'autodétermination le 1er janvier 1960. Le « Oui » à l'indépendance l'emporte massivement. Le \textbf{1er janvier 1960}, le Cameroun oriental (ex-Cameroun français) proclame son indépendance. \textbf{Ahmadou Ahidjo}, alors Premier ministre, devient le premier Président de la République le 5 mai 1960, après l'adoption d'une Constitution le 21 février 1960 instaurant un régime présidentiel fort.

\begin{definition}[Indépendance]
Acte juridique et politique par lequel un territoire accède à la souveraineté internationale, se dotant d'institutions propres, d'une armée, d'une diplomatie et d'une monnaie (ou d'une union monétaire), sans tutelle extérieure.
\end{definition}

\courssub{Le cas particulier du Cameroun britannique : les plébiscites de février 1961}

\noindent Le Cameroun britannique, administré depuis le Nigeria, est divisé en deux zones distinctes : le \textbf{Cameroun septentrional} (au nord, peuplé majoritairement de Foulbés et de groupes soudanais) et le \textbf{Cameroun méridional} (au sud, peuplé de Bantoïs, culturellement proches des populations du Cameroun oriental). L'ONU impose une consultation populaire pour déterminer leur avenir.

\begin{definition}[Plébiscite]
Consultation populaire par laquelle un peuple se prononce directement sur une question précise de rattachement territorial ou de souveraineté (ici, le choix entre l'intégration au Nigeria indépendant ou la réunification avec la République du Cameroun).
\end{definition}

\begin{methode}[Analyser une carte historique : les plébiscites de 1961]
\begin{enumerate}
    \item \textbf{Identifier le titre, la légende et l'échelle :} « Cameroun britannique : résultats des plébiscites du 11 février 1961 ». La légende distingue deux couleurs pour les deux options (Nigeria / Cameroun).
    \item \textbf{Repérer les entités géographiques :} Localiser la ligne de partage (parallèle environ 7° Nord), les capitales régionales (Buea au Sud, Garoua au Nord) et la frontière avec le Nigeria.
    \item \textbf{Lire les résultats chiffrés :} Noter les pourcentages pour chaque zone.
    \item \textbf{Interpréter les facteurs :} Expliquer le vote « Nigeria » au Nord (liens économiques avec le Nord-Nigeria, islam, peur de la domination bamiléké/bassa) et le vote « Cameroun » au Sud (liens culturels, langue anglaise commune avec l'administration, rejet de l'absorption par le Nigeria).
    \item \textbf{Conclure sur les conséquences territoriales :} Naissance du Cameroun occidental (ex-Cameroun méridional britannique) et perte du Septentrional pour la réunification.
\end{enumerate}
\end{methode}

\begin{popfigure}
\begin{tikzpicture}[scale=1.1]
% Contour très simplifié du Cameroun (polygone approximatif)
\draw[popdark, thick, fill=popcream] 
(0,0) -- (8,0) -- (9,2) -- (8.5,5) -- (7,6.5) -- (5,7) -- (3,6.5) -- (1.5,5) -- (1,2) -- (2,0) -- cycle;
% Frontière Nord/Sud (ligne de partage des plébiscites)
\draw[popblue, dashed, thick] (1.5,3.5) -- (8,3.5) node[midway, above, sloped, font=\small, popblue] {Ligne de partage (7° N)};
% Zone Nord (Septentrional) - Vote Nigeria
\fill[poporange!40] (1.5,3.5) -- (8,3.5) -- (8.5,5) -- (7,6.5) -- (5,7) -- (3,6.5) -- (1.5,5) -- cycle;
\node[font=\small\bfseries, popdark] at (5.5, 5.5) {CAMEROUN SEPTENTRIONAL};
\node[font=\small, popdark, align=center] at (5.5, 5.0) {Plébiscite 11 fév. 1961\\ \textbf{60\% Nigeria / 40\% Cameroun}\\ \rightarrow Intégration au Nigeria (juin 1961)};
% Zone Sud (Méridional) - Vote Cameroun
\fill[popgreen!30] (0,0) -- (8,0) -- (9,2) -- (8.5,5) -- (8,3.5) -- (1.5,3.5) -- (1,2) -- (2,0) -- cycle;
\node[font=\small\bfseries, popdark] at (4.5, 1.8) {CAMEROUN MÉRIDIONAL};
\node[font=\small, popdark, align=center] at (4.5, 1.2) {Plébiscite 11 fév. 1961\\ \textbf{70,5\% Cameroun / 29,5\% Nigeria}\\ \rightarrow Réunification (1er oct. 1961)};
% Villes repères
\foreach \x/\y/\name/\col in {2.5/5.5/Garoua/poporange, 6.5/2.5/Buea/popgreen, 4/0.5/Douala/popdark, 5.5/4.5/Yaoundé/popdark}
    \draw[fill=\col] (\x,\y) circle (2.5pt) node[right, font=\tiny, \col] {\name};
% Flèche Nord
\draw[-{Stealth[length=3mm]}, poporange, thick] (5.5, 6.2) -- (7.5, 7) node[right, font=\small, poporange] {Vers le Nigeria};
% Légende
\begin{scope}[shift={(9.5, 6)}]
\draw[fill=poporange!40] (0,0) rectangle (0.5,0.3) node[right, font=\small] {Vote Nigeria (Nord)};
\draw[fill=popgreen!30] (0,-0.5) rectangle (0.5,-0.2) node[right, font=\small] {Vote Cameroun (Sud)};
\draw[popblue, dashed, thick] (0,-1) -- (0.5,-1) node[right, font=\small] {Frontière plébiscite};
\draw[popdark, thick] (0,-1.5) -- (0.5,-1.5) node[right, font=\small] {Frontière internationale 1960};
\end{scope}
% Nord
\node[font=\small, rotate=90, anchor=south] at (-0.5, 3.5) {NORD};
% Echelle indicative
\draw[popdark, |-|] (0.5, -0.5) -- (2.5, -0.5) node[midway, below, font=\tiny] {~200 km};
\end{tikzpicture}
\legende{Carte schématique du Cameroun britannique montrant la partition issue des plébiscites du 11 février 1961. Au nord, le rattachement au Nigeria ; au sud, le choix de la réunification avec la République du Cameroun.}
\end{popfigure}

\noindent Le \textbf{11 février 1961}, sous l'égide de l'ONU, les populations votent :
\begin{itemize}
    \item Au \textbf{Cameroun septentrional} : 60 \% choisissent l'intégration à la Fédération du Nigeria (indépendant depuis octobre 1960). L'intégration est effective le 1er juin 1961.
    \item Au \textbf{Cameroun méridional} : 70,5 \% choisissent la réunification avec la République du Cameroun.
\end{itemize}

\begin{attention}[Piège chronologique : Indépendance vs Réunification]
Ne confondez jamais les deux dates fondatrices :
\begin{itemize}
    \item \textbf{1er janvier 1960} : Indépendance du Cameroun \emph{oriental} (ex-français). Fin de la tutelle française.
    \item \textbf{1er octobre 1961} : Réunification du Cameroun \emph{occidental} (ex-britannique méridional) avec le Cameroun oriental. Naissance de la Fédération.
\end{itemize}
L'examen (Probatoire/Bac) sanctionne systématiquement cette confusion dans les dissertations ou les commentaires de documents.
\end{attention}

\courssub{La naissance de la République fédérale du Cameroun (1er octobre 1961)}

\noindent La réunification ne se fait pas par absorption, mais par une \textbf{fédération} négociée lors de la Conférence de Foumban (juillet 1961). Les délégués du Cameroun oriental (Ahidjo) et du Cameroun méridional (John Ngu Foncha) définissent les institutions de la \textbf{République fédérale du Cameroun}.

\begin{definition}[Réunification]
Réunion en un seul État souverain des deux parties du Cameroun séparées depuis la défaite allemande de 1916 et le partage franco-britannique de 1919 (confirmé par la SDN en 1922). Elle s'opère le 1er octobre 1961 par l'adhésion de l'État fédéré du Cameroun occidental à la République du Cameroun.
\end{definition}

\noindent La structure fédérale (Constitution du 1er septembre 1961) repose sur :
\begin{itemize}
    \item \textbf{Deux États fédérés :} l'\textbf{État du Cameroun oriental} (capital Yaoundé, langue administrative française) et l'\textbf{État du Cameroun occidental} (capital Buea, langue administrative anglaise).
    \item \textbf{Un pouvoir fédéral fort :} Président de la Fédération (Ahidjo), Vice-président (Foncha), Assemblée fédérale, gouvernement fédéral (défense, monnaie, diplomatie, finances).
    \item \textbf{Bilinguisme officiel :} Français et anglais langues de travail et d'enseignement.
    \item \textbf{John Ngu Foncha} devient Premier ministre du Cameroun occidental et Vice-président de la République fédérale.
\end{itemize}

\begin{aretenir}[Les piliers de la Fédération (1961-1972)]
\begin{itemize}
    \item \textbf{État fédéral :} Souveraineté extérieure, armée, monnaie, douane, justice suprême.
    \item \textbf{États fédérés :} Assemblées législatives propres, gouvernements propres (Premiers ministres), administration locale, éducation primaire/secondaire, justice coutumière.
    \item \textbf{Symboles :} Drapeau vert-rouge-jaune \emph{à deux étoiles} (une par État fédéré), hymne « Ô Cameroun, berceau de nos ancêtres ».
\end{itemize}
\end{aretenir}

\courssub{La répression de la rébellion UPC et la consolidation du pouvoir d'Ahidjo}

\noindent Paradoxalement, l'indépendance et la réunification s'accompagnent d'une guerre civile larvée. L'UPC, interdite en 1955 (France) puis en 1957 (Royaume-Uni), refuse la « néocolonisation » et la fédération jugée « Ahidjo-Foncha ». Ses maquis (Bassa : Ernest Ouandié, Bamiléké : Wambo le Roi, Ndeh Ntumazah) contrôlent des zones rurales.

\begin{experience}[Analyse de document : Le discours d'Ahidjo du 1er janvier 1960 (extrait)]
\textbf{Document :} « Camerounais, l'indépendance n'est pas un but, c'est un point de départ. Elle nous impose le devoir de construire, dans l'ordre et la discipline, une nation prospère... Tout acte de sabotage, toute atteinte à l'unité nationale sera réprimé sans faiblesse. »
\begin{enumerate}
    \item \textbf{Nature :} Discours officiel d'investiture (source primaire, discours politique).
    \item \textbf{Contexte :} Jour de l'indépendance, insurrection UPC active, guerre froide.
    \item \textbf{Analyse :} Le mot « ordre » et « discipline » annonce le parti unique (1966) et la répression. « Unité nationale » justifie la lutte contre l'UPC (qualifiée de tribale/séparatiste).
    \item \textbf{Conclusion :} L'indépendance est confisquée par un pouvoir autoritaire au nom de la stabilité.
\end{enumerate}
\end{experience}

\noindent La stratégie d'Ahidjo combine action militaire (aide française, armée camerounaise) et isolation politique :
\begin{itemize}
    \item Assassinat de \textbf{Félix-Roland Moumié} (Secrétaire général UPC) à Genève en 1960 (empoisonnement par agent SDECE).
    \item Exécution des derniers chefs historiques : \textbf{Ernest Ouandié} (janvier 1971) et \textbf{Wambo le Roi} (1971) après procès publics.
    \item Création de l'\textbf{Union Camerounaise (UC)} (1958), fusion avec le KNDP de Foncha (1966) pour former l'\textbf{Union Nationale Camerounaise (UNC)}, parti unique \emph{de facto} (1966) puis \emph{de jure} (1972).
\end{itemize}

\begin{exoresolu}[Sujet type Bac : « Expliquez comment Ahmadou Ahidjo a consolidé son pouvoir entre 1960 et 1966. »]
\textbf{Éléments de réponse structurée :}
\begin{enumerate}
    \item \textbf{Légitimité électorale :} Élection présidentielle de 1960 (candidat unique), référendum constitutionnel 1960.
    \item \textbf{Gestion de la réunification :} Négociation de Foumban (1961), intégration de Foncha (Vice-présidence), marginalisation des leaders anglophones hostiles (Endeley).
    \item \textbf{Répression de l'opposition armée :} Liquidation physique des chefs UPC (Moumié 1960, Ouandié 1971), quadrillage militaire des zones rebelles (Bassa, Bamiléké) avec appui français.
    \item \textbf{Construction du parti unique :} Fusion UC-KNDP (1966) $\rightarrow$ UNC. Interdiction des partis d'opposition (1966). Contrôle de la société civile (syndicats, jeunesse, femmes) via les organisations de masse de l'UNC.
    \item \textbf{Concentration institutionnelle :} Constitution de 1961 (fédérale mais présidentialiste), suppression du poste de Vice-président (1966, remplacé par un Premier ministre fédéral), réforme de 1969 renforçant l'exécutif.
\end{enumerate}
\textbf{Conclusion :} Ahidjo transforme l'indépendance formelle en pouvoir personnel autoritaire, utilisant l'argument de l'unité nationale pour éliminer toute alternative politique.
\tcblower
\textbf{Corrigé détaillé :} La réponse doit articuler \textbf{violence légitime (monopole de la force)} et \textbf{ingénierie institutionnelle}. Noter le rôle clé de l'armée (créée 1960, encadrée par officiers français) et de la diplomatie (reconnaissance internationale rapide). L'erreur fréquente est d'oublier l'apport de la réunification (apport de Foncha et de l'appareil administratif britannique) dans la consolidation.
\end{exoresolu}

\courssub{Frise chronologique de synthèse (1955-1972)}

\begin{popfigure}
\begin{tikzpicture}[x=0.45cm, y=1.2cm, font=\small]
% Axe principal
\draw[popdark, thick, ->] (0,0) -- (42,0) node[right] {Années};
% Graduations principales
\foreach \x/\year in {0/1955, 5/1960, 10/1965, 15/1970, 20/1972, 25/1975, 30/1980, 35/1985, 40/1990} {
    \draw (\x, -0.1) -- (\x, 0.1) node[below, font=\tiny] {\year};
}
% Événements clés - Ligne du haut (Politique/Institutions)
\draw[popblue, thick] (0, 1) node[left, align=right, font=\tiny, popblue] {Politique /\\ Institutions} -- (40, 1);
% 1955
\draw[popblue, fill=popblue!20] (0, 1.3) rectangle (1.5, 1.7) node[midway, font=\tiny, align=center] {Mai 55 :\\ Insurrection UPC};
% 1960
\draw[popblue, fill=popblue!20] (5, 1.3) rectangle (6.5, 1.7) node[midway, font=\tiny, align=center] {1er Janv 60 :\\ Indépendance};
\draw[popblue, fill=popblue!20] (5.2, 1.8) rectangle (6.7, 2.2) node[midway, font=\tiny, align=center] {Mai 60 :\\ Ahidjo Président};
% 1961
\draw[popblue, fill=popgreen!30] (6, 1.3) rectangle (7.5, 1.7) node[midway, font=\tiny, align=center] {Fév 61 :\\ Plébiscites};
\draw[popblue, fill=popgreen!30] (6.5, 1.8) rectangle (8, 2.2) node[midway, font=\tiny, align=center] {1er Oct 61 :\\ Réunification /\\ Fédération};
% 1966
\draw[popblue, fill=poporange!30] (11, 1.3) rectangle (12.5, 1.7) node[midway, font=\tiny, align=center] {1966 :\\ Parti unique (UNC)};
% 1972
\draw[popblue, fill=poppurple!30] (20, 1.3) rectangle (21.5, 1.7) node[midway, font=\tiny, align=center] {20 Mai 72 :\\ Référendum /\\ État unitaire};
% Événements clés - Ligne du bas (Conflit UPC)
\draw[poporange, thick] (0, -1) node[left, align=right, font=\tiny, poporange] {Conflit UPC /\\ Répression} -- (40, -1);
\draw[poporange, fill=poporange!20] (0, -1.4) rectangle (2, -1.7) node[midway, font=\tiny, align=center] {1955-58 :\\ Maquis Bassa /\\ Bamiléké};
\draw[poporange, fill=poporange!20] (2, -1.4) rectangle (4, -1.7) node[midway, font=\tiny, align=center] {1960 :\\ Assassinat Moumié};
\draw[poporange, fill=poporange!20] (16, -1.4) rectangle (18, -1.7) node[midway, font=\tiny, align=center] {1970-71 :\\ Exécution Ouandié /\\ Wambo (Fin maquis)};
\end{tikzpicture}
\legende{Frise chronologique : de l'insurrection UPC (1955) à la fin de la Fédération (1972). En bleu : temps forts institutionnels (indépendance, réunification, parti unique, État unitaire). En orange : chronologie de la rébellion UPC et de sa répression.}
\end{popfigure}

\begin{savaistu}[Le 1er octobre : une date mémorielle controversée]
Aujourd'hui, le 1er octobre est férié (Fête de l'Unité). Cependant, depuis la crise anglophone de 2016, des mouvements séparatistes (Ambazonia) organisent des « journées ville morte » et proclament symboliquement l'indépendance de l'ex-Cameroun méridional ce jour-là. Ils estiment que la fédération a été trahie par l'État unitaire de 1972. Cette actualité montre que la réunification de 1961 n'a pas réglé durablement la question du « vivre-ensemble » entre les deux héritages coloniaux.
\end{savaistu}

\begin{exercice}[Entraînement Probatoire : Commentaire de carte]
\emph{À l'aide de la carte ci-dessus (plébiscites 1961) et de vos connaissances, montrez que le résultat des plébiscites de février 1961 a consacré une partition du Cameroun britannique plutôt qu'une réunification intégrale.}
\end{exercice}

\begin{corrige}[Exercice 1]
\textbf{Introduction :} Présenter le contexte (ONU, tutelle britannique, deux zones). Annoncer la thèse : le vote a scindé le territoire sous tutelle britannique.
\textbf{Développement :}
\begin{enumerate}
    \item \textbf{Le Nord (Septentrional) :} Vote majoritaire pour le Nigeria (60\%). Facteurs : continuité administrative avec le Nord-Nigeria, religion musulmane, économie tournée vers le Bénoué/Ngoundéré, peur de la domination « bassa/bamiléké » du Sud. Conséquence : perte définitive pour le Cameroun (intégration juin 1961).
    \item \textbf{Le Sud (Méridional) :} Vote massif pour le Cameroun (70,5\%). Facteurs : liens ethniques/culturels avec l'Ouest camerounais (Bamiléké, Bassa), élite anglophone formée au Cameroun oriental (Foncha), hostilité à l'hégémonie ibo/yoruba du Nigeria. Conséquence : réunification le 1er octobre 1961.
\end{enumerate}
\textbf{Conclusion :} La réunification est « incomplète » (perte du Nord). Elle crée un État fédéral asymétrique (2 États fédérés au lieu de 3 envisagés initialement). Cette fracture Nord/Sud héritée du vote de 1961 pèse encore sur l'équilibre régional actuel.
\end{corrige}

\coursec{Le Cameroun sous Ahidjo : la construction de l'État unitaire (1961-1972)}

Après la proclamation de l'indépendance du Cameroun oriental (1\ier{} janvier 1960) et la réunification consacrée par le plébiscite du 11 février 1961, le jeune État camerounais naît sous une forme fragile : une République fédérale composée de deux États aux héritages coloniaux différents. Le défi majeur du président Ahmadou Ahidjo est alors de transformer ce pays bilingue, multiethnique et traversé par la rébellion upéciste, en une nation unie dotée d'un État fort. Cette section étudie comment, entre 1961 et 1972, Ahidjo construit méthodiquement un État unitaire et centralisé, en passant par trois étapes : la Constitution fédérale de Foumban, le parti unique de 1966 et le référendum du 20 mai 1972.

\courssub{La Conférence de Foumban (juillet 1961) : la naissance de la République fédérale}

Du 17 au 21 juillet 1961, les délégations de la République du Cameroun (orientale, francophone) et du Southern Cameroons (occidental, anglophone) se réunissent à \cle{Foumban}, dans la région de l'Ouest, pour élaborer la Constitution du nouvel État réunifié. La conférence est dominée par la délégation d'Ahidjo : le projet de constitution présenté est largement inspiré de la Constitution française de 1958 et du modèle de la République du Cameroun déjà indépendante.

\begin{definition}[État fédéral]
Un \cle{État fédéral} est un État composé de plusieurs États fédérés qui disposent de compétences propres (parlement, gouvernement, administration locales) tout en reconnaissant l'autorité d'un État fédéral central pour les domaines communs (défense, monnaie, affaires étrangères).
\end{definition}

La Constitution adoptée entre en vigueur le \cle{1\ier{} octobre 1961} : naît la \textbf{République fédérale du Cameroun}, composée de deux États fédérés :
\begin{itemize}
\item l'\textbf{État fédéré du Cameroun oriental} (francophone, capitale Yaoundé) ;
\item l'\textbf{État fédéré du Cameroun occidental} (anglophone, capitale Buea).
\end{itemize}

Au sommet, un pouvoir fédéral : un président de la République (Ahmadou Ahidjo), un vice-président (John Ngu Foncha, dirigeant du Southern Cameroons), une Assemblée nationale fédérale et un gouvernement fédéral.

\begin{attention}[Déséquilibre institutionnel du fédéralisme]
Le fédéralisme camerounais est \emph{asymétrique}. L'État oriental représente environ les quatre cinquièmes du territoire et de la population : il pèse donc lourdement plus que l'État occidental. De plus, le président fédéral concentre des pouvoirs considérables (nomination, promulgation des lois, état d'urgence), ce qui marginalise rapidement les institutions fédérées. Dès 1961, le fédéralisme camerounais est donc moins un partage du pouvoir qu'une transition vers la centralisation.
\end{attention}

\courssub{Le 1\ier{} septembre 1966 : la création de l'UNC, parti unique}

La construction de l'unité passe d'abord par le terrain politique. En 1961, le Cameroun compte plusieurs partis : l'UC d'Ahidjo à l'Est, le KNDP de Foncha et le CPNC à l'Ouest, sans oublier l'UPC, interdite et en rébellion. Ahidjo mène une politique habile de fusion et d'absorption : dès 1962, plusieurs petits partis rejoignent l'UC. Le \cle{1\ier{} septembre 1966}, lors d'un congrès à Yaoundé, l'Union camerounaise (UC), le KNDP et les derniers partis restants fusionnent pour former l'\textbf{Union nationale camerounaise (UNC)}. Le Cameroun devient de fait un État à \cle{parti unique}.

\begin{definition}[Parti unique]
Un \cle{parti unique} désigne un système politique dans lequel un seul parti est autorisé et confisque le pouvoir : toute opposition organisée est interdite, et le parti encadre l'État, l'administration et la société (jeunesse, syndicats, femmes).
\end{definition}

Ahidjo justifie le parti unique par plusieurs arguments, qu'il développe dans ses discours et dans son ouvrage \emph{Contribution à la construction nationale} (1964) :
\begin{enumerate}
\item \textbf{L'unité nationale} : dans un pays de plus de 200 groupes ethniques et de deux héritages linguistiques, la compétition entre partis risque, selon lui, de raviver les clivages régionaux et ethniques et de déchirer le pays.
\item \textbf{Le développement} : un parti unique rassemblerait toutes les énergies derrière un projet commun, au lieu de les disperser dans des querelles politiciennes.
\item \textbf{La stabilité} : face à la rébellion upéciste et aux coups d'État qui secouent l'Afrique des années 1960, le parti unique garantirait l'ordre et la paix, conditions du progrès.
\end{enumerate}

\begin{exemplebox}[Un raisonnement à interroger]
L'argument d'Ahidjo est un raisonnement par la finalité : la fin (unité, développement) justifierait le moyen (suppression du pluralisme). L'élève doit savoir le présenter fidèlement, puis le discuter : le parti unique a effectivement assuré la stabilité, mais au prix de l'étouffement du débat démocratique, de l'arrestation des opposants et de la confusion entre le parti et l'État.
\end{exemplebox}

\courssub{Le référendum du 20 mai 1972 : la fin du fédéralisme}

Dernière étape de la construction unitaire : la suppression du fédéralisme lui-même. Le \cle{20 mai 1972}, un référendum est organisé sur l'ensemble du territoire pour approuver une nouvelle Constitution instituant la \textbf{République unie du Cameroun}. Les résultats officiels annoncent environ \num{99,99} \% de « oui ». Les deux États fédérés disparaissent ; le pays est désormais divisé en provinces administrées par des gouverneurs nommés par le président. La date du 20 mai devient la fête nationale du Cameroun.

\begin{exoresolu}[Analyser un résultat référendaire]
Énoncé : le référendum du 20 mai 1972 recueille officiellement 99,99 \% de suffrages favorables. Que signifie un tel score ? Peut-on le prendre pour argent comptant ?
\tcblower
Solution : un score de 99,99 \% est statistiquement exceptionnel dans un scrutin libre, où les opinions divergent naturellement. Il s'explique par le contexte : parti unique depuis 1966, absence d'opposition légale, encadrement de la population par l'UNC, campagne menée uniquement pour le « oui ». Ce résultat exprime donc moins un consensus spontané qu'un plébiscite organisé par un régime qui contrôle l'ensemble du processus électoral. Il traduit l'achèvement de la concentration du pouvoir entre les mains d'Ahidjo. L'historien doit donc croiser ce chiffre officiel avec le contexte politique pour l'interpréter correctement.
\end{exoresolu}

\begin{popfigure}
\begin{center}
\begin{tikzpicture}[scale=0.75, transform shape,
  box/.style={draw=popblue, fill=popblueL, rounded corners, minimum height=0.9cm, text width=3.6cm, align=center, font=\small},
  boxo/.style={draw=poporange, fill=poporangeL, rounded corners, minimum height=0.9cm, text width=3.6cm, align=center, font=\small},
  arr/.style={-{Stealth[length=2.5mm]}, thick, popdark}]
% État fédéral 1961
\node[box, fill=popblue!40, text width=4.2cm] (pf) at (0,3.4) {\textbf{Président fédéral}\\ (Ahidjo) + Vice-président (Foncha)};
\node[box, align=center] (eo) at (-3.2,1.2){\textbf{État fédéré oriental}\\ Assemblée + gouvernement (Yaoundé)};
\node[box, align=center] (eoc) at (3.2,1.2){\textbf{État fédéré occidental}\\ Assemblée + gouvernement (Buea)};
\draw[arr] (pf) -- (eo);
\draw[arr] (pf) -- (eoc);
\node[font=\bfseries\color{popblue}] at (0,4.6) {République fédérale (1961)};
% Flèche de transition
\draw[arr, poporange, line width=2pt] (5.6,2.3) -- (7.4,2.3) node[midway, above, font=\small\bfseries, text=poporange]{20 mai 1972};
% État unitaire 1972
\node[boxo, fill=poporange!40, text width=4.2cm] (pu) at (10.4,3.4) {\textbf{Président de la République}\\ pouvoir concentré (Yaoundé)};
\node[boxo, align=center] (gov) at (10.4,1.2){\textbf{Gouverneurs de provinces}\\ nommés par le président};
\node[boxo, align=center] (pref) at (10.4,-0.8){\textbf{Préfets et sous-préfets}\\ administration déconcentrée};
\draw[arr] (pu) -- (gov);
\draw[arr] (gov) -- (pref);
\node[font=\bfseries\color{poporange}] at (10.4,4.6) {République unie (1972)};
\end{tikzpicture}
\end{center}
\legende{De l'État fédéral (1961) à l'État unitaire (1972) : en 1961, deux États fédérés dotés d'institutions propres ; en 1972, une chaîne de commandement unique, du président aux sous-préfets.}
\end{popfigure}

\begin{popfigure}
\begin{center}
\begin{tikzpicture}
\begin{axis}[
  width=0.55\linewidth,
  xbar, xmin=0, xmax=100,
  symbolic y coords={Non, Oui},
  ytick=data,
  xlabel={Part des suffrages exprimés (\%)},
  bar width=0.9cm,
  nodes near coords, every node near coord/.append style={font=\small\bfseries},
  xlabel style={font=\small},
]
\addplot[fill=popgreen] coordinates {(99.99,Oui)};
\addplot[fill=poppink] coordinates {(0.01,Non)};
\end{axis}
\end{tikzpicture}
\end{center}
\legende{Résultats officiels du référendum du 20 mai 1972 : environ 99,99 \% de « oui » à la République unie du Cameroun (source : résultats officiels). Un tel score doit être interrogé à la lumière du contexte de parti unique.}
\end{popfigure}

\courssub{La centralisation du pouvoir et l'idéologie de l'intégration nationale}

\begin{definition}[État unitaire]
Un \cle{État unitaire} est un État dans lequel le pouvoir est concentré au niveau central : il n'existe pas d'États fédérés autonomes, et les collectivités locales (provinces, départements) ne disposent que de compétences déléguées par le pouvoir central.
\end{definition}

Après 1972, la centralisation s'organise autour de trois piliers :
\begin{itemize}
\item \textbf{Le présidentialisme} : le président de la République concentre l'essentiel des pouvoirs. Il nomme les ministres, les gouverneurs, les préfets, les magistrats et les dirigeants des entreprises publiques. Le Parlement vote les lois mais ne contrôle pas réellement l'exécutif.
\item \textbf{L'administration préfectorale} : le territoire est quadrillé par des gouverneurs, préfets et sous-préfets, tous nommés et révocables par le président. Ils transmettent les décisions du centre et surveillent la population.
\item \textbf{Le contrôle des élites} : Ahidjo associe au pouvoir des notables de toutes les régions (équilibre régional dans les nominations), tout en les rendant dépendantes de sa faveur. L'UNC encadre la jeunesse, les femmes et les travailleurs à travers ses organisations satellites.
\end{itemize}

Cet appareil est légitimé par une idéologie officielle : le \cle{nationalisme camerounais}, résumé par la formule d'\textbf{intégration nationale}. Il s'agit de fondre les identités ethniques, régionales et linguistiques dans une identité camerounaise unique, autour de l'État, du parti et du président. Le « camerounais d'abord » doit primer sur l'appartenance communautaire. Cette idéologie, martelée à l'école, dans les médias d'État et lors des fêtes nationales, a permis de consolider un sentiment national réel, mais a aussi servi à disqualifier toute contestation comme « tribalisme » ou « subversion ».

\begin{attention}[Piège chronologique]
Ne pas inverser l'ordre des étapes : le passage au \textbf{parti unique} date de \textbf{1966} (création de l'UNC), tandis que l'\textbf{unité étatique} (fin du fédéralisme) date de \textbf{1972} (référendum du 20 mai). Ahidjo unifie d'abord la scène politique, puis les institutions. Inverser ces dates dans une copie est une erreur grave.
\end{attention}

\begin{methode}[Rédiger un paragraphe argumenté en Histoire]
\begin{enumerate}
\item \textbf{Idée principale} : annoncer clairement l'argument du paragraphe en une phrase (« Ahidjo construit l'unité nationale en supprimant progressivement le pluralisme »).
\item \textbf{Fait historique daté} : appuyer l'idée sur un fait précis, daté et localisé (« le 1\ier{} septembre 1966, la création de l'UNC instaure le parti unique »).
\item \textbf{Explication} : montrer le lien de cause à effet (« privé d'opposition légale, le régime peut ensuite réformer les institutions sans obstacle, ce qui prépare le référendum de 1972 »).
\item \textbf{Transition} : conclure en ouvrant sur l'idée suivante (« cette unification politique débouche naturellement sur l'unification institutionnelle »).
\end{enumerate}
\end{methode}

\begin{aretenir}[L'essentiel]
Entre 1961 et 1972, Ahidjo transforme un État fédéral fragile en État unitaire centralisé : la Constitution de Foumban (1961) crée une fédération déséquilibrée au profit du pouvoir central ; l'UNC (1\ier{} septembre 1966) instaure le parti unique au nom de l'unité nationale et du développement ; le référendum du 20 mai 1972 (99,99 \% de « oui » officiels) fonde la République unie du Cameroun. Le régime repose sur le présidentialisme, l'administration préfectorale et l'idéologie de l'intégration nationale.
\end{aretenir}

\begin{exercice}[Paragraphe argumenté]
À l'aide de la méthode ci-dessus, rédigez un paragraphe argumenté (8 à 12 lignes) répondant à la question : « Le parti unique a-t-il été un instrument d'unité nationale ou un instrument de domination personnelle d'Ahidjo ? » Vous utiliserez au moins deux faits datés entre 1961 et 1972.
\end{exercice}

\begin{corrige}[Exercice]
Éléments de correction : l'élève doit présenter les deux faces de la question. Instrument d'unité : le parti unique a mis fin aux compétitions partisanes à base régionale (fusion de l'UC et du KNDP en 1966), a permis la stabilité et préparé l'unification de 1972. Instrument de domination : l'UNC confisque le pouvoir, interdit toute opposition, encadre la société, et le référendum de 1972 (99,99 \%) se déroule sans débat possible. Une conclusion nuancée est attendue : les deux dimensions sont inséparables, l'unité nationale servant aussi à légitimer la concentration du pouvoir.
\end{corrige}

\coursec{Le bilan économique et social de l'ère Ahidjo (1960-1982)}

Dans la section précédente, nous avons vu comment Ahmadou Ahidjo a construit, entre 1961 et 1972, un État unitaire fort : réunification des deux Cameroun en 1961, parti unique en 1966, République unie du Cameroun en 1972. Mais un État ne se résume pas à ses institutions politiques. Il faut encore nourrir la population, scolariser les enfants, créer des emplois et des richesses. C'est précisément le pari d'Ahidjo : faire du Cameroun indépendant un pays qui se développe et qui compte en Afrique. Vingt-deux ans après son arrivée au pouvoir, quel bilan peut-on dresser ? Cette section évalue les réussites économiques et sociales de l'ère Ahidjo, sans oublier leurs limites, car un historien ne juge jamais en noir et blanc.

\begin{popfigure}
\begin{center}
\begin{tikzpicture}[scale=0.9, every node/.style={font=\small}]
% Frise chronologique 1960-1982
\draw[popdark, thick] (0,0) -- (16,0);
\foreach \x/\year in {0/1960, 2/1962, 4/1966, 6/1970, 8/1972, 10/1974, 12/1977, 14/1980, 16/1982} {
    \draw[popdark, thick] (\x, -0.15) -- (\x, 0.15);
    \node[below=2pt, popdark] at (\x, -0.15) {\year};
}
% Événements au-dessus
\node[above=6pt, align=center, popblue] at (0, 0.5) {Indépendance\\1er janv. 1960};
\node[above=6pt, align=center, popgreen] at (2, 0.5) {1er Plan\\1961-1965};
\node[above=6pt, align=center, poppurple] at (4, 0.5) {Parti unique\\1966};
\node[above=6pt, align=center, poporange] at (6, 0.5) {2e Plan\\1966-1970};
\node[above=6pt, align=center, poppink] at (8, 0.5) {Rép. Unie\\20 mai 1972};
\node[above=6pt, align=center, popblue] at (10, 0.5) {Transcamerounais\\1974};
\node[above=6pt, align=center, popgreen] at (12, 0.5) {Pétrole\\1977-1978};
\node[above=6pt, align=center, poppurple] at (14, 0.5) {4e Plan\\1976-1980};
\node[above=6pt, align=center, poporange] at (16, 0.5) {Démission\\Ahidjo 1982};
\end{tikzpicture}
\end{center}
\legende{Frise chronologique des grands repères de l'ère Ahidjo (1960-1982).}
\end{popfigure}

\courssub{Le libéralisme tempéré : la doctrine économique d'Ahidjo}

Au moment de l'indépendance, les jeunes États africains hésitent entre deux modèles : le socialisme planifié (choisi par le Ghana de Nkrumah ou la Guinée de Sékou Touré) et le libéralisme pur (économie de marché sans intervention de l'État). Ahidjo refuse ce dilemme. Il invente une voie originale, qu'il théorise dans son ouvrage \emph{Contribution à la construction nationale} (1964) : le \cle{libéralisme tempéré}.

\begin{definition}[Le libéralisme tempéré]
Le \cle{libéralisme tempéré} est la doctrine économique d'Ahmadou Ahidjo qui combine l'\cle{économie de marché} (propriété privée, entreprises privées, ouverture aux capitaux étrangers) et l'\cle{intervention de l'État planificateur} (l'État fixe les grandes orientations, investit dans les infrastructures et participe au capital des entreprises stratégiques). Le libéralisme est donc « tempéré », c'est-à-dire adouci, encadré, par l'action publique.
\end{definition}

L'intuition est simple : le Cameroun de 1960 est un pays pauvre, sans industrie, sans ingénieurs en nombre suffisant. Le secteur privé national est trop faible pour tout faire, mais l'État seul ne peut pas tout faire non plus. Il faut donc attirer les capitaux étrangers tout en gardant la main sur les secteurs stratégiques.

Concrètement, cette doctrine se traduit par des \cle{plans quinquennaux de développement} : tous les cinq ans, l'État publie un plan qui fixe les objectifs prioritaires (routes, écoles, barrages, industrialisation) et répartit les investissements. Les quatre plans de l'ère Ahidjo sont :
\begin{itemize}
    \item \textbf{1er Plan (1961-1965)} : infrastructures de base, éducation, santé, mise en valeur agricole.
    \item \textbf{2e Plan (1966-1970)} : consolidation, début d'industrialisation, formation des cadres.
    \item \textbf{3e Plan (1971-1975)} : industrialisation lourde, grands barrages, Transcamerounais.
    \item \textbf{4e Plan (1976-1980)} : développement intégré, pétrochimie, national-cameroonisme.
\end{itemize}

\begin{aretenir}[L'essentiel sur le libéralisme tempéré]
Doctrine économique d'Ahidjo : marché + planification. Outils : les plans quinquennaux (1961-1965, 1966-1970, 1971-1975, 1976-1980). Objectif : un développement équilibré et autocentré, sans rupture brutale avec les partenaires étrangers, notamment la France.
\end{aretenir}

\courssub{Une croissance économique remarquable : le Cameroun, « locomotive » de l'Afrique centrale}

Les résultats ne se font pas attendre. Entre 1960 et 1982, l'économie camerounaise connaît une \cle{croissance} régulière et soutenue, de l'ordre de 5 à 7 \% par an en moyenne. Le produit intérieur brut (PIB) par habitant progresse nettement, comme le montre l'histogramme ci-dessous.

\begin{popfigure}
\begin{center}
\begin{tikzpicture}
\begin{axis}[
    ybar,
    width=0.95\linewidth,
    height=7cm,
    bar width=18pt,
    ymin=0, ymax=1200,
    symbolic x coords={1960,1965,1970,1975,1980,1982},
    xtick=data,
    ylabel={PIB par habitant (dollars US courants, valeurs approchées)},
    xlabel={Année},
    nodes near coords,
    every node near coord/.append style={font=\small, rotate=90, anchor=west, yshift=2pt},
    ymajorgrids=true,
    grid style={popline!40},
    axis line style={popdark},
    enlarge x limits=0.15,
]
\addplot[fill=popblue!70, draw=popblue] coordinates {(1960,150) (1965,220) (1970,300) (1975,520) (1980,900) (1982,1050)};
\end{axis}
\end{tikzpicture}
\end{center}
\legende{Évolution du PIB par habitant du Cameroun entre 1960 et 1982 (valeurs approchées, en dollars US courants). La hausse s'accélère nettement après 1975 avec la manne pétrolière. Source : données indicatives Banque mondiale.}
\end{popfigure}

Grâce à cette dynamique, le Cameroun devient la première économie de l'Afrique centrale : on le surnomme la \cle{locomotive de l'Afrique centrale}. Son budget est équilibré, sa monnaie (le franc CFA) est stable, et le pays attire les investisseurs. À la fin des années 1970, le Cameroun fait figure de modèle de réussite économique en Afrique francophone.

\begin{savaistu}[Le Cameroun, « l'Afrique en miniature »]
Le Cameroun est souvent surnommé \emph{l'Afrique en miniature}. Pourquoi ? Parce qu'on y trouve presque tous les paysages du continent (forêt équatoriale au sud, savane au centre, sahel à l'extrême-nord, montagnes à l'ouest) et presque toutes les productions agricoles africaines : cacao, café, coton, bananes, arachides, palmier à huile, hévéa, thé... Cette diversité exceptionnelle est un atout majeur pour l'économie du pays.
\end{savaistu}

\courssub{Les grandes réalisations : infrastructures et industrie}

Le libéralisme tempéré se concrétise par des réalisations spectaculaires, encore visibles aujourd'hui.

\textbf{Le barrage hydroélectrique d'Édéa.} Construit sur le fleuve Sanaga et mis en service en 1958 (période coloniale), il est agrandi après l'indépendance (ajout de turbines). Le barrage d'Édéa fournit l'électricité indispensable à l'industrie et aux villes. L'énergie bon marché d'Édéa rend possible l'industrie lourde.

\textbf{Le complexe ALUCAM.} À Édéa même, l'usine \cle{ALUCAM} (Aluminium du Cameroun), dont la convention est signée en 1957 et qui démarre sa production la même année, transforme l'alumine importée en aluminium grâce à l'électricité du barrage. Développée sous Ahidjo (extensions successives), c'est l'une des toutes premières grandes usines d'Afrique centrale, symbole de l'industrialisation camerounaise.

\textbf{Le Transcamerounais.} C'est le chantier le plus emblématique de l'ère Ahidjo : prolonger le chemin de fer vers le nord pour désenclaver les régions septentrionales et intégrer économiquement l'ensemble du pays.

\begin{exemplebox}[Le Transcamerounais en chiffres]
Achevé en \cle{1974}, le \cle{Transcamerounais} relie \cle{Douala} à \cle{Ngaoundéré} sur environ \SI{940}{km} de voie ferrée (dont le tronçon nouveau Yaoundé--Ngaoundéré, environ \SI{622}{km}, ouvert en 1974). Coût : plusieurs dizaines de milliards de francs CFA, financés par l'État et des prêts extérieurs. Effet : le coton du Nord descend vers le port de Douala, les produits manufacturés montent vers Ngaoundéré ; les régions du Nord, autrefois isolées, sont rattachées à l'économie nationale.
\end{exemplebox}

\textbf{Le port de Douala.} Modernisé et agrandi (notamment le quai à conteneurs dans les années 1970), le port de Douala (alors géré par l'Office des Ports du Cameroun, le statut de Port Autonome date de 1999) devient le premier port d'Afrique centrale : par lui transitent non seulement les exportations camerounaises (cacao, café, bois, coton), mais aussi une partie des marchandises du Tchad et de la Centrafrique, pays enclavés.

\begin{popfigure}
\begin{center}
\begin{tikzpicture}[scale=0.9, every node/.style={font=\small}]
% Contour simplifié du Cameroun
\draw[thick, popdark, fill=popcream] (2.2,0) -- (3.4,0.2) -- (4.2,1.2) -- (4.4,2.4) -- (4.0,3.2) -- (4.3,4.2) -- (3.8,5.2) -- (3.2,6.0) -- (2.6,6.8) -- (2.0,7.6) -- (1.6,8.4) -- (1.2,7.4) -- (1.4,6.2) -- (1.0,5.2) -- (1.3,4.2) -- (0.8,3.4) -- (0.4,2.6) -- (0.2,1.6) -- (0.6,0.6) -- (1.4,0.1) -- cycle;
% Nord
\draw[->, popdark] (5.2,7.6) -- (5.2,8.4) node[above]{N};
% Océan
\node[popblue] at (0.3,0.6) {\emph{Golfe de Guinée}};
% Villes
\fill[popblue] (0.7,1.1) circle (2.5pt) node[right=2pt, popdark]{\textbf{Douala (port)}};
\fill[popgreen] (1.5,1.9) circle (2.5pt) node[right=2pt, popdark]{\textbf{Édéa (barrage + ALUCAM)}};
\fill[popink] (1.8,2.6) circle (2.5pt) node[right=2pt, popdark]{\textbf{Yaoundé}};
\fill[poporange] (2.9,5.6) circle (2.5pt) node[right=2pt, popdark]{\textbf{Ngaoundéré}};
% Transcamerounais
\draw[line width=2.2pt, poporange, decorate, decoration={zigzag, segment length=4pt, amplitude=1pt}] (0.7,1.1) -- (1.8,2.6) -- (2.9,5.6);
\node[poporange, rotate=62] at (2.1,4.3) {\textbf{Transcamerounais (1974)}};
% Sanaga
\draw[popblue, thick] (1.0,2.3) -- (1.5,1.9) -- (2.2,1.5);
\node[popblue] at (2.4,1.8) {\emph{Sanaga}};
% Légende
\node[anchor=west] at (4.6,3.4) {\textcolor{poporange}{\rule{14pt}{2.2pt}} Chemin de fer};
\node[anchor=west] at (4.6,2.9) {\textcolor{popblue}{$\bullet$} Port / \textcolor{popgreen}{$\bullet$} Industrie};
\node[anchor=west] at (4.6,2.4) {\textcolor{popink}{$\bullet$} Capitale \textcolor{poporange}{$\bullet$} Terminus};
\end{tikzpicture}
\end{center}
\legende{Croquis schématique du Cameroun : localisation des grandes réalisations économiques de l'ère Ahidjo (port de Douala, barrage d'Édéa et usine ALUCAM sur la Sanaga, ligne du Transcamerounais Douala--Yaoundé--Ngaoundéré achevée en 1974). Croquis sans prétention d'exactitude cartographique.}
\end{popfigure}

\begin{methode}[Construire un croquis géographique (méthode exigible au Bac)]
Au Baccalauréat technique, un croquis est une épreuve codifiée. Respecte toujours les étapes suivantes :
\begin{enumerate}
\item Tracer le \cle{fond de carte} : contour simplifié mais reconnaissable du territoire, avec l'orientation (flèche du nord).
\item Placer les \cle{informations essentielles} demandées par le sujet (villes, infrastructures, flux) à l'aide de symboles simples et proportionnés (points, traits, flèches d'épaisseur variable).
\item Rédiger un \cle{titre} précis au-dessus du croquis (quoi ? où ? quand ?).
\item Construire une \cle{légende ordonnée} : les symboles sont classés par thèmes et chaque symbole utilisé sur la carte figure dans la légende, et réciproquement.
\item Vérifier la cohérence : pas d'information sans symbole, pas de symbole sans légende.
\end{enumerate}
\end{methode}

\courssub{L'agriculture de rente et les sociétés de développement}

L'économie camerounaise repose d'abord sur l'agriculture. On distingue l'agriculture vivrière (mil, sorgho, manioc, plantain, destinée à nourrir les familles) et l'\cle{agriculture de rente}, c'est-à-dire les cultures destinées à l'exportation et qui rapportent des devises : le \cle{cacao} (Centre, Sud, Est), le \cle{café} (Ouest, Littoral, Nord-Ouest), le \cle{coton} (Nord, Adamaoua), les \cle{bananes} (Littoral, Sud-Ouest), mais aussi l'hévéa, le palmier à huile et le thé.

Pour encadrer les paysans et organiser la production, l'État crée ou restructure des \cle{sociétés de développement} paraétatiques. La plus célèbre est la \cle{SODECAO} (Société de développement du cacao), créée en 1974 : elle distribue les plants et les pesticides, encadre techniquement les producteurs et stabilise les prix. Sur le même modèle existent la \cle{SODECOTON} (coton, Nord, créée 1974), la \cle{SODEPALM} (palmier à huile), la \cle{CDC} (Cameroon Development Corporation, bananes et palmiers, région du Sud-Ouest, héritée de l'époque coloniale mais confortée), ou encore l'\cle{ONCPB} (Office National de Commercialisation des Produits de Base, créé 1977) pour la commercialisation et la stabilisation des prix. Grâce à ces structures, le Cameroun devient l'un des premiers producteurs mondiaux de cacao et de café.

\courssub{Le pétrole et le « national-cameroonisme » économique}

Deux événements marquent la fin de l'ère Ahidjo.

\textbf{Le début de l'exploitation pétrolière.} En \cle{1977-1978}, la compagnie ELF-SEREPCA commence l'exploitation du pétrole au large de la côte camerounaise (gisements du Rio del Rey, près de Limbé, puis Kome-Kribi). Le pétrole devient rapidement la première source de devises du pays et doppe la croissance : c'est lui qui explique l'envolée du PIB par habitant visible sur l'histogramme après 1975. Mais Ahidjo choisit de garder les revenus pétroliers partiellement « hors budget », dans des comptes spéciaux, ce qui manquera de transparence.

\textbf{Le « national-cameroonisme ».} Ahidjo veut que les Camerounais participent davantage à l'économie dominée par les capitaux étrangers. Sa politique du \cle{national-cameroonisme} économique (lancée officiellement en 1976) prévoit la participation de l'État et des nationaux au capital des entreprises (souvent à hauteur de 30 à 50 \%), la formation de cadres camerounais et la « camerounisation » des postes de direction. L'État devient ainsi actionnaire dans de nombreuses sociétés (banques, industries, transports).

\begin{exoresolu}[Calculer et interpréter un taux de croissance]
Énoncé : le PIB par habitant du Cameroun passe d'environ 150 dollars en 1960 à environ 1050 dollars en 1982. 1) Calcule le taux de variation sur la période. 2) Propose une explication historique de cette progression.
\tcblower
Solution détaillée :
1) Le taux de variation se calcule par la formule :
\[ t = \frac{V_A - V_D}{V_D} \times 100 = \frac{1050 - 150}{150} \times 100 = \frac{900}{150} \times 100 = 600\,\%. \]
Le PIB par habitant a donc été multiplié par 7 (augmentation de 600 \%) entre 1960 et 1982.
2) Explication historique : cette progression s'explique par la stabilité politique retrouvée après 1971, la politique du libéralisme tempéré (plans quinquennaux), les bonnes recettes de l'agriculture de rente (cacao, café, coton), les grandes infrastructures (Transcamerounais, port de Douala) et, à partir de 1977-1978, les revenus du pétrole. Attention toutefois : il s'agit de valeurs en dollars courants, qui intègrent aussi l'inflation mondiale des années 1970.
\end{exoresolu}

\courssub{Les progrès sociaux : éducation, santé, urbanisation}

Le développement économique s'accompagne de progrès sociaux sensibles.

\textbf{L'éducation.} La scolarisation progresse fortement : le taux brut de scolarisation au primaire passe d'environ 60 \% au début des années 1960 à plus de 90 \% au début des années 1980, faisant du Cameroun l'un des pays les plus scolarisés d'Afrique. Des collèges et lycées sont créés dans toutes les provinces, et l'Université fédérale du Cameroun (créée en 1962 à Yaoundé) forme les cadres de la nation.

\textbf{La santé.} Le réseau sanitaire se développe : hôpitaux provinciaux, centres de santé dans les arrondissements, campagnes de vaccination (variole, rougeole, BCG). L'espérance de vie progresse (de ~40 ans en 1960 à ~50 ans en 1982) et la mortalité infantile recule, même si les campagnes restent mal desservies.

\textbf{L'urbanisation.} Conséquence de la croissance et de l'exode rural, les villes explosent : \cle{Douala} passe d'environ 200 000 habitants en 1960 à plus de 600 000 en 1982 (recensement 1976 : ~450 000), et \cle{Yaoundé} suit le même rythme. Cette urbanisation rapide crée des emplois mais aussi des quartiers spontanés sans eau ni électricité.

\courssub{Les limites du bilan : un développement inachevé et inégal}

Un bilan d'historien doit être nuancé. Quatre limites majeures ternissent le tableau.

\textbf{Les inégalités régionales.} Le développement profite d'abord au Littoral et au Centre (Douala, Yaoundé, Édéa), tandis que l'Ouest anglophone se sent marginalisé et que l'Extrême-Nord reste en retard malgré le Transcamerounais. Les inégalités entre villes et campagnes se creusent.

\textbf{La dépendance extérieure.} L'économie reste dépendante des cours mondiaux des matières premières (cacao, café) et des capitaux, techniciens et marchés étrangers, notamment français. Le Cameroun exporte des produits bruts et importe des produits manufacturés : il ne maîtrise pas sa propre croissance.

\textbf{La corruption naissante.} La concentration des pouvoirs et l'opacité des revenus pétroliers favorisent les détournements et le clientélisme au sein de l'administration et du parti unique.

\textbf{L'absence de libertés politiques.} Le développement économique s'accompagne d'un régime autoritaire : parti unique, censure, répression de toute opposition (ex. : exécution d'Ernest Ouandié en 1971, arrestation d'Albert Mukong). Le « contrat social » d'Ahidjo peut se résumer ainsi : le développement en échange du silence politique.

\begin{attention}[Piège de rédaction au Bac]
Ne présente jamais le bilan d'Ahidjo comme uniquement positif (« le Cameroun était un paradis ») ou uniquement négatif (« tout était mauvais »). Le Baccalauréat valorise l'\cle{analyse nuancée et documentée} : cite des réalisations précises et chiffrées (Transcamerounais 1974, SODECAO, pétrole 1977-1978), puis des limites précises (inégalités régionales, dépendance, autoritarisme), et conclus par un jugement équilibré. Une copie manichéenne perd des points, même si ses informations sont exactes.
\end{attention}

\begin{aretenir}[Bilan de l'ère Ahidjo en quatre points]
\begin{itemize}
\item \textbf{Économie} : libéralisme tempéré, plans quinquennaux, croissance régulière, Cameroun « locomotive » de l'Afrique centrale ; grandes réalisations (Édéa, ALUCAM, Transcamerounais 1974, port de Douala) ; pétrole dès 1977-1978.
\item \textbf{Agriculture} : cultures de rente encadrées par les sociétés de développement (SODECAO, SODECOTON, ONCPB).
\item \textbf{Social} : scolarisation massive, santé en progrès, urbanisation rapide de Douala et Yaoundé.
\item \textbf{Limites} : inégalités régionales, dépendance extérieure, corruption naissante, absence de libertés politiques.
\end{itemize}
\end{aretenir}

Ce bilan en demi-teinte est l'héritage que reçoit Paul Biya lorsqu'Ahidjo lui transmet le pouvoir le 6 novembre 1982. C'est précisément cette succession, et le nouveau cours qu'elle imprime au pays, que nous étudierons dans la section suivante.

\coursec{La succession de 1982 et l'avènement de Paul Biya}

Après vingt-deux années d'un pouvoir sans partage, l'ère Ahidjo s'achève non par une révolution ni par un coup d'État, mais par une démission. La section précédente a montré un bilan économique et social en demi-teinte : croissance réelle mais inégale, État fort mais autoritaire, unité construite mais au prix d'un parti unique et d'une répression. C'est au sommet de cet État hyper-centralisé que se joue, le 4 novembre 1982, l'un des événements les plus surprenants de l'histoire du Cameroun indépendant : le départ volontaire d'Ahmadou Ahidjo et l'arrivée au pouvoir de Paul Biya. Comprendre cette succession, sa rupture ultérieure et ses conséquences, c'est comprendre les fondements du régime camerounais contemporain.

\courssub{Le 4 novembre 1982 : une démission surprise}

\textbf{Les faits}\par

Le \cle{4 novembre 1982}, Ahmadou Ahidjo, premier président de la République du Cameroun, au pouvoir depuis 1960 (soit 22 ans), annonce à la radio sa \cle{démission}. La nouvelle stupéfie le pays : rien, dans l'actualité politique, ne laissait prévoir un tel départ. Ahidjo, âgé de 58 ans, contrôle alors totalement l'appareil d'État et le parti unique, l'UNC (Union nationale camerounaise).

Conformément à la \cle{Constitution} de 1972, qui prévoit qu'en cas de vacance de la présidence de la République le Premier ministre assure l'intérim puis la succession, \cle{Paul Biya}, \cle{Premier ministre depuis 1975}, devient président de la République. Il prête serment le 6 novembre 1982. La transmission du pouvoir se déroule donc dans le calme, sans violence et sans interruption des institutions.

\begin{definition}[Succession constitutionnelle]
La \cle{succession constitutionnelle} est la transmission du pouvoir selon les règles prévues par la Constitution, sans interruption du régime. Elle se distingue du coup d'État (prise du pouvoir par la force, hors du cadre légal) et de la révolution (renversement de l'ordre politique par le soulèvement populaire). Au Cameroun, la Constitution de 1972 désigne le Premier ministre comme successeur en cas de vacance présidentielle.
\end{definition}

\begin{attention}[Ce que la démission de 1982 n'est pas]
La démission d'Ahidjo est \textbf{volontaire} et \textbf{constitutionnelle}. Ce n'est ni un coup d'État ni une révolution : aucune force armée n'a contraint le président à partir, aucune loi n'a été suspendue, les institutions ont continué de fonctionner normalement. Dans une copie, ne jamais écrire que Biya a « pris le pouvoir » : il l'a \emph{reçu} selon la procédure légale. Confondre succession constitutionnelle et coup d'État est une erreur grave d'analyse historique.
\end{attention}

\textbf{Les raisons de la démission}\par

Dans son message à la nation, Ahidjo invoque des raisons personnelles : la \cle{fatigue} accumulée après 22 ans d'exercice du pouvoir et des problèmes de \cle{santé} (il évoque une « fatigue physique et nerveuse » constatée par ses médecins lors d'un séjour en France).

Mais les historiens avancent aussi des \cle{hypothèses politiques} :
\begin{itemize}
\item Ahidjo aurait pensé pouvoir \textbf{contrôler le pouvoir en coulisses} : en restant président du parti unique (l'UNC), il imaginait que Biya, réputé discret et loyal, gouvernerait sous sa tutelle, selon le modèle d'un « président du parti » dominant un « président de la République » docile ;
\item des \textbf{tensions internes} au régime et des rivalités entre clans auraient rendu sa position plus inconfortable qu'il n'y paraissait ;
\item certains évoquent des \textbf{pressions ou des conseils de l'entourage} (médecins, proches, réseaux français) l'incitant à un départ anticipé.
\end{itemize}

L'hypothèse de la « télécommande » est la plus répandue : Ahidjo aurait voulu organiser une succession apparente tout en conservant le pouvoir réel. La suite des événements montrera que ce calcul était erroné.

\begin{popfigure}
\begin{center}
\begin{tikzpicture}[scale=0.9, every node/.style={font=\small}]
% Frise chronologique nov 1982 - mars 1985
\draw[thick, popdark, -{Stealth}] (0,0) -- (14.5,0) node[right]{temps};
\foreach \x/\d/\lab in {
  0.8/4 nov. 1982/Démission d'Ahidjo ; Biya devient président (6 nov.),
  3.2/22 août 1983/Rupture : Biya révèle un « complot » ; Ahidjo démissionne de l'UNC,
  5.8/28 fév. 1984/Ahidjo condamné à mort par contumace (peine commuée),
  8.4/6 avril 1984/Tentative de coup d'État de la Garde républicaine : échec,
  10.8/1984/Le pays redevient « République du Cameroun »,
  13.2/24 mars 1985/Congrès de Bamenda : l'UNC devient le RDPC}
{
  \draw[popblue, thick] (\x,0) -- (\x,0.25);
  \fill[poporange] (\x,0) circle (2.5pt);
  \node[above, align=center, text width=2.1cm] at (\x,0.3) {\textbf{\d}\\ \lab};
}
\end{tikzpicture}
\end{center}
\legende{Frise chronologique de la succession et de la rupture (novembre 1982 -- mars 1985). En moins de trois ans, le Cameroun passe d'une passation de pouvoir pacifique à une rupture violente entre Ahidjo et Biya, puis à une refonte des symboles du régime (nom du pays, nom du parti).}
\end{popfigure}

\courssub{La rupture Ahidjo--Biya (1983-1984)}

\textbf{De la succession à l'affrontement}\par

Dès les premiers mois, le scénario imaginé par Ahidjo échoue. Paul Biya entend exercer \textbf{pleinement} le pouvoir présidentiel et refuse la tutelle de l'ancien président, resté à la tête de l'UNC. La coexistence de deux « chefs » — le président de la République et le président du parti unique — devient intenable.

Au cours de l'année 1983, la tension monte : Ahidjo, depuis l'étranger où il séjourne de plus en plus, multiplie les critiques et tente de faire prévaloir la suprématie du parti sur l'État. En \cle{août 1983}, Biya dénonce publiquement un \cle{complot présumé} fomenté par Ahidjo et ses partisans. Ahidjo démissionne alors de la présidence de l'UNC et s'installe durablement en exil, en France puis au \cle{Sénégal}.

En février 1984, un tribunal camerounais condamne Ahidjo à mort \cle{par contumace} (c'est-à-dire en son absence, l'accusé n'ayant pas comparu) pour complot contre la sûreté de l'État. La peine sera ensuite commuée en détention à perpétuité, mais l'ancien président ne remettra jamais les pieds au Cameroun.

\begin{methode}[Analyser un événement politique]
Pour analyser rigoureusement un événement politique comme la rupture de 1983, suivre quatre étapes :
\begin{enumerate}
\item \textbf{Les faits établis} : que s'est-il passé, quand, où ? (Démission le 4 novembre 1982, investiture de Biya le 6 novembre, rupture publique en août 1983.)
\item \textbf{Les acteurs} : qui agit, avec quels moyens et quels intérêts ? (Ahidjo, Biya, le parti UNC, l'armée, les réseaux fidèles à chacun.)
\item \textbf{Les causes} : distinguer les causes \emph{immédiates} (le refus de Biya d'obéir, la dénonciation du complot) des causes \emph{profondes} (l'ambiguïté d'une succession où l'ancien chef garde le parti, la concentration des pouvoirs, les rivalités de réseaux).
\item \textbf{Les conséquences} : à court terme (exil d'Ahidjo, épuration de l'appareil) et à long terme (consolidation du pouvoir de Biya, changement des symboles du régime).
\end{enumerate}
\end{methode}

\textbf{La tentative de coup d'État du 6 avril 1984}\par

La rupture atteint son paroxysme avec la \cle{tentative de coup d'État du 6 avril 1984}. Des éléments de la \cle{Garde républicaine}, l'unité d'élite chargée de la protection du président et composée en majorité de militaires originaires du Nord (la région d'Ahidjo), attaquent le palais présidentiel et tentent de renverser Paul Biya.

Après des combats à Yaoundé, les forces loyalistes reprennent le contrôle de la capitale. Le putsch échoue. Les putschistes présumés sont arrêtés ; plusieurs officiers sont jugés et exécutés. L'attribution réelle du coup reste débattue : Biya accuse les partisans d'Ahidjo, ce que l'ancien président, depuis Dakar, dément formellement.

\begin{exemplebox}[Le coût humain du 6 avril 1984]
La tentative de coup d'État d'avril 1984 fait \textbf{plusieurs dizaines de morts} dans les combats à Yaoundé (les bilans varient selon les sources, de quelques dizaines à plus d'une centaine de victimes). Dans les semaines qui suivent, plusieurs officiers reconnus coupables de participation au putsch sont \textbf{exécutés} après leur condamnation par un tribunal militaire. Cet épisode est le dernier affrontement armé majeur pour le pouvoir au Cameroun indépendant.
\end{exemplebox}

\begin{savaistu}[Un exil sans retour]
Ahmadou Ahidjo meurt en exil à \textbf{Dakar (Sénégal)} le 30 novembre 1989, sans être jamais revenu au Cameroun. Sa dépouille n'y a toujours pas été rapatriée : le premier président du Cameroun indépendant repose au cimetière de Yoff, à Dakar. Le retour de sa dépouille reste un sujet récurrent du débat public camerounais.
\end{savaistu}

\courssub{Le « Renouveau » : nouveaux symboles, nouvelles promesses}

\textbf{Changer les symboles du régime}\par

Pour marquer la rupture avec l'ère Ahidjo et affirmer son autorité, Paul Biya modifie les symboles de l'État et du parti :
\begin{itemize}
\item par la loi du \cle{4 février 1984}, le pays abandonne le nom de « République unie du Cameroun » (hérité du fédéralisme aboli en 1972) et redevient la \cle{République du Cameroun} ; le drapeau et l'hymne sont également modifiés (retour au drapeau tricolore avec l'étoile unique) ;
\item au congrès de \cle{Bamenda}, les 21-24 \cle{mars 1985}, l'UNC est dissoute et remplacée par le \cle{RDPC} (Rassemblement démocratique du peuple camerounais), nouveau parti unique dont Biya est le président.
\end{itemize}

Ces changements ont une double fonction : effacer l'héritage d'Ahidjo (l'UNC était son parti) et légitimer le nouveau chef en lui donnant ses propres institutions.

\textbf{Les promesses du Renouveau}\par

Paul Biya baptise son projet politique le \cle{Renouveau}. Il promet trois grandes orientations :
\begin{enumerate}
\item la \cle{rigueur} : discipline dans la gestion de l'État, lutte contre le laxisme administratif ;
\item la \cle{moralisation} : combat contre la corruption et les détournements de fonds publics ;
\item la \cle{libéralisation} : assouplissement du régime, promesse d'une démocratisation progressive et d'un État de droit, résumée par le slogan du « \emph{libéralisme intégral et communautaire} ».
\end{enumerate}

Ces promesses suscitent un réel espoir dans la population, d'autant que Biya libère des détenus politiques et adopte un style plus ouvert que celui, très autoritaire, d'Ahidjo. Mais les pratiques du pouvoir — parti unique, centralisation, répression des opposants — demeureront largement inchangées, comme le montrera la section suivante.

\begin{popfigure}
\begin{center}
\begin{tikzpicture}[scale=1, every node/.style={font=\small}]
% Schéma de la passation de pouvoir de novembre 1982
\node[draw, rounded corners, fill=popblueL, text width=4.2cm, align=center, minimum height=2.2cm] (ahidjo) at (0,0) {\textbf{Ahmadou Ahidjo}\\ Président de la République\\ (1960-1982)\\ \emph{démissionne le 4 nov. 1982}};
\node[draw, rounded corners, fill=popgreenL, text width=4.2cm, align=center, minimum height=2.2cm] (biya) at (10.5,0) {\textbf{Paul Biya}\\ Premier ministre\\ (depuis 1975)\\ \emph{investi président le 6 nov. 1982}};
\node[draw, rounded corners, fill=popgoldL, text width=3.6cm, align=center] (constit) at (5.25,2.6) {\textbf{Constitution de 1972}\\ vacance de la présidence\\ $\rightarrow$ le Premier ministre succède};
\draw[-{Stealth}, very thick, poporange] (ahidjo.east) -- (biya.west) node[midway, below, align=center, text width=3cm]{transmission pacifique et légale du pouvoir};
\draw[-{Stealth}, thick, popdark] (constit.south west) -- (ahidjo.north east);
\draw[-{Stealth}, thick, popdark] (constit.south east) -- (biya.north west);
\node[below, align=center, text width=4cm] at (0,-1.6) {\small 22 ans de pouvoir ; parti unique UNC ; État unitaire depuis 1972};
\node[below, align=center, text width=4cm] at (10.5,-1.6) {\small continuité institutionnelle ; programme du « Renouveau »};
\end{tikzpicture}
\end{center}
\legende{Schéma commenté de la passation de pouvoir de novembre 1982 (à la manière d'une légende de photographie). Au centre, la Constitution de 1972 sert de cadre juridique : c'est elle qui désigne le Premier ministre Paul Biya comme successeur d'Ahmadou Ahidjo. La flèche centrale souligne le caractère pacifique et légal de la transmission, qui la distingue d'un coup d'État.}
\end{popfigure}

\courssub{Exercice résolu et bilan}

\begin{exoresolu}[Analyser la succession de 1982]
\textbf{Énoncé.} Un élève écrit : « En 1982, Paul Biya a pris le pouvoir par un coup d'État, puis Ahidjo est parti en exil. » Relevez les erreurs de cette phrase et rédigez une version correcte en appliquant la méthode d'analyse d'un événement politique (faits, acteurs, causes, conséquences).
\tcblower
\textbf{Solution détaillée.}
\begin{enumerate}
\item \textbf{Erreur de nature juridique} : il n'y a pas eu de coup d'État en 1982. La démission d'Ahidjo (4 novembre 1982) est volontaire, et l'accession de Biya à la présidence (6 novembre 1982) s'est faite conformément à la Constitution de 1972, qui prévoit la succession par le Premier ministre. Il s'agit d'une \cle{succession constitutionnelle}.
\item \textbf{Erreur de chronologie et de causalité} : l'exil d'Ahidjo n'est pas la conséquence directe de l'arrivée de Biya au pouvoir, mais de la \emph{rupture} entre les deux hommes en 1983 (complot présumé dénoncé en août 1983, démission d'Ahidjo de l'UNC, condamnation à mort par contumace en février 1984).
\item \textbf{Version correcte} : « Le 4 novembre 1982, Ahmadou Ahidjo démissionne volontairement de la présidence de la République, invoquant la fatigue et la santé ; conformément à la Constitution, le Premier ministre Paul Biya lui succède. La rupture entre les deux hommes, en 1983-1984, conduit Ahidjo à l'exil au Sénégal, où il meurt en 1989. »
\end{enumerate}
\end{exoresolu}

\begin{aretenir}[L'essentiel de la section]
\begin{itemize}
\item \textbf{4 novembre 1982} : démission surprise d'Ahidjo après 22 ans de pouvoir ; \textbf{Paul Biya}, Premier ministre depuis 1975, devient président conformément à la Constitution de 1972 : c'est une succession constitutionnelle, pas un coup d'État.
\item Raisons officielles : fatigue et santé ; hypothèse politique majeure : Ahidjo pensait garder le pouvoir réel via le parti unique.
\item \textbf{1983-1984} : rupture Ahidjo--Biya ; complot présumé, condamnation d'Ahidjo à mort par contumace, exil au Sénégal (il meurt à Dakar en 1989).
\item \textbf{6 avril 1984} : tentative de coup d'État de la Garde républicaine ; échec, dizaines de morts, officiers exécutés.
\item Nouveaux symboles : \textbf{République du Cameroun} (1984) et \textbf{RDPC} (congrès de Bamenda, mars 1985).
\item Le \textbf{Renouveau} promet rigueur, moralisation et libéralisation.
\end{itemize}
\end{aretenir}

\begin{exercice}[Pour s'entraîner]
\begin{enumerate}
\item Définis la succession constitutionnelle et explique en quoi l'événement de novembre 1982 en est une illustration.
\item Construis un tableau à deux colonnes distinguant les causes immédiates et les causes profondes de la rupture Ahidjo--Biya.
\item Rédige un paragraphe argumenté (10 à 15 lignes) répondant à la question : « En quoi les années 1982-1985 marquent-elles à la fois une continuité et une rupture dans l'histoire politique du Cameroun ? »
\end{enumerate}
\end{exercice}

\begin{corrige}[Exercice — éléments de réponse]
\begin{enumerate}
\item La succession constitutionnelle est la transmission du pouvoir selon les règles prévues par la Constitution, sans interruption du régime. En novembre 1982, la démission volontaire d'Ahidjo entraîne l'application de la Constitution de 1972 : le Premier ministre Paul Biya devient président, sans violence ni suspension des institutions.
\item Causes immédiates : refus de Biya d'obéir à Ahidjo, dénonciation du complot d'août 1983. Causes profondes : ambiguïté d'une succession où l'ancien président conserve le parti unique, concentration des pouvoirs, rivalités de réseaux politiques et régionaux.
\item Continuité : même régime issu du parti unique, même personnel politique (Biya fut Premier ministre d'Ahidjo), même État unitaire centralisé. Rupture : changement de chef après 22 ans, élimination politique d'Ahidjo, nouveaux symboles (République du Cameroun, RDPC) et nouveau discours (Renouveau : rigueur, moralisation, libéralisation).
\end{enumerate}
\end{corrige}

\coursec{Le Cameroun de 1982 à nos jours : mutations politiques}

Après la démission surprise du président Ahmadou Ahidjo le 4 novembre 1982, son Premier ministre, Paul Biya, assure l'intérim puis est élu président de la République le 14 janvier 1984. Les premières années (1982-1986) sont marquées par le discours de la « \cle{rigueur} » et de la « \cle{moralisation} » de la vie publique, ainsi que par la consolidation du pouvoir face à la tentative de coup d'État de la garde républicaine en avril 1984 (procès d'Ahidjo par contumace). Mais dès 1986, le pays bascule dans une crise économique majeure qui va reconfigurer l'espace politique et social.

\courssub{La succession de 1982 et les premières années du pouvoir Biya (1982-1986)}

\textbf{Une transition constitutionnelle.} Le 4 novembre 1982, le président Ahidjo démissionne pour raisons de santé. Conformément à la Constitution, le Premier ministre Paul Biya devient président par intérim. Il est élu au suffrage universel le 14 janvier 1984 avec 99,98 \% des voix (candidat unique du parti unique, le \cle{RDPC}, né en 1985 de la transformation de l'UNC).

\textbf{La « rigueur » et la « moralisation ».} Le nouveau chef de l'État lance une campagne contre la corruption et la gabegie, fait arrêter d'anciens dignitaires et réduit le train de vie de l'État. Il remplace progressivement les cadres ahidjistes par ses proches.

\textbf{Le coup d'État manqué d'avril 1984.} Du 4 au 6 avril 1984, des éléments de la garde républicaine (majoritairement originaires du Nord) tentent de renverser le pouvoir. L'armée régulière, appuyée par la population de Yaoundé, mater la rébellion en trois jours. Bilan : plusieurs centaines de morts. Ahidjo, exilé en France, est condamné à mort par contumace (peine commuée en détention à perpétuité). Cet événement scelle l'autorité de Biya et accélère la « \cle{camerounisation} » des cadres de l'armée et de l'administration.

\begin{aretenir}[Les jalons 1982-1986]
\begin{itemize}
\item 4 nov. 1982 : démission d'Ahidjo ; 14 janv. 1984 : élection de Paul Biya.
\item 1984 : transformation de l'UNC en RDPC (24 mars) ; tentative de coup d'État (avril).
\item 1985-1986 : apogée des recettes pétrolières ; début de la chute des cours (1986).
\end{itemize}
\end{aretenir}

\courssub{La crise économique de 1986-1993 et l'ajustement structurel}

\textbf{Une crise venue des matières premières.} À partir de 1986, le Cameroun subit de plein fouet l'effondrement des cours mondiaux de ses quatre piliers d'exportation : le pétrole, le café, le cacao et le coton. Les recettes budgétaires chutent de plus de 50 \% en quelques années. Les entreprises publiques ferment, les arriérés de salaire s'accumulent, l'État ne peut plus honorer sa dette extérieure.

\textbf{Le plan de redressement « maison » (1987-1988).} Avant de céder aux injonctions extérieures, le gouvernement tente un plan d'austérité interne : gel des recrutements, suppression de primes, réduction des subventions. L'échec de ce plan contraint le Cameroun à négocier avec le FMI et la Banque mondiale.

\begin{definition}[L'ajustement structurel]
Un \cle{ajustement structurel} est un programme de réformes économiques imposé par le \cle{FMI} et la \cle{Banque mondiale} aux États endettés, en échange de nouveaux prêts. Il comporte généralement : la réduction drastique des dépenses publiques (fonctionnaires, investissements), la privatisation des entreprises d'État, la libéralisation des prix et du commerce, et la suppression des subventions aux produits de base.
\end{definition}

Le Cameroun signe son premier accord avec le FMI en \cle{juillet 1988}. Les plans successifs (1988, 1990, 1991) imposent : licenciements massifs dans la fonction publique et les entreprises publiques (\cle{SONEL}, \cle{CAMAIR}, \cle{SOCATUR}, \cle{REGIFERCAM}), baisses successives des salaires (\cle{-20 \%} en 1993, puis \cle{-50 \%} sur les primes), gel des recrutements, dévaluation rampante avant la dévaluation officielle.

\textbf{La dévaluation du FCFA (janvier 1994).} Mesure la plus spectaculaire de la période, la dévaluation du franc CFA est décidée le 11 janvier 1994 à Dakar par les chefs d'État de la zone franc, sous l'égide de la France et du FMI.

\begin{exemplebox}[La dévaluation du FCFA de janvier 1994]
Le 12 janvier 1994, la parité fixe passe de \num{50} FCFA à \num{100} FCFA pour 1 franc français : la valeur de la monnaie est \cle{divisée par deux}.
\begin{itemize}
\item \textbf{Conséquences immédiates :} les importations (médicaments, hydrocarbures, produits manufacturés, blé) coûtent deux fois plus cher en monnaie locale ; le pouvoir d'achat des salariés et des retraités s'effondre ; l'inflation dépasse \num{30} \% en 1994.
\item \textbf{Effets recherchés :} les exportations camerounaises (cacao, café, bois, coton, pétrole) deviennent plus compétitives sur le marché mondial ; les recettes d'exportation en FCFA augmentent mécaniquement.
\end{itemize}
\end{exemplebox}

\begin{popfigure}
\begin{center}
\begin{tikzpicture}
\begin{axis}[
width=0.9\linewidth,height=6.5cm,
xlabel={Ann\'ee},ylabel={PIB par habitant (USD, valeurs indicatives)},
xmin=1979,xmax=2001,ymin=300,ymax=1050,
xtick={1980,1984,1988,1992,1996,2000},
grid=major,grid style={popline!40},
legend style={at={(0.03,0.97)},anchor=north west,font=\small}
]
\addplot[popcurve,mark=*,mark size=1.2pt] coordinates {
(1980,800) (1982,880) (1984,950) (1986,1000) (1988,900)
(1990,800) (1992,650) (1994,450) (1996,500) (1998,560) (2000,600)
};
\addlegendentry{PIB par habitant}
\draw[poporange,dashed,thick] (axis cs:1986,300) -- (axis cs:1986,1050);
\node[poporange,font=\small,anchor=west] at (axis cs:1986.2,980) {d\'ebut de la crise (1986)};
\draw[poppink,dashed,thick] (axis cs:1994,300) -- (axis cs:1994,1050);
\node[poppink,font=\small,anchor=west] at (axis cs:1994.2,380) {d\'evaluation du FCFA (1994)};
\end{axis}
\end{tikzpicture}
\end{center}
\legende{Évolution indicative du PIB par habitant au Cameroun (1980-2000). On observe un sommet vers 1986 (effet pétrole), une chute continue jusqu'à la dévaluation de 1994 (point bas), puis une reprise lente et incomplète. Source : données Banque mondiale, valeurs simplifiées.}
\end{popfigure}

\begin{methode}[Commenter une courbe économique en Histoire]
Pour commenter une courbe économique en Histoire :
\begin{enumerate}
\item \textbf{Présenter} la courbe : nature des données, période couverte, source, unité de mesure.
\item \textbf{Repérer les dates charnières} : sommets, points bas, ruptures de pente (ici 1986 et 1994).
\item \textbf{Décrire les tendances} : hausse, stagnation, chute, reprise — en citant des valeurs chiffrées approximatives.
\item \textbf{Expliquer historiquement} : relier chaque tendance aux événements politiques et économiques étudiés (chute des cours, ajustement structurel, dévaluation, reprise des exportations).
\item \textbf{Conclure} : dégager le sens général (ici, une « décennie perdue » 1986-1994 dont les effets sociaux nourrissent l'explosion démocratique).
\end{enumerate}
\end{methode}

\courssub{Les « années de braise » et le retour du multipartisme (1990-1992)}

\textbf{Le contexte international et interne.} La crise économique radicalise le mécontentement social (grèves étudiantes, mutineries de garnison). Parallèlement, la chute du mur de Berlin (novembre 1989) et le discours de La Baule de François Mitterrand (juin 1990), conditionnant l'aide française à la démocratisation, créent une fenêtre d'opportunité pour l'opposition.

\textbf{La création du SDF.} Le 26 mai 1990, à Bamenda (chef-lieu de la province du Nord-Ouest), John \cle{Fru Ndi}, libraire anglophone, lance le \cle{Social Democratic Front} (SDF) malgré l'interdiction préfectorale. La répression brutale du meeting inaugural fait plusieurs morts (bilan officiel : 6, opposition : une vingtaine). L'événement marque la naissance de la première grande formation d'opposition structurée, ancrée dans l'espace anglophone mais à vocation nationale.

\begin{definition}[Le multipartisme]
Le \cle{multipartisme} est un système politique qui autorise l'existence, l'activité légale et la compétition électorale de plusieurs partis politiques. Il s'oppose au \cle{parti unique}, où une seule formation est autorisée (cas du Cameroun de 1966 à 1990 avec l'UNC puis le RDPC).
\end{definition}

Sous la pression de la rue (manifestations, grèves, « \cle{villes mortes} » précoces) et de la communauté internationale, le président Biya promulgue en \cle{décembre 1990} un train de « lois sur les libertés » : la loi n°90/053 du 19 décembre 1990 rétablit officiellement le multipartisme intégral. En quelques mois, plus de cent partis politiques se créent (UPC, MDR, CDU, etc.).

\textbf{Les « années de braise » (1991-1992).} L'opposition, regroupée au sein de la \cle{Coordination des forces de l'opposition} (SDF, UPC, MDR, CDU...), exige la tenue d'une \cle{Conférence nationale souveraine} sur le modèle béninois (février 1990). Le pouvoir refuse, proposant un simple « Grand débat national » (mars-avril 1991) sans pouvoir constituant. L'opposition lance alors l'opération \cle{« villes mortes »} (novembre 1991 - mars 1992) : grèves générales illimitées, fermeture des commerces et marchés, paralysie des transports dans les grandes villes (Douala, Bamenda, Bafoussam, Yaoundé, Limbe). L'économie est asphyxiée. Le pouvoir répond par la \cle{répression} : arrestations massives de militants, interdiction de meetings, état d'urgence déclaré dans sept provinces sur dix en octobre 1992, assignation à résidence de John Fru Ndi. Le bilan humain est lourd : plusieurs dizaines de morts (ONG : une centaine).

\courssub{Des élections contestées : le contentieux électoral permanent}

\textbf{La présidentielle du 11 octobre 1992.} Première élection présidentielle pluraliste de l'histoire du Cameroun. Elle oppose Paul Biya (RDPC) à John Fru Ndi (SDF), Bello Bouba Maigari (UNDP) et d'autres. Les résultats officiels proclamés par la Cour suprême donnent : Paul Biya \num{39,98} \%, John Fru Ndi \num{35,97} \%, Bello Bouba Maigari \num{19,20} \%. L'écart est inférieur à 4 points. L'opposition dénonce des fraudes massives (bourrage d'urnes, vote multiple, disparitions de procès-verbaux dans le Nord-Ouest) et proclame la victoire de Fru Ndi. La tension aboutit à l'instauration de l'état d'urgence dans la province du Nord-Ouest, fief du SDF.

\textbf{Un contentieux électoral récurrent.} Toutes les élections suivantes donnent lieu à des contestations systématiques : fiabilité des listes électorales (omissions, doublons), bourrage d'urnes allégué, faiblesse structurelle des organismes de gestion (l'ONEL, puis \cle{ELECAM} créé en 2006, placé sous la tutelle du pouvoir exécutif), rôle controversé du Conseil constitutionnel (jugé partial) dans la proclamation des résultats définitifs.

\begin{popfigure}
\begin{center}
\begin{tikzpicture}
\begin{axis}[
width=0.95\linewidth,height=6.5cm,
ybar,bar width=8pt,
ylabel={Score officiel (\% des suffrages exprimés)},
symbolic x coords={1992,1997,2004,2011,2018},
xtick=data,ymin=0,ymax=100,
nodes near coords,nodes near coords style={font=\scriptsize},
grid=major,grid style={popline!40},
legend style={at={(0.5,1.05)},anchor=south,legend columns=2,font=\small},
enlarge x limits=0.15
]
\addplot[fill=popblue] coordinates {(1992,39.98) (1997,92.57) (2004,70.92) (2011,77.99) (2018,71.28)};
\addlegendentry{Paul Biya (RDPC)}
\addplot[fill=poporange] coordinates {(1992,35.97) (1997,0) (2004,17.40) (2011,10.71) (2018,14.23)};
\addlegendentry{Principal opposant (Fru Ndi 1992-2011, Kamto 2018)}
\end{axis}
\end{tikzpicture}
\end{center}
\legende{Scores officiels des élections présidentielles depuis le retour au multipartisme. En 1997, le SDF, l'UNDP et l'UPC boycottent le scrutin (d'où l'absence de barre orange). En 2018, Maurice Kamto (MRC) conteste officiellement les résultats devant le Conseil constitutionnel. Source : résultats officiels proclamés, valeurs arrondies à deux décimales.}
\end{popfigure}

\textbf{Les réélections de Paul Biya.} Réélu en 1997 (scrutin boycotté par les principaux opposants, taux de participation officiel \num{84} \%), 2004 (\num{70,9} \%), 2011 (\num{78,0} \%) et 2018 (\num{71,3} \%), Paul Biya cumule plus de quarante-deux ans au pouvoir (2024), ce qui constitue le plus long règne d'un chef d'État en exercice en Afrique centrale et l'un des plus longs du continent.

\courssub{La Constitution du 18 janvier 1996 et la décentralisation}

\textbf{Une nouvelle loi fondamentale.} Adoptée par référendum le 18 janvier 1996 (promulguée le 24 janvier), elle remplace la Constitution de 1972 révisée. Elle institue un \cle{État décentralisé} : « une République une et indivisible, laïque, démocratique et sociale » (art. 1). Les collectivités territoriales décentralisées sont la \cle{région} et la \cle{commune} (art. 55). Elle prévoit un Parlement bicaméral : l'Assemblée nationale (députés élus au suffrage direct) et un \cle{Sénat} (sénateurs élus au suffrage indirect par les conseillers municipaux et régionaux — 70 élus, 30 nommés par le président de la République). Le mandat présidentiel est fixé à 7 ans, renouvelable une seule fois (art. 6 al. 2).

\textbf{Une mise en œuvre très lente (17 ans pour le Sénat, 24 ans pour les régions).}
\begin{itemize}
\item Le Sénat n'est effectivement installé qu'en \cle{juin 2013} (premières élections sénatoriales le 14 avril 2013), soit dix-sept ans après la Constitution.
\item Les premières \cle{élections régionales} se tiennent le \cle{6 décembre 2020} : elles élisent les conseillers régionaux (90 par région) qui doivent gérer les compétences transférées (développement économique local, équipements, formation professionnelle, environnement). Le RDPC remporte neuf des dix régions ; le SDF l'emporte dans la région du Nord-Ouest.
\end{itemize}

\begin{attention}[Ne pas confondre décentralisation et fédéralisme]
La \cle{décentralisation} est un transfert de compétences et de moyens de l'État central vers des collectivités locales (communes, régions) qui restent sous la tutelle administrative de l'État (le préfet représente l'État dans le département, le gouverneur dans la région) : le Cameroun reste un État \emph{unitaire}. Le \cle{fédéralisme}, en revanche, associe des États fédérés dotés de leur propre constitution, de leur gouvernement, de leur parlement et de larges autonomies législatives et fiscales (cas du Cameroun fédéral de 1961 à 1972). La Constitution de 1996 choisit la décentralisation, non le retour au fédéralisme.
\end{attention}

\courssub{La révision constitutionnelle de 2008 : le verrouillage du pouvoir}

En avril 2008, dans un contexte social explosif (émeutes de février 2008 contre la vie chère et la révision constitutionnelle annoncée), l'Assemblée nationale (dominée à plus de 90 \% par le RDPC) adopte la révision de l'article 6 alinéa 2 de la Constitution : la \cle{limitation à deux} du nombre de mandats présidentiels (7 ans renouvelable une fois) est \cle{supprimée}. Paul Biya peut désormais se représenter indéfiniment. Cette révision (loi n°2008/001 du 14 avril 2008), votée à la hâte sans consensus national, est vivement critiquée par l'opposition et la société civile, qui y voient un « coup d'État constitutionnel » verrouillant l'alternance. Elle ouvre la voie aux candidatures de 2011 et 2018.

\courssub{Les défis du Cameroun contemporain (années 2010-2020)}

Au-delà des mutations institutionnelles, le Cameroun fait face à des défis majeurs qui structurent l'actualité politique et sociale.

\begin{itemize}
\item \textbf{La crise anglophone (depuis 2016).} Née de revendications corporatistes (avocats, enseignants du Common Law) sur la marginalisation du système juridique et éducatif anglophone, elle bascule dans l'insurrection armée fin 2017 (création de groupes séparatistes « Amba » proclamant l'indépendance de l'« Ambazonie »). Bilan (estimations ONG, 2024) : plus de 6 000 morts, plus de 700 000 déplacés internes, centaines de milliers de réfugiés au Nigeria, destruction massive d'infrastructures scolaires et sanitaires dans le Nord-Ouest et le Sud-Ouest. Le « Grand Dialogue National » (sept.-oct. 2019) aboutit à un statut spécial pour les deux régions, jugé insuffisant par les séparatistes.
\item \textbf{L'insécurité dans l'Extrême-Nord (Boko Haram / ISWAP).} Depuis 2014, la secte terroriste nigériane Boko Haram (puis sa scission ISWAP) mène des attaques meurtrières, des enlèvements et du pillage dans la région de l'Extrême-Nord (frontière nigériane et lac Tchad). Bilan : plusieurs milliers de morts, plus de 300 000 déplacés internes, effondrement de l'économie locale (agriculture, commerce transfrontalier). L'armée camerounaise, appuyée par la Force multinationale mixte (FMM), riposte mais la menace persiste.
\item \textbf{Gouvernance, corruption et dette.} Malgré l'opération « \cle{Épervier} » (lancée 2006, coups de filet médiatisés contre des hauts fonctionnaires), la corruption reste endémique (indices Transparency International). La dette publique dépasse \num{60} \% du PIB (2023), limitant les marges de manœuvre budgétaires. La dépendance aux matières premières (pétrole, bois, cacao, aluminium) persiste malgré les discours sur la transformation structurelle (Stratégie nationale de développement 2020-2030).
\item \textbf{Jeunesse, emploi et urbanisation.} Avec plus de 60 \% de la population âgée de moins de 25 ans, l'insertion professionnelle des diplômés (chômage des jeunes \textgreater \num{13} \% officiel, bien plus en réel) et l'informel (\textgreater \num{90} \% de l'emploi) sont des bombes à retardement sociales.
\end{itemize}

\begin{exoresolu}[Commenter la courbe du PIB par habitant (1980-2000)]
Énoncé : à partir de la courbe du PIB par habitant au Cameroun (1980-2000), décris les grandes phases et propose une explication historique pour chacune.
\tcblower
\textbf{Solution.}
\textbf{Présentation} : la courbe représente l'évolution du PIB par habitant en dollars US constants (valeurs indicatives Banque mondiale) de 1980 à 2000.
\textbf{Description} : on distingue trois phases majeures.
\begin{enumerate}
\item \textbf{1980-1986 : Croissance soutenue.} Le PIB par habitant passe d'environ \num{800} à \num{1000} USD (+25 \%). Sommet atteint en 1986.
\item \textbf{1986-1994 : Récession profonde (« décennie perdue »)}. Chute brutale et continue jusqu'à \num{450} USD en 1994 (-55 \% par rapport au pic de 1986). Point bas coïncidant avec la dévaluation.
\item \textbf{1994-2000 : Reprise lente et fragile.} Remontée progressive jusqu'à \num{600} USD en 2000, sans retrouver le niveau de 1986.
\end{enumerate}
\textbf{Explication historique} :
\begin{itemize}
\item La croissance 1980-1986 repose sur la rente pétrolière (début production 1977, pic 1985) et des cours élevés du café/cacao.
\item La récession 1986-1994 s'explique par l'effondrement simultané des cours du pétrole (\num{-50} \% en 1986), du café et du cacao, aggravée par la rigidité des dépenses publiques, puis par les plans d'ajustement structurel (1988 ss.) : compression salariale, licenciements, baisse de l'investissement public.
\item Le point bas de 1994 est l'effet mécanique de la dévaluation du FCFA (division par deux de la valeur en USD).
\item La reprise post-1994 tient au gain de compétitivité des exportations (cacao, bois, coton) et à la reprise des cours du pétrole, mais reste fragile (croissance extensive, pas structurelle).
\end{itemize}
\textbf{Conclusion} : la courbe illustre une « décennie perdue » (1986-1994) dont les effets sociaux (pauvreté, chômage, inégalités) constituent le terreau des revendications démocratiques des « années de braise » (1991-1992).
\end{exoresolu}

\begin{aretenir}[L'essentiel : Le Cameroun de 1982 à nos jours]
\begin{itemize}
\item \textbf{1982-1986 :} Succession Ahidjo $\to$ Biya ; « rigueur », « moralisation » ; coup d'État avorté avril 1984.
\item \textbf{1986-1994 :} Crise économique (chute cours matières premières) ; plans d'ajustement structurel (FMI/BM) ; dévaluation FCFA (janv. 1994, divisée par 2).
\item \textbf{1990-1992 :} Création SDF (26 mai 1990) ; lois libertés déc. 1990 (multipartisme) ; « années de braise » (villes mortes, répression) ; présidentielle contestée oct. 1992 (Biya 39,98 \% vs Fru Ndi 35,97 \%).
\item \textbf{1996 :} Constitution : État décentralisé (régions, communes), Sénat (installé 2013), mandat 7 ans renouvelable 1 fois.
\item \textbf{2008 :} Révision constitutionnelle : suppression limitation mandats (verrouillage pouvoir).
\item \textbf{Élections :} Biya réélu 1997 (boycott), 2004, 2011, 2018 ; élections régionales déc. 2020.
\item \textbf{Défis actuels :} Crise anglophone (depuis 2016, conflit armé NO/SO), Boko Haram (Extrême-Nord), corruption/dette, chômage des jeunes.
\end{itemize}
\end{aretenir}

\begin{exercice}[Pour s'entraîner — Type Bac / Probatoire]
1. Définis le multipartisme et explique en quoi la loi du 19 décembre 1990 marque une rupture majeure dans l'histoire politique du Cameroun.
2. À l'aide du diagramme en barres (scrutins présidentiels 1992-2018), compare le scrutin de 1992 à celui de 1997 et explique la différence des scores obtenus par Paul Biya.
3. Rédige un paragraphe argumenté (15-20 lignes) répondant à la question : \emph{La Constitution du 18 janvier 1996 a-t-elle réellement démocratisé la vie politique camerounaise ?} (Tu t'appuieras sur les avancées institutionnelles et sur les limites persistantes).
\end{exercice}

\begin{corrige}[Exercice 1]
\textbf{1.} Le multipartisme est un système politique qui autorise l'existence, l'activité légale et la compétition électorale de plusieurs partis politiques. La loi n°90/053 du 19 décembre 1990 met fin à vingt-quatre ans de parti unique (UNC de 1966 à 1985, puis RDPC) et légalise l'opposition politique. C'est une rupture majeure obtenue sous la double pression de la mobilisation populaire (« villes mortes », création du SDF) et du contexte international (fin de la Guerre froide, discours de La Baule).

\textbf{2.} En 1992, le scrutin est serré : Paul Biya obtient \num{39,98} \% contre \num{35,97} \% à John Fru Ndi (SDF), dans un contexte de forte mobilisation de l'opposition et de participation massive. En 1997, les principaux partis d'opposition (SDF, UNDP, UPC) boycottent l'élection pour protester contre l'absence de commission électorale indépendante et le non-respect des engagements du « Grand Dialogue » ; Paul Biya obtient ainsi un score pléthorique de \num{92,57} \% face à des candidats de second plan. La différence s'explique donc par la participation (1992) vs le boycott (1997) de l'opposition majeure.

\textbf{3.} \emph{Thèse attendue (nuancée) :} La Constitution de 1996 pose des jalons démocratiques réels : consécration du multipartisme (art. 3), création du Sénat et des régions (décentralisation), limitation du mandat présidentiel (art. 6 al. 2 initial), indépendance théorique de la justice. \emph{Mais} sa mise en œuvre est tronquée : Sénat installé 17 ans plus tard (2013), régions élues 24 ans plus tard (2020), révision de 2008 supprimant la limitation des mandats (verrouillage de l'alternance), contentieux électoral récurrent (ELECAM contesté, Conseil constitutionnel perçu comme aux ordres), répression des libertés (crise anglophone, arrestations opposants). \emph{Conclusion :} Cadre formel démocratique, mais pratique politique encore largement autoritaire ; la démocratisation reste inachevée.
\end{corrige}

\coursec{Les défis du Cameroun contemporain}

Depuis l'avènement de Paul Biya en 1982, le Cameroun a connu de profondes mutations politiques : retour au multipartisme en 1990, révision constitutionnelle de 1996, régionalisation inscrite dans la Constitution. Mais au-delà de ces évolutions institutionnelles, le pays fait face, au début du XXI\textsuperscript{e} siècle, à des défis considérables qui engagent son avenir : la sécurité des populations, la cohésion nationale, le développement économique, la qualité de la gouvernance et la place du pays dans le monde. Cette dernière partie de la leçon les examine un à un, avec la rigueur et la distance qui s'imposent à l'historien, même lorsqu'il traite de l'actualité la plus récente.

\begin{attention}[Analyser l'actualité sans jugement partisan]
Lorsqu'on étudie la période contemporaine, on est tenté de prendre parti, de louer ou de condamner. Ce n'est pas le rôle de l'historien, ni celui du candidat au Baccalauréat. Il faut \cle{analyser les faits} avec distance : dater les événements, identifier les acteurs, croiser les sources (discours officiels, rapports d'organisations internationales, presse, témoignages), présenter les causes et les conséquences sans slogan. Une copie qui ressemble à un tract politique est toujours sanctionnée ; une copie qui explique et nuance est toujours valorisée.
\end{attention}

\courssub{Le défi sécuritaire : Boko Haram dans l'Extrême-Nord}

Le premier défi est sécuritaire. La secte islamiste \cle{Boko Haram}, née au Nigeria dans les années 2000, étend ses exactions au bassin du lac Tchad. À partir de \cle{2013-2014}, elle multiplie les incursions dans la région de l'\cle{Extrême-Nord} du Cameroun : attaques de villages, enlèvements (dont celui de la famille Moulin-Fournier en 2013 et de dix religieuses en 2014), attentats-suicides dans les marchés et les mosquées de Fotokol, Maroua ou Mora.

La réponse camerounaise s'organise : déploiement de la force multinationale mixte (Cameroun, Nigeria, Niger, Tchad, Bénin) à partir de 2015, opérations militaires dans le massif de Mandara et autour du lac Tchad. Mais les conséquences humaines sont lourdes : des centaines de milliers de \cle{déplacés internes} et de réfugiés nigérians affluent, les écoles et les marchés ferment, les échanges transfrontaliers s'effondrent dans une région déjà pauvre. La sécurisation de l'Extrême-Nord mobilise durablement des ressources qui manquent ailleurs.

\courssub{La crise anglophone : origines et radicalisation (depuis 2016)}

Le second défi, plus inattendu, touche les deux régions issues du Cameroun britannique : le \cle{Nord-Ouest} (capitale Bamenda) et le \cle{Sud-Ouest} (capitale Buéa).

\begin{definition}[La crise anglophone]
La \cle{crise anglophone} est un conflit né des revendications des populations des régions du \cle{Nord-Ouest} et du \cle{Sud-Ouest}, qui dénoncent leur \cle{marginalisation} au sein de l'État unitaire camerounais, largement francophone dans son fonctionnement : faible place de l'anglais dans l'administration, nomination de magistrats et d'enseignants francophones dans les juridictions et les écoles anglophones, sentiment de relégation économique et politique.
\end{definition}

La chronologie est essentielle à retenir :
\begin{enumerate}
\item \cle{Octobre-novembre 2016} : les \cle{avocats} de common law, puis les \cle{enseignants} anglophones, entament des grèves pour protester contre la « francophonisation » des tribunaux et des universités. Leurs revendications sont d'abord corporatistes et réformistes.
\item \cle{2017} : la répression des manifestations, les arrestations de leaders et la coupure d'Internet dans les deux régions radicalisent le mouvement. Des groupes séparatistes proclament symboliquement, le 1\textsuperscript{er} octobre 2017, l'indépendance d'une « Ambazonie » imaginaire.
\item \cle{Depuis 2018} : affrontements armés entre groupes séparatistes et forces de défense, opérations de « villes mortes », boycotts scolaires, exactions contre les civils. Le bilan humain est lourd : des milliers de morts, des centaines de milliers de déplacés internes et de réfugiés au Nigeria.
\end{enumerate}

\begin{popfigure}
\begin{center}
\begin{tikzpicture}[scale=1.05,>=Stealth]
% Contour très simplifié du Cameroun
\draw[thick,popdark,fill=popcream] (2.2,6.2) -- (3.4,6.6) -- (4.6,5.6) -- (4.4,4.6) -- (3.6,3.9) -- (4.3,3.2) -- (4.1,1.9) -- (3.0,0.6) -- (2.0,0.2) -- (1.0,0.6) -- (0.4,1.8) -- (0.2,3.0) -- (0.8,4.0) -- (0.6,5.0) -- cycle;
% Foyers de crise
\fill[poporange!45] (3.6,5.4) circle (0.85);
\node[popdark,font=\bfseries\small] at (3.6,5.4) {Extrême-Nord};
\fill[poppink!55] (0.75,4.4) circle (0.55);
\fill[poppink!55] (0.65,3.3) circle (0.55);
\node[popdark,font=\bfseries\small] at (1.7,4.4) {Nord-Ouest};
\node[popdark,font=\bfseries\small] at (1.7,3.3) {Sud-Ouest};
% Villes
\fill[popblue] (2.3,2.3) circle (0.07); \node[right,font=\small,popblue] at (2.4,2.3) {Yaoundé};
\fill[popblue] (0.9,1.3) circle (0.07); \node[right,font=\small,popblue] at (1.0,1.15) {Douala};
\fill[popblue] (3.5,5.9) circle (0.07); \node[above,font=\small,popblue] at (3.5,6.0) {Maroua};
% Grands projets
\draw[popgreen,very thick] (0.55,0.85) circle (0.16);
\node[popgreen,font=\bfseries\small,anchor=west] at (0.8,0.55) {Port de Kribi (2018)};
\draw[popgreen,very thick] (3.35,3.55) circle (0.16);
\node[popgreen,font=\bfseries\small,anchor=west] at (3.55,3.5) {Barrage de Lom Pangar (2017)};
\draw[popgreen,very thick] (2.6,1.0) circle (0.16);
\node[popgreen,font=\bfseries\small,anchor=west] at (2.85,0.95) {Barrage de Memve'ele (2019)};
% Nord
\draw[->,thick,popdark] (0.35,5.6) -- (0.35,6.3) node[above,font=\small]{N};
\end{tikzpicture}
\end{center}
\legende{Croquis du Cameroun contemporain : en rose, les régions du Nord-Ouest et du Sud-Ouest touchées par la crise anglophone ; en orange, l'Extrême-Nord exposé à Boko Haram ; en vert, les grands projets structurants (port de Kribi, barrages de Lom Pangar et de Memve'ele). Carte schématique, sans précision cartographique.}
\end{popfigure}

\courssub{Le Grand Dialogue National et le statut spécial (2019)}

Face à l'impasse, le chef de l'État convoque un \cle{Grand Dialogue National}, qui se tient à Yaoundé du \cle{30 septembre au 4 octobre 2019}, sous la présidence du Premier ministre Joseph Dion Ngute. Il réunit délégués des dix régions, partis politiques, société civile et diaspora, mais une partie des séparatistes le boycotte.

Sa principale traduction concrète est la loi du \cle{24 décembre 2019} portant \cle{statut spécial} des régions du Nord-Ouest et du Sud-Ouest : elle reconnaît leur spécificité historique, linguistique et éducative (systèmes scolaire anglo-saxon et judiciaire de common law) et prévoit des organes régionaux renforcés, dans le cadre de la décentralisation. Ce statut constitue une avancée institutionnelle, mais sa mise en œuvre reste lente et les violences n'ont pas totalement cessé : la crise illustre la difficulté de résoudre un conflit identitaire par les seuls instruments juridiques.

\courssub{Le défi économique : l'émergence à l'horizon 2035}

\begin{definition}[L'émergence]
L'\cle{émergence} désigne, dans le vocabulaire du développement, l'objectif de faire passer un pays du statut de pays en développement à celui de \cle{pays à revenu intermédiaire}, industrialisé et intégré à l'économie mondiale. Le Cameroun s'est fixé cet objectif dans le document stratégique « \cle{Cameroun émergent à l'horizon 2035} » (2009), décliné en stratégie de croissance et d'emploi.
\end{definition}

Cette vision repose sur de \cle{grands projets structurants} :
\begin{itemize}
\item le \cle{port en eau profonde de Kribi} (Sud), inauguré en 2018, destiné à désengorger Douala et à desservir le Tchad et la Centrafrique ;
\item le \cle{barrage de Lom Pangar} (Est), achevé en 2017, qui régule le fleuve Sanaga pour sécuriser la production hydroélectrique ;
\item le \cle{barrage de Memve'ele} (Sud), sur la Ntem, inauguré en 2019, et le chantier de Nachtigal sur la Sanaga, pour résoudre le déficit énergétique qui freine l'industrie.
\end{itemize}

\begin{savaistu}[Le port de Kribi, géant d'Afrique centrale]
Inauguré en \cle{2018}, le port en eau profonde de \cle{Kribi} est l'un des plus grands projets portuaires d'Afrique centrale. Sa profondeur (environ 16 mètres) permet d'accueillir de très gros porte-conteneurs, impossibles à recevoir à Douala. Il est adossé à une zone industrielle et portuaire censée attirer des usines de transformation locale du bois, du cacao et du coton.
\end{savaistu}

Mais le chemin vers l'émergence est semé d'obstacles : croissance insuffisante pour réduire massivement la pauvreté, dette publique croissante, économie encore dépendante des matières premières (pétrole, cacao, bois) et du secteur informel.

\courssub{Le défi de la gouvernance et de la cohésion nationale}

Le quatrième défi est celui de la \cle{gouvernance}. La \cle{corruption}, régulièrement dénoncée par les organisations internationales et combattue par l'opération « Épervier » lancée en 2006, détourne des ressources publiques et décourage les investisseurs. Le \cle{chômage et le sous-emploi des jeunes} — plus des trois quarts des actifs travaillent dans l'informel — alimentent l'\cle{exode des compétences} : médecins, ingénieurs et enseignants partent chercher en Europe ou en Amérique du Nord les conditions de travail que le pays ne leur offre pas.

Le cinquième défi est celui de la \cle{cohésion nationale}. Le Cameroun compte environ \cle{250 groupes ethniques}, deux langues officielles (français et anglais) et plusieurs grandes religions. Ce pluralisme est une richesse, mais aussi une fragilité lorsque les appartenances communautaires sont instrumentalisées à des fins politiques. Le \cle{bilinguisme}, inscrit dans la Constitution, et la promotion du \cle{vivre-ensemble} (Commission nationale pour la promotion du bilinguisme et du multiculturalisme, créée en 2017) sont présentés comme les remparts de l'unité nationale.

\begin{popfigure}
\begin{center}
\begin{tikzpicture}[scale=1.0,>=Stealth]
\node[circle,draw=popdark,fill=popblue!25,thick,minimum size=2.6cm,align=center,font=\bfseries] (c) at (0,0) {Cameroun\\contemporain};
\node[draw=poporange,fill=poporange!25,rounded corners,align=center,font=\small] (d1) at (90:3.1) {\textbf{Sécurité}\\Boko Haram,\\crise anglophone};
\node[draw=popgreen,fill=popgreen!25,rounded corners,align=center,font=\small] (d2) at (18:3.3) {\textbf{Économie}\\Émergence 2035,\\grands projets};
\node[draw=poppurple,fill=poppurple!25,rounded corners,align=center,font=\small] (d3) at (-54:3.3) {\textbf{Gouvernance}\\corruption, chômage,\\exode des compétences};
\node[draw=poppink,fill=poppink!30,rounded corners,align=center,font=\small] (d4) at (-126:3.3) {\textbf{Cohésion nationale}\\bilinguisme, vivre-ensemble,\\250 groupes ethniques};
\node[draw=popblue,fill=popblue!20,rounded corners,align=center,font=\small] (d5) at (162:3.3) {\textbf{Rayonnement}\\CEMAC, CEEAC, UA,\\France, Chine};
\draw[thick,popmuted] (c) -- (d1); \draw[thick,popmuted] (c) -- (d2);
\draw[thick,popmuted] (c) -- (d3); \draw[thick,popmuted] (c) -- (d4);
\draw[thick,popmuted] (c) -- (d5);
\end{tikzpicture}
\end{center}
\legende{Schéma en étoile des cinq grands défis du Cameroun contemporain : sécurité, économie, gouvernance, cohésion nationale et rayonnement international.}
\end{popfigure}

\courssub{Le Cameroun sur la scène internationale et la question de la succession}

Malgré ces difficultés internes, le Cameroun demeure un acteur diplomatique de poids en Afrique centrale. Il est membre de la \cle{CEMAC} (Communauté économique et monétaire de l'Afrique centrale), dont il est la première économie, de la \cle{CEEAC} et de l'\cle{Union africaine}. Sa coopération est diversifiée : la \cle{France} reste un partenaire historique (coopération militaire contre Boko Haram, investissements), tandis que la \cle{Chine} s'impose comme premier bailleur d'infrastructures (port de Kribi, routes, stades, immeubles publics). Cette diversification illustre la nouvelle donne des relations internationales africaines.

Enfin, une question traverse toute la vie politique : celle de la \cle{succession présidentielle}. Paul Biya, né en 1933, au pouvoir depuis 1982 et réélu en 2018 pour un septième mandat, est l'un des plus vieux chefs d'État du monde. La Constitution prévoit qu'en cas de vacance, le président du Sénat assure l'intérim et organise de nouvelles élections. Mais l'application effective de ces dispositions, jamais testée, suscite des débats sur l'\cle{avenir institutionnel} du pays : l'expérience de 1982 (passation pacifique entre Ahidjo et Biya) servira-t-elle de précédent, ou la transition sera-t-elle conflictuelle ? C'est une question ouverte, que l'historien se contente de poser sans y répondre par des pronostics.

\begin{exoresolu}[Sujet type Bac : dissertation]
\textbf{Sujet :} « Le Cameroun contemporain est-il un État en crise ou un État en construction ? » Rédiger l'introduction et la conclusion de cette dissertation.
\tcblower
\textbf{Introduction (corrigé).} Depuis son indépendance en 1960, le Cameroun s'est doté d'un État unitaire et d'institutions stables. Pourtant, depuis 2014, il fait face à la menace de Boko Haram dans l'Extrême-Nord et, depuis 2016, à une crise anglophone qui a dégénéré en conflit armé dans le Nord-Ouest et le Sud-Ouest. On peut alors se demander si ces crises révèlent la faiblesse de l'État camerounais ou si elles s'accompagnent d'une construction continue de la nation et de l'économie. Nous verrons d'abord la réalité des crises (sécuritaires, politiques, sociales), puis les signes d'une construction en marche (grands projets, décentralisation, intégration régionale), avant de nuancer : construction et fragilités coexistent.

\textbf{Conclusion (corrigé).} Le Cameroun contemporain est traversé par de graves crises — Boko Haram, conflit anglophone, corruption — qui révèlent les fragilités de l'État. Mais il poursuit en même temps sa construction : statut spécial de 2019, grands projets d'infrastructures, ambition d'émergence en 2035. Il est donc moins un État en faillite qu'un État inachevé, dont l'avenir dépendra de sa capacité à transformer la paix et la bonne gouvernance en développement partagé. Cette question rejoint celle, plus large, de la consolidation des États africains nés des indépendances.
\end{exoresolu}

\begin{methode}[Construire une conclusion de dissertation]
Au Baccalauréat, la conclusion se construit en trois temps :
\begin{enumerate}
\item \textbf{Rappeler le problème} posé dans l'introduction, en une phrase reformulée (ne jamais recopier l'introduction).
\item \textbf{Dresser un bilan nuancé} : répondre clairement à la question en synthétisant les deux parties du développement, avec un mot de mesure (« cependant », « toutefois », « il reste que »).
\item \textbf{Ouvrir sur un défi d'avenir} : élargir vers une question nouvelle mais liée (par exemple : la succession présidentielle, l'intégration régionale africaine), sans développer.
\end{enumerate}
\end{methode}

\begin{exercice}[Pour s'entraîner]
\begin{enumerate}
\item Définis la crise anglophone et rappelle ses deux déclencheurs de 2016.
\item Cite deux grands projets liés à la vision « émergence 2035 » et explique leur utilité économique.
\item Rédige en dix lignes un paragraphe argumenté et non partisan sur le Grand Dialogue National de 2019 : contexte, décision principale, limites.
\end{enumerate}
\end{exercice}

\begin{corrige}[Exercice]
\begin{enumerate}
\item La crise anglophone est le conflit né des revendications des populations du Nord-Ouest et du Sud-Ouest, qui dénoncent leur marginalisation dans l'État unitaire. Elle a été déclenchée en 2016 par la grève des avocats de common law (contre la nomination de magistrats francophones) puis par celle des enseignants anglophones (contre la « francophonisation » des écoles et universités).
\item Le port en eau profonde de Kribi (2018) permet d'accueillir de grands navires et de développer le commerce et l'industrie ; le barrage de Lom Pangar (2017) régule la Sanaga et sécurise la production d'électricité, condition de l'industrialisation. (Memve'ele est aussi acceptable.)
\item Exemple de réponse attendue : convoqué fin septembre 2019 pour répondre à la crise anglophone, le Grand Dialogue National a abouti à la loi de décembre 2019 accordant un statut spécial au Nord-Ouest et au Sud-Ouest, reconnaissant leurs spécificités linguistiques, scolaires et judiciaires. C'est une avancée institutionnelle réelle, mais limitée par le boycott d'une partie des séparatistes et par la lenteur de sa mise en œuvre, les violences n'ayant pas cessé.
\end{enumerate}
\end{corrige}

\begin{aretenir}[L'essentiel]
\begin{itemize}
\item Depuis \cle{2013-2014}, Boko Haram menace l'Extrême-Nord (incursions, attentats, déplacés internes).
\item Depuis \cle{2016}, la crise anglophone (grèves des avocats et enseignants, puis radicalisation séparatiste) déstabilise le Nord-Ouest et le Sud-Ouest ; le \cle{Grand Dialogue National} (2019) débouche sur un \cle{statut spécial} pour ces deux régions.
\item La vision « \cle{émergence 2035} » s'appuie sur de grands projets : port de Kribi, barrages de Lom Pangar et Memve'ele.
\item Gouvernance (corruption, chômage des jeunes, exode des compétences) et cohésion nationale (bilinguisme, vivre-ensemble des 250 groupes ethniques) restent des chantiers majeurs.
\item Membre actif de la CEMAC, de la CEEAC et de l'Union africaine, le Cameroun coopère avec la France et la Chine ; la question de la succession présidentielle demeure ouverte.
\end{itemize}
\end{aretenir}

\newpage

\coursec{Activité d'intégration}

\coursec{Le Cameroun indépendant}

\courssub{Objectifs de l'activité}

À l'issue de ce travail pratique, réalisé en groupes de quatre élèves, tu seras capable de :
\begin{itemize}
\item construire un \cle{croquis cartographique} ordonné (titre, légende numérotée, orientation, échelle indicative) retraçant la formation territoriale du Cameroun indépendant, des tutelles de 1946 à l'État unitaire de 1972 ;
\item exploiter un \cle{tableau statistique} pour tracer un graphique d'évolution et y identifier des phases (croissance, crise, reprise) ;
\item analyser et croiser des \cle{documents historiques} (discours, résultats de scrutins) en relevant les faits, les dates et les acteurs ;
\item rédiger un \cle{texte argumenté} d'une quinzaine de lignes, répondant à une question de bilan et justifié par au moins quatre faits datés ;
\item présenter oralement une production collective et participer à la correction au tableau.
\end{itemize}

\begin{aretenir}[Savoir-faire évalués]
Cette activité mobilise les deux savoir-faire officiels de l'unité : appliquer les notions de l'unité à des situations concrètes, et rédiger puis justifier une démarche, une réponse ou une conclusion.
\end{aretenir}

\courssub{Situation}

Dans le cadre de la semaine de la citoyenneté organisée par ton lycée technique de Douala, le club d'histoire doit réaliser une exposition intitulée « \emph{Le Cameroun indépendant : une réussite ?} ». Le proviseur a confié à ta classe la production des panneaux de l'exposition. Ton groupe de quatre élèves reçoit le dossier documentaire suivant.

\textbf{Document 1 — Carte muette du Cameroun} (à compléter) : contour du pays, frontières, position de la ligne de partage de 1919 entre zone française (à l'est) et zone britannique (à l'ouest, le long du Nigeria), villes repères : Yaoundé, Douala, Buea, Bamenda, Garoua, Foumban.

\textbf{Document 2 — Évolution du PIB par habitant au Cameroun} (valeurs indicatives, en dollars courants) :

\begin{center}
\begin{tabularx}{\linewidth}{|l|X|X|X|X|X|X|X|X|X|X|}
\hline
\rowcolor{popblueL} Année & 1960 & 1970 & 1980 & 1986 & 1990 & 1994 & 2000 & 2010 & 2020 \\
\hline PIB/habitant (USD) & 70 & 160 & 890 & 1000 & 960 & 630 & 700 & 1250 & 1500 \\
\hline
\end{tabularx}
\end{center}

\textbf{Document 3 — Extraits de discours.} \emph{« Nous voulons bâtir un Cameroun unifié, où fils du Nord et du Sud, de l'Est et de l'Ouest, se reconnaîtront frères dans une même nation »} (Ahmadou Ahidjo, conférence de Foumban, juillet 1961). \emph{« Je m'engage à poursuivre l'œuvre de construction nationale dans la rigueur et le moralisme, au service de tous les Camerounais »} (Paul Biya, discours d'investiture, novembre 1982).

\textbf{Document 4 — Résultats de scrutins.} Plébiscite du 11 février 1961 : le Cameroun septentrional britannique choisit le Nigeria (environ 60 \%), le Cameroun méridional britannique choisit le Cameroun (environ 70 \%). Élection présidentielle du 11 octobre 1992 (première élection pluraliste) : Paul Biya environ 40 \%, John Fru Ndi environ 36 \%.

\textbf{Travail attendu :} trois productions — un croquis, un graphique, un texte argumenté — puis une restitution orale de cinq minutes.

\courssub{Consignes}

\begin{enumerate}
\item Sur la carte muette, trace la ligne de partage de 1919, colorie différemment la zone française et la zone britannique, puis indique par des flèches et des dates les étapes : indépendance du Cameroun français (1\ier{} janvier 1960), plébiscites du 11 février 1961, République fédérale (1\ier{} octobre 1961), République unie (1972). Rédige une légende numérotée et ordonnée.
\item À partir du document 2, trace le graphique de l'évolution du PIB par habitant de 1960 à 2020. Calcule le taux de variation entre 1960 et 1980, puis entre 1986 et 1994. Identifie trois phases et propose pour chacune une explication historique.
\item Dans le document 3, relève pour chaque discours : l'auteur, la date, le contexte, l'idée principale. Quel objectif commun assignent les deux présidents à l'État camerounais ?
\item À l'aide du document 4, explique ce que révèle le plébiscite de 1961 sur la construction territoriale du pays, et ce que le score de 1992 révèle de l'ouverture démocratique.
\item Rédige en 15 lignes environ une réponse argumentée à la question : « Le Cameroun indépendant est-il une réussite ? ». Ta réponse doit comporter une position claire, au moins quatre faits datés tirés du dossier, et une nuance (limites ou défis).
\item Prépare la restitution orale : répartissez la parole entre les quatre membres (présentation du croquis, du graphique, du texte, réponse aux questions).
\end{enumerate}

\courssub{Ressources à mobiliser}

\begin{itemize}
\item \cle{Chronologie de l'unité} : tutelles de 1946 ; indépendance du Cameroun français le 1\ier{} janvier 1960 ; plébiscites du 11 février 1961 ; conférence de Foumban (juillet 1961) ; République fédérale du Cameroun (1\ier{} octobre 1961) ; référendum du 20 mai 1972 instituant la République unie ; démission d'Ahidjo et investiture de Paul Biya (4-6 novembre 1982) ; retour au multipartisme (1990) ; élection pluraliste de 1992.
\item \cle{Bilan de l'ère Ahidjo} : boom pétrolier à partir de 1977-1978, « années de prospérité » ; crise économique à partir de 1986-1987 ; ajustement structurel et dévaluation du franc CFA (janvier 1994).
\item \cle{Technique du croquis} : titre, flèche du nord, échelle indicative, légende numérotée classée de l'amont vers l'aval (ou du passé vers le présent).
\item \cle{Calcul de taux de variation} : $t = \dfrac{V_A - V_D}{V_D} \times 100$, où $V_D$ est la valeur de départ et $V_A$ la valeur d'arrivée.
\item \cle{Méthode du texte argumenté} : thèse annoncée, arguments illustrés de faits datés, nuance, conclusion.
\end{itemize}

\courssub{Démarche et corrigé commenté}

\begin{exoresolu}[Corrigé commenté du TP]
\textbf{Étape 1 — Le croquis de la formation territoriale.} On trace d'abord le contour simplifié du Cameroun, la frontière avec le Nigeria à l'ouest et la ligne de partage de 1919. On colorie en deux teintes la zone française (Est) et la zone britannique (Ouest). On place ensuite les flèches et les dates dans l'ordre chronologique, et l'on rédige la légende numérotée.

\begin{popfigure}
\begin{tikzpicture}[scale=0.9,font=\small]
% contour simplifié du Cameroun
\draw[thick,popdark,fill=popcream] (2,6.5) -- (3.2,7.2) -- (3.4,8.6) -- (2.6,9.6) -- (2.8,10.6) -- (4.2,10.4) -- (5.4,9.4) -- (6.4,8.8) -- (7.2,7.4) -- (7.6,5.6) -- (6.8,4.2) -- (6.4,2.6) -- (5.2,1.2) -- (3.4,0.6) -- (1.8,1.4) -- (1.2,3.2) -- (0.8,4.8) -- (1.4,6.0) -- cycle;
% ligne de partage 1919
\draw[very thick,popink,dashed] (1.6,1.2) -- (1.9,3.0) -- (2.1,5.0) -- (2.4,6.6) -- (2.9,7.4);
% zone britannique
\fill[poporange!40] (1.6,1.2) -- (1.9,3.0) -- (2.1,5.0) -- (2.4,6.6) -- (2.9,7.4) -- (2,6.5) -- (1.4,6.0) -- (0.8,4.8) -- (1.2,3.2) -- (1.8,1.4) -- cycle;
\draw[thick,popdark] (2,6.5) -- (3.2,7.2) -- (3.4,8.6) -- (2.6,9.6) -- (2.8,10.6) -- (4.2,10.4) -- (5.4,9.4) -- (6.4,8.8) -- (7.2,7.4) -- (7.6,5.6) -- (6.8,4.2) -- (6.4,2.6) -- (5.2,1.2) -- (3.4,0.6) -- (1.8,1.4) -- (1.2,3.2) -- (0.8,4.8) -- (1.4,6.0) -- cycle;
% villes
\fill[popblue] (4.6,3.2) circle (0.08) node[right,popdark]{Yaoundé};
\fill[popblue] (2.6,2.2) circle (0.08) node[below right,popdark]{Douala};
\fill[popblue] (1.5,2.6) circle (0.08) node[left,popdark]{Buea};
\fill[popblue] (2.2,5.4) circle (0.08) node[left,popdark]{Bamenda};
\fill[popblue] (4.8,8.6) circle (0.08) node[right,popdark]{Garoua};
\fill[popblue] (3.6,4.6) circle (0.08) node[right,popdark]{Foumban};
% annotations
\node[poporange,align=center] at (1.1,4.2) {\scriptsize Zone\\ \scriptsize britannique};
\node[popgreen!60!black] at (5.6,6.4) {\scriptsize Zone française};
\draw[->,thick,popink] (8.6,9.8) -- (8.6,10.6) node[above]{\scriptsize N};
\end{tikzpicture}
\legende{Croquis de la formation territoriale du Cameroun. 1 : ligne de partage de 1919 ; 2 : zone sous tutelle française, indépendante le 1\ier{} janvier 1960 ; 3 : zone sous tutelle britannique, plébiscite du 11 février 1961 (le Nord rejoint le Nigeria, le Sud rejoint le Cameroun) ; 4 : République fédérale du 1\ier{} octobre 1961 ; 5 : République unie du 20 mai 1972. Échelle indicative.}
\end{popfigure}

\textbf{Étape 2 — Le graphique économique et les calculs.} On reporte les valeurs du document 2.

\begin{popfigure}
\begin{tikzpicture}
\begin{axis}[width=0.85\linewidth,height=6cm,xlabel={Année},ylabel={PIB/habitant (USD)},xmin=1958,xmax=2022,ymin=0,ymax=1700,xtick={1960,1970,1980,1990,2000,2010,2020},grid=both,grid style={popline!40},legend pos=north west]
\addplot[popcurve,mark=*,mark size=1.5pt] coordinates {(1960,70)(1970,160)(1980,890)(1986,1000)(1990,960)(1994,630)(2000,700)(2010,1250)(2020,1500)};
\legend{PIB par habitant}
\end{axis}
\end{tikzpicture}
\legende{Évolution du PIB par habitant au Cameroun, 1960-2020 (dollars courants, valeurs indicatives du dossier).}
\end{popfigure}

Taux de variation entre 1960 et 1980 :
\begin{align*}
t &= \frac{890 - 70}{70} \times 100 = \frac{820}{70} \times 100 \approx 1171\,\%
\end{align*}
Taux de variation entre 1986 et 1994 :
\begin{align*}
t &= \frac{630 - 1000}{1000} \times 100 = -37\,\%
\end{align*}
On distingue trois phases : \textbf{phase 1 (1960-1986)} : forte croissance, portée par le boom pétrolier à partir de 1977-1978 ; \textbf{phase 2 (1986-1994)} : crise économique grave, chute du PIB par habitant de 37 \%, ajustement structurel puis dévaluation du franc CFA en janvier 1994 ; \textbf{phase 3 (depuis 1994-2000)} : reprise lente et irrégulière.

\textbf{Étape 3 — Analyse des discours.} Document 3a : Ahidjo, juillet 1961, conférence de Foumban, contexte de la réunification ; idée principale : bâtir l'unité nationale. Document 3b : Biya, novembre 1982, investiture après la démission d'Ahidjo ; idée principale : continuité de la construction nationale, avec « rigueur et moralisme ». Objectif commun : consolider un État national unifié et stable.

\textbf{Étape 4 — Lecture des scrutins.} En 1961, le plébiscite partage l'ancien Cameroun britannique : le Nord choisit le Nigeria, le Sud choisit le Cameroun, ce qui fixe les frontières actuelles. En 1992, le score serré (environ 40 \% contre 36 \%) montre que le multipartisme rétabli en 1990 a fait de l'élection une compétition réelle, même si les résultats ont été contestés.

\textbf{Étape 5 — Le texte argumenté (modèle).} \emph{Le Cameroun indépendant est une réussite réelle mais inachevée. Sur le plan politique, le pays a réussi sa réunification : le 11 février 1961, le Cameroun méridional britannique a choisi par plébiscite de rejoindre la République du Cameroun, et la conférence de Foumban (juillet 1961) a fondé la République fédérale. En 1972, le référendum du 20 mai achève l'unité de l'État. Sur le plan économique, le PIB par habitant est passé d'environ 70 dollars en 1960 à près de 1000 dollars en 1986, grâce notamment au pétrole. La succession pacifique de 1982, avec l'investiture de Paul Biya, témoigne d'une stabilité rare en Afrique. Cependant, des limites existent : la crise de 1986-1994 a fait chuter le PIB par habitant de 37 \%, et l'élection contestée de 1992 révèle une démocratie encore fragile. Le Cameroun indépendant est donc une réussite à consolider.}
\tcblower
La solution ci-dessus suit la démarche exigée : croquis avec légende ordonnée, calculs exacts des taux de variation, identification des trois phases économiques, analyse contextualisée des discours et des scrutins, puis texte argumenté structuré (thèse nuancée, quatre faits datés au moins, conclusion). Veille à ce que chaque production porte un titre et que la parole soit répartie entre les quatre membres du groupe lors de la restitution orale.
\end{exoresolu}

\courssub{Bilan de l'activité}

Cette activité a permis de mettre en évidence trois idées essentielles de la leçon. D'abord, l'État camerounais actuel est le produit d'un \cle{processus} et non d'un acte unique : tutelles, indépendance de 1960, plébiscites de 1961, fédération, puis unité en 1972. Ensuite, l'histoire économique du pays se lit en \cle{phases} : la prospérité pétrolière de l'ère Ahidjo, la crise de 1986-1994, puis la reprise sous l'ère Biya ; les données chiffrées rendent ces phases visibles et mesurables. Enfin, l'exercice du texte argumenté montre qu'un \cle{bilan historique} n'est jamais un simple « oui » ou « non » : il exige une position justifiée par des faits datés et ouverte à la nuance. Ces compétences — construire un croquis, exploiter des statistiques, argumenter par écrit — sont directement mobilisables à l'épreuve d'Histoire du Probatoire et du Baccalauréat technique.

\newpage

\coursec{Méthodes et exercices résolus}

\coursec{Le Cameroun indépendant}

\begin{popfigure}
\begin{tikzpicture}[
    scale=0.9,
    every node/.style={font=\footnotesize},
    date/.style={circle, fill=popblue, text=white, minimum width=6pt, inner sep=1pt},
    event/.style={anchor=west, text width=5.5cm, align=left},
    period/.style={anchor=west, text width=5.5cm, align=left, text=popmuted, font=\footnotesize\itshape}
]
% Axe principal
\draw[popline, thick] (0,0) -- (16,0);
% Dates et événements
\foreach \x/\year/\lab/\desc [count=\i from 0] in {
    0/1960/1er janv./{Indépendance Cameroun oriental (Ahidjo)},
    1.5/1961/11 fév./{Plébiscite ONU (Sud : réunification)},
    2.5/1961/1er oct./{Réunification / République Fédérale},
    4.5/1966/1966/{Parti unique UNC (fusion UC+KNDP)},
    6.5/1972/20 mai/{Référendum État unitaire (99,99\% OUI)},
    7.5/1972/2 juin/{Constitution République Unie},
    9.5/1982/4 nov./{Démission Ahidjo / Biya Président},
    10.5/1983/juil./{Démission Ahidjo tête UNC / Biya contrôle parti},
    11.5/1984/6 avr./{Tentative coup d'État Garde Présidentielle},
    12.5/1984/1984/{Loi n°84-01 : République du Cameroun},
    13.5/1990/1990/{Loi libertés / Multipartisme (RDPC)},
    14.5/1996/1996/{Constitution actuelle (Décentralisation)},
    15.5/2008/avril/{Révision Constitution : suppression limite mandats},
    16.5/2016/2016/{Début crise anglophone (grèves avocats/enseignants)},
    17.5/2019/2019/{Grand Dialogue National / Statut spécial NO/SO},
    18.5/2025/2025/{Prochaine élection présidentielle}} {
    \node[date] (d\i) at (\x,0) {};
    \node[event, above=4pt of d\i, rotate=45, anchor=south west] {\textbf{\year} : \lab};
    \node[period, below=4pt of d\i, rotate=-45, anchor=north west] {\desc};
}
% Bornes
\node[below=1.2cm of d0, anchor=north, font=\small\bfseries] {1960};
\node[below=1.2cm of d15, anchor=north, font=\small\bfseries] {2025};
\end{tikzpicture}
\legende{Frise chronologique synthétique : les grandes dates du Cameroun indépendant (1960--2025).}
\end{popfigure}

\courssub{Méthodes et exercices résolus}

% --- PAIRE 1 : Indépendances et Réunification ---
\begin{methode}[Analyser un processus d'accession à la souveraineté et à l'unité nationale]
\textbf{Objectif :} Maîtriser la chronologie et les acteurs clés du passage de la tutelle à l'indépendance, puis de la fédération à l'unité (compétence : \emph{Appliquer les notions... / Rédiger et justifier}).
\par
\textbf{Démarche en 4 étapes :}
\begin{enumerate}
    \item \textbf{Identifier les dates charnières :} 1er janvier 1960 (Indépendance Cameroun oriental), 11 février 1961 (Plébiscite ONU), 1er octobre 1961 (Réunification / Indépendance Cameroun occidental), 20 mai 1972 (Référendum État unitaire).
    \item \textbf{Distinquer les deux espaces :} Cameroun oriental (ex-français, Ahidjo) et Cameroun occidental (ex-britannique, Foncha/Endeley). Noter le rôle de l'ONU (plébiscite février 1961).
    \item \textbf{Analyser le texte constitutionnel :} Constitution fédérale du 1er septembre 1961 (État fédéral, deux assemblées, un président fédéral). Repérer les articles clés (ex: Art. 1, Art. 14).
    \item \textbf{Structurer la réponse :} Introduction (contexte), Développement chronologique (Indépendance $\rightarrow$ Réunification), Conclusion (héritage : dualité linguistique/juridique).
\end{enumerate}
\textbf{Réflexe Bac :} Toujours citer le \textbf{plébiscite du 11 février 1961} comme acte juridique fondateur de la réunification, et non seulement la date du 1er octobre.
\end{methode}

\begin{exoresolu}[Commentaire de document : La proclamation de l'indépendance et la réunification]
\textbf{Énoncé :} Vous disposez des extraits suivants :
\begin{itemize}
    \item \textbf{Doc 1 :} Extrait du discours d'Ahmadou Ahidjo, 1er janvier 1960 : « Le Cameroun est né à la vie internationale... »
    \item \textbf{Doc 2 :} Résultats du plébiscite du 11 février 1961 (ONU) : Nord : 60\% pour le Nigeria / 40\% pour le Cameroun ; Sud : 70,5\% pour le Cameroun / 29,5\% pour le Nigeria.
    \item \textbf{Doc 3 :} Article 1 de la Constitution fédérale du 1er septembre 1961 : « Le Cameroun est une République fédérale... composée de l'État de l'Est et de l'État de l'Ouest. »
\end{itemize}
\textbf{Questions :}
\begin{enumerate}
    \item Présentez la nature et l'origine des documents.
    \item Expliquez pourquoi le 1er janvier 1960 ne marque pas l'indépendance de tout le Cameroun actuel.
    \item Montrez comment le plébiscite de 1961 a scellé le destin du Cameroun occidental.
    \item En vous appuyant sur le Doc 3, précisez la forme de l'État issue de la réunification.
\end{enumerate}
\tcblower
\textbf{Solution détaillée :}
\begin{enumerate}
    \item \textbf{Nature/Origine :}
    \begin{itemize}
        \item Doc 1 : Discours officiel (source primaire, discours politique), auteur : Ahmadou Ahidjo (Premier Ministre puis Président), date : 1er janv. 1960, lieu : Yaoundé. Contexte : Proclamation de l'indépendance du Cameroun oriental (ex-tutelle française).
        \item Doc 2 : Tableau statistique / Rapport officiel (source primaire, administrative), origine : ONU (organe de tutelle), date : 11 fév. 1961. Contexte : Plébiscite organisé conformément à la résolution 1350 (XIV) de l'AG ONU.
        \item Doc 3 : Texte constitutionnel (source primaire, juridique), origine : Conférence constitutionnelle de Foumban (juillet 1961), promulguée 1er sept. 1961.
    \end{itemize}
    \item \textbf{1er janvier 1960 : indépendance partielle.} Cette date ne concerne que le \textbf{Cameroun oriental} (ex-Cameroun sous tutelle française), devenu République du Cameroun. Le \textbf{Cameroun occidental} (ex-Cameroun britannique, administré depuis le Nigeria) est encore sous tutelle britannique. Il attend l'application de la résolution de l'ONU prévoyant un plébiscite pour choisir son avenir (intégration au Nigeria ou réunification au Cameroun indépendant).
    \item \textbf{Rôle du plébiscite (11 fév. 1961).} Le Doc 2 montre un vote contrasté :
    \begin{itemize}
        \item \textbf{Nord (Nord-Ouest actuel) :} Majorité pour l'intégration au Nigeria (60\%). Conséquence : Rattachement au Nigeria le 1er juin 1961 (perte territoriale pour le projet national).
        \item \textbf{Sud (Sud-Ouest actuel) :} Majorité nette pour la réunification (70,5\%). Conséquence : Indépendance par la réunification au Cameroun oriental le 1er octobre 1961.
    \end{itemize}
    C'est l'acte de libre autodétermination qui valide juridiquement la réunification du Sud-Ouest.
    \item \textbf{Forme de l'État (Doc 3).} L'article 1 institue une \textbf{République fédérale}. Deux entités fédérées : \textbf{État de l'Est} (ex-Oriental, capitale Yaoundé) et \textbf{État de l'Ouest} (ex-Occidental, capitale Buea). Pouvoirs partagés : Président fédéral (Ahidjo), Assemblée fédérale, mais chaque État garde son exécutif (Premier Ministre), son assemblée, sa fonction publique, son système éducatif/juridique (Common Law vs Droit civil).
\end{enumerate}
\textbf{Conclusion :} La réunification (1er oct. 1961) est un acte négocié (Conférence de Foumban) sanctionné par un plébiscite onusien, aboutissant à un fédéralisme asymétrique (deux États aux historiques juridiques distincts) qui durera jusqu'au référendum du 20 mai 1972.
\end{exoresolu}

% --- PAIRE 2 : Construction de l'État Unitaire (Ahidjo) ---
\begin{methode}[Expliquer la centralisation du pouvoir et la naissance de l'État unitaire]
\textbf{Objectif :} Comprendre la logique politique (stabilité, unité) et les étapes juridiques (1966, 1972) de la transition fédérale $\rightarrow$ unitaire (compétence : \emph{Rédiger et justifier}).
\par
\textbf{Démarche en 4 étapes :}
\begin{enumerate}
    \item \textbf{Contextualiser 1966 :} Crise UPC (rébellion), pluralisme partisan perçu comme facteur de division. Création du parti unique \textbf{UNC} (fusion KNDP + UC) $\rightarrow$ contrôle politique total.
    \item \textbf{Analyser le Référendum du 20 mai 1972 :} Question unique (« Voulez-vous que le Cameroun soit un État unitaire ? »). Résultats officiels : 99,99\% OUI. Adoption Constitution 2 juin 1972.
    \item \textbf{Lister les suppressions :} Poste de Vice-Président (Foncha), Assemblées fédérées, Premiers Ministres fédérés, Gouverneurs remplacent les PM, Cour suprême unique, unique parti (UNC $\rightarrow$ RDPC 1985).
    \item \textbf{Justifier le « OUI » massif :} Argument officiel (unité, développement, fin du particularisme) vs Réalité (contrôle militaire, absence d'opposition, pression administrative).
\end{enumerate}
\textbf{Formule à retenir :} \textbf{Fédéralisme (1961) $\xrightarrow[\text{1966: Parti unique}]{\text{1972: Référendum}}$ État Unitaire (1972) $\xrightarrow{\text{1984: Loi n°84-01}}$ République du Cameroun.}
\end{methode}

\begin{exoresolu}[Dissertation : Les étapes de la construction de l'État unitaire sous Ahidjo (1961-1972)]
\textbf{Énoncé :} « De la Fédération à l'Unité : analysez les motivations et les modalités de la centralisation du pouvoir au Cameroun sous la présidence d'Ahmadou Ahidjo (1961-1972). »
\tcblower
\textbf{Solution détaillée (Plan structuré) :}
\textbf{Introduction :}
\begin{itemize}
    \item \textbf{Accroche :} 1er octobre 1961, naissance de la République Fédérale du Cameroun. 20 mai 1972, référendum actant la mort de la Fédération.
    \item \textbf{Définition :} État fédéral (dualisme institutionnel) vs État unitaire (centralisation).
    \item \textbf{Problématique :} Comment et pourquoi le régime d'Ahidjo a-t-il démantelé l'édifice fédéral en une décennie ?
    \item \textbf{Annonce du plan :} I. La consolidation du pouvoir personnel et partisan (1961-1966). II. L'abolition référendaire de la Fédération (1972). III. Les conséquences institutionnelles majeures.
\end{itemize}
\textbf{Développement :}
\textbf{I. La consolidation du pouvoir personnel et partisan (1961-1966)}
\begin{itemize}
    \item \textbf{Héritage fédéral fragile :} Deux États (Est/Ouest), deux systèmes (droit civil/Common law, français/anglais), deux Premiers Ministres (Ahidjo/Foncha). Vice-présidence réservée à l'Ouest (Foncha).
    \item \textbf{Lutte contre l'opposition (UPC) :} Rébellion maquisarde (Bassa, Bamileke) perçue comme menace séparatiste. Répression militaire (1960-1970). Argument : l'unité nationale prime sur les libertés.
    \item \textbf{1966 : Fusion des partis $\rightarrow$ UNC.} KNDP (Foncha, Ouest) + UC (Ahidjo, Est) = Union Nationale Camerounaise. Parti unique \textbf{de facto}. Fin du pluralisme, musellement de l'opposition parlementaire (ex: affaire Okala, 1962).
\end{itemize}
\textbf{II. L'abolition référendaire de la Fédération (1972)}
\begin{itemize}
    \item \textbf{Contexte :} Démission de Foncha (1970), remplacé par Muna. Tensions sur la gestion des finances fédérales. Ahidjo juge la fédération « coûteuse » et « source de division ».
    \item \textbf{20 mai 1972 :} Référendum sur l'État unitaire. Campagne « OUI » massive (moyens d'État, pas de campagne « NON »). Résultat : 99,99\% OUI (chiffre officiel).
    \item \textbf{2 juin 1972 :} Promulgation Constitution de la République Unie du Cameroun. Disparition : Vice-Présidence, Assemblées fédérées, Premiers Ministres fédérés. Création : Poste de Premier Ministre unique (nommé par le Président), Gouverneurs de provinces (nommés par le Président), Assemblée nationale unique.
\end{itemize}
\textbf{III. Conséquences institutionnelles majeures}
\begin{itemize}
    \item \textbf{Centralisation absolue :} Tout pouvoir émane de Yaoundé (Président). Administration unique, justice unique (Cour Suprême), parti unique (UNC).
    \item \textbf{Gestion de la diversité :} Maintien théorique du bilinguisme et du bijuridisme (Art. 1er Const 1972), mais érosion pratique (nomination de francophones à la tête de l'Ouest, harmonisation législative).
    \item \textbf{Stabilité vs Autoritarisme :} Fin des sécessions (UPC vaincue), mais verrouillage démocratique.
\end{itemize}
\textbf{Conclusion :} La construction de l'État unitaire répond à une double logique : \textbf{sécuritaire} (fin de la rébellion UPC, peur du séparatisme occidental) et \textbf{personnelle} (concentration des pouvoirs par Ahidjo). Le référendum de 1972 est l'acte juridique final d'une centralisation entamée dès 1961. Héritage : un État fort, centralisé, mais fragilisé par l'absence de contre-pouvoirs.
\end{exoresolu}

% --- PAIRE 3 : Bilan Économique et Social (Ère Ahidjo) ---
\begin{methode}[Évaluer un bilan économique et social sur une période longue]
\textbf{Objectif :} Nuancer le discours officiel (« développement ») par l'analyse chiffrée et structurelle (compétence : \emph{Appliquer les notions... / Rédiger et justifier}).
\par
\textbf{Démarche en 4 étapes :}
\begin{enumerate}
    \item \textbf{Séparer les deux sous-périodes :} \textbf{1960-1977} (Croissance, investissements, planification) vs \textbf{1977-1982} (Début crise, endettement, pétrole).
    \item \textbf{Utiliser les indicateurs macroéconomiques :} Taux de croissance PIB, part secteur primaire/secondaire/tertiaire, dette extérieure, balance commerciale.
    \item \textbf{Analyser les politiques sectorielles :} Agriculture (SODECAO, SEMRY, UNVDA), Industrie (SOCAPALM, HEVECAM, ALUCAM, SONARA), Infrastructures (Barrage Lagdo, Port Kribi, Transcam).
    \item \textbf{Évaluer le social :} Santé (gratuité relative, campagne vaccination), Éducation (scolarisation, université Yaoundé 1962), Emploi (fonctionnaire = modèle), Inégalités (urbain/rural, Est/Ouest).
\end{enumerate}
\textbf{Piège à éviter :} Ne pas dire « tout était parfait » ni « tout a échoué ». Utiliser \textbf{« Malgré des performances réelles (infrastructures, scolarisation), les bases structurelles restaient fragiles (dépendance extérieures, mono-culture) »}.
\end{methode}

\begin{exoresolu}[Analyse de données : L'économie camerounaise sous Ahidjo (1960-1982)]
\textbf{Énoncé :} À partir du tableau ci-dessous et de vos connaissances, rédigez une réponse argumentée de 20 lignes maximum : « Le bilan économique de l'ère Ahidjo est-il positif ? »
\begin{center}
\begin{tabularx}{\linewidth}{|l|X|X|X|}\hline
\textbf{Indicateur} & \textbf{1960} & \textbf{1977 (Pic)} & \textbf{1982 (Fin règne)} \\ \hline
PIB/habitant (USD) & $\sim$ 150 & $\sim$ 1\,200 & $\sim$ 1\,100 \\ \hline
Taux croissance annuel & -- & +7\% (moy. 70-77) & -2\% (1981-82) \\ \hline
Part Pétrole (Export) & 0\% & 10\% & 45\% \\ \hline
Dette Extérieure (Mds F CFA) & Faible & 300 & 800 \\ \hline
Scolarisation Primaire & 45\% & 95\% & 100\% (objectif) \\ \hline
\end{tabularx}
\end{center}
\tcblower
\textbf{Solution détaillée :}
\textbf{Thèse :} Le bilan est \textbf{contrasté} : une « décennie glorieuse » (1960-1977) masquant des failles structurelles explosant dès 1978.
\par
\textbf{Arguments positifs (1960-1977) :}
\begin{itemize}
    \item \textbf{Croissance forte :} PIB/hab $\times$ 8 (150 $\rightarrow$ 1200 USD). Moyenne 5-7\%/an. « Miracle camerounais » salué par BM/FMI.
    \item \textbf{Transformation structurelle :} Lancement industrialisation (ALUCAM 1965, SONARA 1981, agro-industrie : SOCAPALM, HEVECAM, CDC). Infrastructures lourdes (Lagdo, Transcam, Kribi).
    \item \textbf{Social :} Scolarisation primaire quasi universelle (95\% $\rightarrow$ 100\%). Création Univ. Yaoundé (1962). Politique de « camerounisation » des cadres (fonction publique).
\end{itemize}
\textbf{Arguments négatifs / Limites structurelles :}
\begin{itemize}
    \item \textbf{Dépendance extravertie :} Économie de rente (Café, Cacao, Bois, puis Pétrole 45\% export 1982). Vulnérabilité aux cours mondiaux (choc café/cacao 1977).
    \item \textbf{Endettement galopant :} Dette $\times$ 2,6 (300 $\rightarrow$ 800 Mds F CFA 1977-1982). Financement de projets à faible rentabilité (Lagdo, usines étatiques déficitaires).
    \item \textbf{Déséquilibres :} Négligence secteur vivrier (insécurité alimentaire), inégalités régionales (Est vs Ouest, Urbain vs Rural), chômage diplômés fin période.
    \item \textbf{Gouvernance :} Gestion opaque (pas de Cour des comptes), népotisme, entreprises publiques déficitaires (subventions budgétaires).
\end{itemize}
\textbf{Conclusion :} Ahidjo a posé les \textbf{infrastructures de base} et formé des \textbf{cadres}, mais le modèle « État-entrepreneur » rentier (pétrole) a masqué l'absence de diversification réelle, préparant la crise structurelle de 1987-1994 sous Biya.
\end{exoresolu}

% --- PAIRE 4 : Succession 1982 et Avènement Biya ---
\begin{methode}[Analyser une transition politique pacifique mais conflictuelle]
\textbf{Objectif :} Expliquer les mécanismes constitutionnels et les luttes de pouvoir « silencieuses » de 1982-1984 (compétence : \emph{Rédiger et justifier}).
\par
\textbf{Démarche en 4 étapes :}
\begin{enumerate}
    \item \textbf{Respecter la Constitution (1972) :} Art. 6 : Vice-Président (ou Premier Ministre) assure l'intérim. Ahidjo nomme Biya (PM depuis 1975) Vice-Président (août 1982) $\rightarrow$ Démission Ahidjo (4 nov 1982) $\rightarrow$ Biya Président (6 nov 1982). \textbf{Légalité formelle respectée.}
    \item \textbf{Identifier le « Double pouvoir » (nov 1982 - août 1983) :} Ahidjo reste Président UNC (parti), Biya Président République. Conflit de légitimité : Ahidjo contrôle le parti (base), Biya l'État (armée, admin).
    \item \textbf{Chronologuer la rupture :} Février 1983 (Voyage Ahidjo en France « santé »), Juillet 1983 (Démission Ahidjo présidence UNC), Août 1983 (Complot « coup d'État » $\rightarrow$ Arrestations), 1984 (Procès Ahidjo/Bello Bouba condamnés par contumace), 6 avril 1984 (Tentative coup d'État Garde Présidentielle).
    \item \textbf{Expliquer la victoire de Biya :} Contrôle armée (général Kadji), soutien France (Mitterrand), habileté politique (ouverture « rigueur/moralisation », libération prisonniers politiques), élimination rivaux (Bello, Sadou, Maikano).
\end{enumerate}
\textbf{Mots-clés :} \cle{Transition constitutionnelle} vs \cle{Lutte de succession}.
\end{methode}

\begin{exoresolu}[Commentaire composé : La succession de 1982, une « transition modelée » ?]
\textbf{Énoncé :} « La passation de pouvoir entre Ahmadou Ahidjo et Paul Biya en 1982 est souvent présentée comme un modèle de transition pacifique en Afrique. Nuancez cette affirmation en analysant la période novembre 1982 - avril 1984. »
\tcblower
\textbf{Solution détaillée :}
\textbf{Introduction :} 4 nov 1982 : Démission surprise Ahidjo (santé). 6 nov : Biya investi (Art. 6 Const). Apparence : légalité, calme, continuité. Réalité : crise de légitimité larvée $\rightarrow$ confrontation ouverte (1984).
\par
\textbf{1. La légalité formelle : un atout majeur.}
\begin{itemize}
    \item Respect strict Art. 6 (Vice-Président $\rightarrow$ Président). Pas de vide pouvoir, pas de militaire au pouvoir.
    \item Biya : technocrate, Premier Ministre depuis 1975, perçu comme loyal, modéré, « homme du Nord » (région d'Ahidjo) $\rightarrow$ rassure l'establishment.
    \item Soutien international immédiat (France, ONU, OUA).
\end{itemize}
\textbf{2. Le « Diarchisme » conflictuel (Nov 1982 - Juil 1983).}
\begin{itemize}
    \item \textbf{Deux centres de pouvoir :} Ahidjo (Président UNC, résidence Mouanko) vs Biya (Président République, Palais Etoudi).
    \item \textbf{Conflits concrets :} Nominations (Biya nomme ses proches, Ahidjo s'oppose via parti), Gestion finances (Ahidjo contrôle trésoriers UNC), Symbolique (portraits Ahidjo maintenus).
    \item \textbf{Base sociale :} Ahidjo garde fidélité Nord, armée (certains), vieux barons. Biya construit sa base (Sud, technocrates, jeunesse, presse « liberté »).
\end{itemize}
\textbf{3. La rupture définitive (Été 1983 - Avril 1984).}
\begin{itemize}
    \item \textbf{Juillet 1983 :} Démission Ahidjo présidence UNC (pression Biya/France). Biya devient Président UNC $\rightarrow$ Fusion total Parti/État.
    \item \textbf{Août 1983 :} « Complot » découvert. Arrestation Bello Bouba, Maikano, Sadou Daoudou (pro-Ahidjo). Ahidjo s'exile en France.
    \item \textbf{Février 1984 :} Procès Yaoundé. Ahidjo, Bello, Maikano condamnés à mort (par contumace pour Ahidjo). Grâce présidentielle pour les autres.
    \item \textbf{6 Avril 1984 :} Tentative coup d'État Garde Présidentielle (majoritairement Nordistes). Répression sanglante (500-1000 morts). Fin de l'opposition armée interne.
\end{itemize}
\textbf{Conclusion :} Transition \textbf{juridiquement parfaite}, mais \textbf{politiquement violente}. Biya a dû « tuer le père » (symboliquement et physiquement via la répression) pour asseoir son autorité. Le « modèle » cache une guerre de palais et une purge ethnico-régionale (Nord vs Sud) aux séquelles durables.
\end{exoresolu}

% --- PAIRE 5 : Mutations Politiques (1982 - nos jours) ---
\begin{methode}[Analyser l'évolution d'un régime autoritaire vers un pluralisme encadré]
\textbf{Objectif :} Maîtriser la chronologie des réformes (1990, 1996, 2008) et la permanence du pouvoir personnel (compétence : \emph{Appliquer les notions...}).
\par
\textbf{Démarche en 4 étapes :}
\begin{enumerate}
    \item \textbf{Périodiser :} \textbf{1982-1990 :} « Rigueur/Moralisation » + Parti unique (RDPC 1985). \textbf{1990-1996 :} Ouverture (Loi 90/052 libertés, 90/053 partis), Élections 1992 (présidentielles), 1997 (boycott), Constitution 1996 (décentralisation, Sénat, Cour Suprême $\rightarrow$ Conseil Constitutionnel). \textbf{1997-2008 :} « Stabilité » (victoires 97, 2004, 2008 révision suppression limite mandats). \textbf{2008-... :} Crise anglophone (2016), Grand Dialogue National (2019), Code électoral 2024, Élections 2025.
    \item \textbf{Distinquer forme et fond :} Multipartisme \textbf{de jure} (300+ partis) vs Hégémonie RDPC \textbf{de facto} (contrôle admin, médias, finances, découpage électoral, Conseil Constitutionnel).
    \item \textbf{Analyser les institutions clés :} Président (mandat 7 ans, renouvelable $\rightarrow$ 2008 illimité), Parlement (Assemblée + Sénat 2013), Justice (non indépendance perçue), Administration territoriale (Gouverneurs, Préfets, DO = bras armé du pouvoir).
    \item \textbf{Citer les crises majeures :} Crise anglophone (revendications fédéralisme/sécession), Boko Haram (Extrême-Nord), Crise post-électorale 2018 (M. Kamto).
\end{enumerate}
\textbf{Repères :} \textbf{1990} (Libertés), \textbf{1996} (Constitution actuelle), \textbf{2008} (Suppression limite mandats), \textbf{2016} (Crise NO/SO), \textbf{2019} (GDN), \textbf{2025} (Prochaine présidentielle).
\end{methode}

\begin{exoresolu}[Étude de cas : La révision constitutionnelle de 2008]
\textbf{Énoncé :} En avril 2008, le Parlement (Assemblée Nationale) vote la révision de l'article 6 alinéa 2 de la Constitution de 1996, supprimant la limitation à deux mandats présidentiels. Analysez les arguments des partisans et des opposants, et les conséquences politiques de cette révision.
\tcblower
\textbf{Solution détaillée :}
\textbf{Contexte :} Février 2008 : Émeutes de la faim (prix carburant, vie chère). Répression dure. Avril 2008 : Session parlementaire extraordinaire. RDPC majorité absolue (153/180 députés).
\par
\textbf{Arguments des Partisans (Majorité RDPC / Gouvernement) :}
\begin{enumerate}
    \item \textbf{Continuité / Stabilité :} Biya = « Homme de la paix », « Architecte du développement ». Changer = risque instabilité (ex: 1984, 1990).
    \item \textbf{Expérience :} Gestion dossiers complexes (Pétrole, Dette, Crise anglophone naissante, Boko Haram). « L'expérience ne se vote pas ».
    \item \textbf{Légalité :} Procédure respectée (Art. 13 Const 1996 : initiative Président/Parlement, vote 3/5èmes). Pas de référendum obligatoire (sauf si Président le décide).
    \item \textbf{Souveraineté :} Rejet ingérence occidentale (UE, France critiques). « Le peuple camerounais décide ».
\end{enumerate}
\textbf{Arguments des Opposants (Opposition SDF, UDC, MRC, Société civile, Évêques) :}
\begin{enumerate}
    \item \textbf{Coup de force constitutionnel :} « Constitution taillée sur mesure ». Violation esprit 1996 (consensus national). « Monarchie républicaine ».
    \item \textbf{Démocratie bloquée :} Verrouillage alternance. Élections sans enjeu (2011, 2018). Démoralisation électorale (abstention forte).
    \item \textbf{Concentration des pouvoirs :} Président = Chef État + Chef Gouvernement (depuis 1984, PM subordonné) + Chef Parti + Chef Armées + Nomination tous hauts postes + Contrôle Conseil Constitutionnel. Absence contre-pouvoirs.
    \item \textbf{Crise de légitimité :} Émeutes 2008 = rejet social. Réponse par la force + révision = déni de démocratie.
\end{enumerate}
\textbf{Conséquences politiques :}
\begin{itemize}
    \item \textbf{Ancrage du pouvoir :} Biya candidat 2011 (victoire 78\%), 2018 (71\%). Longévité record (42 ans+ au pouvoir en 2024).
    \item \textbf{Radicalisation opposition :} Création MRC (Maurice Kamto, 2012), boycotts, contestation rue (2018, 2020).
    \item \textbf{Crise anglophone exacerbée :} Perception « Yaoundé ne réforme pas, il verrouille ». Revendications fédéralisme $\rightarrow$ sécession (2017).
    \item \textbf{Relations internationales :} Tensions ponctuelles (UE, USA droits homme), mais partenariats stratégiques maintenus (sécurité, économie).
\end{itemize}
\textbf{Conclusion :} La révision 2008 est l'acte fondateur de la \textbf{présidence à vie} déguisée. Elle a assuré la stabilité du régime (pas de transition) au prix d'une \textbf{légitimité démocratique érodée} et d'une conflictualité accrue (Anglophone, Post-électoral).
\end{exoresolu}

% --- PAIRE 6 : Défis Contemporains ---
\begin{methode}[Structurer une analyse des défis multidimensionnels d'un pays en développement]
\textbf{Objectif :} Classer, hiérarchiser et lier les défis (sécuritaires, économiques, sociaux, institutionnels) pour une réponse de synthèse (compétence : \emph{Rédiger et justifier}).
\par
\textbf{Démarche en 4 étapes (Grille d'analyse PESTEL simplifiée) :}
\begin{enumerate}
    \item \textbf{Sécuritaire (Priorité 1) :} Crise anglophone (NO/SO : déplacés, écoles fermées, économie locale détruite), Boko Haram (Extrême-Nord : insécurité, déplacés, agriculture paralysée), Insécurité urbaine (coupeurs de route, « microbes »).
    \item \textbf{Économique / Structurel :} Sortie crise (PPTE 2006, mais dette remonte $\sim$ 45\% PIB 2023), Dépendance hydrocarbures/bois/cacao (chocs externes), Sous-industrialisation, Informel ($>$ 90\% emploi), Infrastructures déficitaires (énergie, routes), Corruption (Indice CPI faible), Vision 2035 (objectif pays émergent) vs réalité.
    \item \textbf{Social / Humain :} Démographie galopante ($\sim$ 28M hab, 60\% $<$ 25 ans), Chômage jeunes (diplômés), Santé (mortalité maternelle/infantile encore haute, paludisme, VIH, accès soins), Éducation (qualité, adéquation formation/emploi, grève enseignants), Inégalités (Gini $\sim$ 0,45), Genre.
    \item \textbf{Institutionnel / Gouvernance :} Décentralisation actée (1996, 2004, 2019 GDN, lois 2019/2020, élections régionales 2020) mais transfert compétences/ressources lent. Justice lente/partiale. Élections transparentes ? (Code électoral 2024, ELECAM). Succession présidentielle (2025 ?) : incertitude majeure.
\end{enumerate}
\textbf{Synthèse :} Lier les défis : \textbf{Insécurité $\rightarrow$ Fuite investissements / Fermeture écoles $\rightarrow$ Chômage / Radicalisation $\rightarrow$ Insécurité}. Casser le cercle vicieux = Défi central.
\end{methode}

\begin{exoresolu}[Sujet de synthèse : Les défis du Cameroun contemporain à l'horizon 2035]
\textbf{Énoncé :} « Le Cameroun aspire à devenir un pays émergent à l'horizon 2035 (Document de Stratégie pour la Croissance et l'Emploi - DSCE). Identifiez les trois défis majeurs qui entravent cette ambition et proposez, pour chacun, une piste de solution réaliste. »
\tcblower
\textbf{Solution détaillée :}
\textbf{Introduction :} Vision 2035 (adoptée 2009) : Pays émergent, démocratique, uni dans la diversité. PIB/hab $\sim$ 1500\$ (2023) $\rightarrow$ Objectif $\sim$ 4000\$. Écart immense. Trois défis systémiques.
\par
\textbf{Défi 1 : L'insécurité généralisée (Frein n°1 au développement).}
\begin{itemize}
    \item \textbf{Constat :} Deux guerres internes (Anglophone NO/SO depuis 2016, Boko Haram Extrême-Nord depuis 2014) + insécurité urbaine. Coût humain (6000+ morts, 1M+ déplacés), économique (destruction PIB régional, fermetures entreprises CDC/PAMOL, chute recettes douanières), social (génération sacrifiée : écoles fermées 4-5 ans).
    \item \textbf{Piste réaliste :} \textbf{Approche globale « Sécurité + Développement + Justice »}. Militaire seul échoue. $\rightarrow$ Application effective décentralisation (pouvoirs/ressources aux CTD), Justice transitionnelle (vérité/réparation), Programme DDR (Désarmement, Démobilisation, Réinsertion) crédible + Investissements massifs infrastructures/services dans zones sinistrées (Plan de reconstruction Présidentiel 2020+ accéléré).
\end{itemize}
\textbf{Défi 2 : La transformation structurelle de l'économie (Sortie de la rente).}
\begin{itemize}
    \item \textbf{Constat :} Croissance molle (3-4\%), insuffisante pour absorber 300 000 jeunes/an sur marché travail. Dépendance exportations brutes (Pétrole, Bois, Cacao, Aluminium) $\rightarrow$ Vulnérabilité chocs prix / change. Industrie manufacturière $<$ 15\% PIB. Informel dominant. Dette publique remontée (risque surendettement modéré FMI).
    \item \textbf{Piste réaliste :} \textbf{Industrialisation par substitution importations + ZLECAf}. Soutien ciblé PME/PMI (accès financement, foncier, électricité stable via barrage Nachtigal + renouvelables), Transformation locale matières premières (cacao $\rightarrow$ chocolat, bois $\rightarrow$ meubles, coton $\rightarrow$ textile), Amélioration climat affaires (numérisation procédures, lutte corruption effective CONAC/Tribunaux spécialisés), Formation technique alignée besoins marché (Partenariat École-Entreprise).
\end{itemize}
\textbf{Défi 3 : La gouvernance démocratique et la succession politique (Stabilité institutionnelle).}
\begin{itemize}
    \item \textbf{Constat :} Concentration pouvoirs excessive (Président 90 ans en 2025, 42 ans pouvoir). Conseil Constitutionnel non indépendant (nomination présidentielle). Code électoral 2024 perfectible (ELECAM, biométrie, vote diaspora). Risque crise succession / transition violente. Défiance citoyenne (abstention, défiance institutions).
    \item \textbf{Piste réaliste :} \textbf{Réformes consensuelles pré-électorales (Dialogue inclusif).} Réforme Code électoral (ELECAM indépendant, vote biométrique sécurisé, diaspora), Réforme Conseil Constitutionnel (nomination partagée Parlement/Société civile), Limitation mandats (rétablissement Art. 6 Const 1996), Loi sur statut opposition / Chef de file opposition. Préparer transition apaisée (Primature forte ? Vice-Présidence ?).
\end{itemize}
\textbf{Conclusion :} L'émergence 2035 est \textbf{techniquement possible} (ressources, position géo, jeunesse) mais \textbf{politiquement conditionnée}. La paix (Défi 1) est la condition \textit{sine qua non} de l'investissement (Défi 2), qui nécessite la confiance (Défi 3). Le temps presse : fenêtre démographique (dividende) se referme vers 2040.
\end{exoresolu}

% --- BOÎTE ATTENTION FINALE ---
\begin{attention}[Les 5 erreurs qui coûtent des points au Bac]
\begin{enumerate}
    \item \textbf{Confondre Fédéralisme (1961) et Unitarisme (1972) :} Ne pas citer les institutions propres à chaque régime (Vice-Présidence, Assemblées fédérées, Premiers Ministres fédérés pour le fédéralisme ; Gouverneurs, PM unique, Assemblée unique pour l'unitarisme). \textit{Astuce : Faire un petit tableau comparatif mental.}
    \item \textbf{Oublier le plébiscite du 11 février 1961 :} L'indépendance du 1er janvier 1960 ne concerne que l'Est. La réunification juridique date du plébiscite ONU (choix du Sud) + 1er octobre 1961. Sans plébiscite, pas de réunification légale.
    \item \textbf{Bilan Ahidjo : « Tout bon » ou « Tout mauvais » :} Le Bac attend de la \textbf{nuance}. Séparer 1960-1977 (croissance, infrastructures, scolarisation) et 1977-1982 (choc termes échange, endettement, début crise). Citer des chiffres (PIB/hab, Dette, Pétrole\%).
    \item \textbf{Réduction de la crise anglophone à « problème linguistique » :} C'est une crise \textbf{politique, identitaire, économique et sécuritaire} (Common Law vs Droit civil, marginalisation perçue, ressources, répression). Citer : Grèves 2016, Déclaration indépendance 2017 (Ambazonie), GDN 2019, Statut spécial 2019.
    \item \textbf{Traiter la révision 2008 comme un simple détail :} C'est l'acte majeur du règne Biya (suppression limite mandats). Il faut en expliquer la \textbf{procédure} (Parlement, pas référendum), le \textbf{contexte} (émeutes février), les \textbf{arguments} (stabilité vs verrouillage) et les \textbf{conséquences} (longévité, crise légitimité, radicalisation opposition).
\end{enumerate}
\end{attention}

\newpage

\coursec{Exercices}

\courssub{Je vérifie mes connaissances}

\begin{exercice}[QCM : repères chronologiques]
Pour chaque question, entoure la bonne réponse (une seule réponse exacte).
\begin{enumerate}
\item Le Cameroun oriental (sous administration française) accède à l'indépendance le :
\begin{enumerate}
\item 1\textsuperscript{er} janvier 1960 \item 1\textsuperscript{er} octobre 1961 \item 20 mai 1972
\end{enumerate}
\item La réunification des deux Cameroun est effective le :
\begin{enumerate}
\item 11 février 1961 \item 1\textsuperscript{er} octobre 1961 \item 4 novembre 1982
\end{enumerate}
\item La République unie du Cameroun est proclamée à la suite du référendum du :
\begin{enumerate}
\item 20 mai 1972 \item 6 avril 1984 \item 18 janvier 1996
\end{enumerate}
\item Ahmadou Ahidjo démissionne de la présidence de la République le :
\begin{enumerate}
\item 4 novembre 1982 \item 6 avril 1984 \item 1\textsuperscript{er} juin 1983
\end{enumerate}
\end{enumerate}
\end{exercice}

\begin{exercice}[Vrai ou faux — justifier]
Réponds par VRAI ou FAUX, puis justifie ta réponse en une ou deux phrases.
\begin{enumerate}
\item Les plébiscites des 11 et 12 février 1961 ont été organisés sous l'égide de l'ONU.
\item Le Cameroun septentrional (Northern Cameroons) a choisi de se rattacher au Cameroun en 1961.
\item Le parti unique au pouvoir à partir de 1966 s'appelle l'Union nationale camerounaise (UNC).
\item La tentative de coup d'État contre Paul Biya a eu lieu en avril 1984.
\end{enumerate}
\end{exercice}

\begin{exercice}[Définitions]
Définis en deux ou trois phrases chacun des termes ou expressions suivants :
\begin{enumerate}
\item Réunification ;
\item État fédéral ;
\item Référendum ;
\item Multipartisme.
\end{enumerate}
\end{exercice}

\begin{exercice}[Calculs directs : lire des données économiques]
Le tableau suivant donne l'évolution du PIB par habitant du Cameroun (en dollars constants) :

\begin{center}
\begin{tabularx}{\linewidth}{|l|X|X|X|X|}\hline
\rowcolor{popblueL} Année & 1980 & 1986 & 1990 & 1994 \\ \hline
PIB par habitant (dollars) & 950 & 1100 & 820 & 620 \\ \hline
\end{tabularx}
\end{center}

\begin{enumerate}
\item Calcule le taux de variation du PIB par habitant entre 1980 et 1986, puis entre 1986 et 1994.
\item Exprime en pourcentage la baisse du PIB par habitant entre 1986 et 1994.
\item Que traduisent ces évolutions pour la période 1986-1994 ?
\end{enumerate}
\end{exercice}

\courssub{Je m'entraîne}

\begin{exercice}[Chronologie : classer et dater]
Classe par ordre chronologique les dix événements suivants, en donnant pour chacun la date exacte (jour, mois, année quand elle est connue) :
\begin{itemize}
\item proclamation de la République unie du Cameroun ;
\item indépendance du Cameroun oriental ;
\item tentative de coup d'État contre Paul Biya ;
\item plébiscites du 11 février 1961 ;
\item création de l'Union nationale camerounaise (parti unique) ;
\item démission d'Ahmadou Ahidjo ;
\item promulgation de la Constitution de janvier 1996 ;
\item réunification effective des deux Cameroun ;
\item retour au multipartisme (lois sur les associations et partis politiques) ;
\item premières élections régionales.
\end{itemize}
Présente ta réponse sous la forme d'une frise chronologique légendée.
\end{exercice}

\begin{exercice}[Question à réponse construite : la succession de 1982]
En une quinzaine de lignes, présente les circonstances de l'arrivée de Paul Biya au pouvoir en novembre 1982 : tu préciseras le rôle de la Constitution, la personnalité du successeur désigné et les tensions qui ont suivi la passation de pouvoir (1983-1984).
\end{exercice}

\begin{exercice}[Du fédéralisme à l'unité : compléter un tableau]
Recopie et complète le tableau comparatif suivant à l'aide de tes connaissances :

\begin{center}
\begin{tabularx}{\linewidth}{|l|X|X|}\hline
\rowcolor{popblueL} Critère & République fédérale du Cameroun (1961-1972) & République unie du Cameroun (1972-1984) \\ \hline
Nombre d'États fédérés & & \\ \hline
Institutions propres aux États & & \\ \hline
Symbole de l'unité (drapeau, hymne) & & \\ \hline
Acte fondateur & & \\ \hline
\end{tabularx}
\end{center}

Conclus en deux phrases sur la signification du référendum du 20 mai 1972.
\end{exercice}

\begin{exercice}[Commentaire de données chiffrées : la crise économique]
Voici l'évolution de quelques indicateurs du Cameroun :

\begin{center}
\begin{tabularx}{\linewidth}{|l|X|X|X|}\hline
\rowcolor{popblueL} Indicateur & 1985 & 1990 & 1994 \\ \hline
Croissance du PIB (\%) & +7,0 & $-$6,0 & $-$3,0 \\ \hline
Prix du pétrole (recettes, indice 100 en 1985) & 100 & 55 & 45 \\ \hline
Salaires de la fonction publique (indice 100 en 1985) & 100 & 100 & 30 \\ \hline
\end{tabularx}
\end{center}

\begin{enumerate}
\item Décris l'évolution de chacun des trois indicateurs entre 1985 et 1994.
\item Relève deux causes de la crise économique camerounaise des années 1986-1994.
\item Dégage deux conséquences sociales de cette crise (baisse des salaires, chômage, paupérisation).
\end{enumerate}
\end{exercice}

\courssub{Je me prépare au Bac}

\begin{exercice}[Dissertation type Baccalauréat]
\textbf{Sujet :} « Le Cameroun de 1960 à 1982 : construction d'un État et développement économique. » Discute cette affirmation.

\textbf{Consignes :}
\begin{enumerate}
\item Rédige une introduction comportant une accroche, la présentation du sujet, sa problématisation et l'annonce du plan.
\item Développe deux grandes parties : (I) la construction de l'État camerounais (indépendances, réunification, passage à l'État unitaire, parti unique) ; (II) le développement économique et ses limites (planification, agriculture, pétrole, industrialisation, déséquilibres).
\item Rédige une conclusion qui fait le bilan et ouvre sur la succession de 1982.
\end{enumerate}
Ton devoir sera noté sur la pertinence de la problématique, l'exactitude des faits et des dates, la qualité de l'argumentation et la clarté de l'expression.
\end{exercice}

\begin{exercice}[Étude critique de documents : du fédéralisme à l'unité]
\textbf{Document 1 — Extrait du discours d'Ahmadou Ahidjo à la conférence de Foumban, juillet 1961 :} « Nous avons voulu un État fédéral parce que nos deux territoires ont vécu quarante ans d'administrations différentes... La fédération respecte les particularismes de chacun tout en construisant la nation camerounaise. »

\textbf{Document 2 — Extrait de la Constitution du 2 juin 1972 :} « La République du Cameroun est une et indivisible... Elle adopte le nom de République unie du Cameroun. »

\begin{enumerate}
\item Présente chaque document : nature, auteur, date, contexte.
\item Explique pourquoi le fédéralisme a été choisi en 1961 (document 1).
\item Montre, à partir du document 2, en quoi la Constitution de 1972 rompt avec le choix de 1961.
\item À l'aide de tes connaissances, explique les raisons invoquées par Ahidjo pour justifier le passage à l'État unitaire (coût des institutions fédérales, unité nationale, efficacité de l'État).
\item Rédige une conclusion d'une dizaine de lignes : le passage du fédéralisme à l'unité est-il une rupture ou l'aboutissement d'un processus ?
\end{enumerate}
\end{exercice}

\begin{exercice}[Travail pratique noté : croquis cartographique]
À partir d'une carte muette du Cameroun et du fond d'informations ci-dessous, réalise un croquis légendé intitulé « Les défis du Cameroun contemporain ».

\textbf{Fond d'informations :}
\begin{itemize}
\item Extrême-Nord : insécurité liée au groupe Boko Haram depuis 2014 ; réfugiés nigérians et déplacés internes (ville de Maroua, camp de Minawao) ;
\item Nord-Ouest et Sud-Ouest : crise anglophone depuis 2016-2017 (villes de Bamenda et Buéa) ;
\item Est : afflux de réfugiés centrafricains ;
\item Littoral : pôle économique (port de Douala), enclavement des régions septentrionales ;
\item Axes de désenclavement : corridor Douala-N'Djaména, chemin de fer Transcamerounais.
\end{itemize}

\begin{enumerate}
\item Trace le contour simplifié du Cameroun, place la flèche du nord et une échelle indicative.
\item Reporte les informations à l'aide de trois ou quatre figurés (zones colorées, points, flèches).
\item Rédige une légende organisée et numérotée.
\item En cinq lignes, dégage du croquis les principaux défis du Cameroun contemporain (sécuritaires, politiques, économiques).
\end{enumerate}
\end{exercice}

\newpage

\coursec{Corrigés des exercices}

\begin{corrige}[Exercice 1 — QCM : repères chronologiques]
\begin{enumerate}
\item \textbf{Réponse a : 1\textsuperscript{er} janvier 1960.} Le Cameroun oriental, sous administration française, accède à l'indépendance le 1\textsuperscript{er} janvier 1960 avec Ahmadou Ahidjo comme chef de l'État. Le 1\textsuperscript{er} octobre 1961 est la date de la réunification ; le 20 mai 1972 celle du référendum sur l'État unitaire.
\item \textbf{Réponse b : 1\textsuperscript{er} octobre 1961.} Après les plébiscites des 11 et 12 février 1961, le Cameroun méridional (Southern Cameroons) rejoint la République du Cameroun le 1\textsuperscript{er} octobre 1961 : naît la République fédérale du Cameroun.
\item \textbf{Réponse a : 20 mai 1972.} Le référendum du 20 mai 1972 approuve la suppression de la fédération ; la République unie du Cameroun est proclamée par la Constitution du 2 juin 1972.
\item \textbf{Réponse a : 4 novembre 1982.} Ahmadou Ahidjo démissionne le 4 novembre 1982 ; conformément à la Constitution, le Premier ministre Paul Biya lui succède le 6 novembre 1982.
\end{enumerate}
\textbf{Point de vigilance :} ne pas confondre la date des plébiscites (11-12 février 1961) avec celle de la réunification effective (1\textsuperscript{er} octobre 1961) : c'est la confusion la plus fréquente à l'examen.
\end{corrige}

\begin{corrige}[Exercice 2 — Vrai ou faux — justifier]
\begin{enumerate}
\item \textbf{VRAI.} Les plébiscites des 11 et 12 février 1961 ont été organisés sous l'égide de l'Organisation des Nations unies (ONU), qui administrait la tutelle britannique et devait trancher le sort du Cameroun occidental.
\item \textbf{FAUX.} Le Cameroun septentrional (Northern Cameroons) a choisi de se rattacher au Nigeria ; c'est le Cameroun méridional (Southern Cameroons) qui a opté pour le rattachement à la République du Cameroun.
\item \textbf{VRAI.} L'Union nationale camerounaise (UNC) est créée le 1\textsuperscript{er} septembre 1966 par fusion des partis politiques ; elle devient le parti unique. En 1985, elle prend le nom de Rassemblement démocratique du peuple camerounais (RDPC).
\item \textbf{VRAI.} La tentative de coup d'État menée par des officiers de la garde présidentielle contre Paul Biya a eu lieu les 6 et 7 avril 1984 ; elle a été réprimée et ses auteurs jugés.
\end{enumerate}
\end{corrige}

\begin{corrige}[Exercice 3 — Définitions]
\begin{enumerate}
\item \textbf{\cle{Réunification} :} réunion en un seul État de deux territoires auparavant séparés. Au Cameroun, elle désigne le rattachement, le 1\textsuperscript{er} octobre 1961, du Cameroun méridional (ex-britannique) à la République du Cameroun (ex-française), formant la République fédérale du Cameroun.
\item \textbf{\cle{État fédéral} :} État composé de plusieurs États fédérés dotés d'institutions propres (assemblée, gouvernement), qui délèguent une partie de leurs compétences à un État fédéral central. De 1961 à 1972, le Cameroun comptait deux États fédérés : le Cameroun oriental et le Cameroun occidental.
\item \textbf{\cle{Référendum} :} consultation par laquelle l'ensemble des citoyens se prononce directement, par oui ou par non, sur une question ou un texte. Exemples : les plébiscites de février 1961 et le référendum du 20 mai 1972.
\item \textbf{\cle{Multipartisme} :} régime politique dans lequel plusieurs partis peuvent légalement exister et concourir aux élections. Au Cameroun, il est rétabli par les lois de décembre 1990 sur les associations et les partis politiques, après la période du parti unique (1966-1990).
\end{enumerate}
\end{corrige}

\begin{corrige}[Exercice 4 — Calculs directs : lire des données économiques]
\begin{enumerate}
\item Taux de variation entre 1980 et 1986 :
\[ \frac{1100-950}{950}\times 100 = \frac{150}{950}\times 100 \approx +15{,}8\,\% \]
Taux de variation entre 1986 et 1994 :
\[ \frac{620-1100}{1100}\times 100 = \frac{-480}{1100}\times 100 \approx -43{,}6\,\% \]
\item La baisse du PIB par habitant entre 1986 et 1994 est d'environ \textbf{43,6\,\%} (soit une perte de 480 dollars par habitant).
\item Ces évolutions traduisent d'abord la \textbf{prospérité} de la fin de l'ère Ahidjo (croissance tirée par le pétrole et les exportations agricoles jusqu'en 1985-1986), puis la \textbf{grave crise économique} de 1986-1994 : chute des cours des matières premières (pétrole, cacao, café), récession, dévaluation du franc CFA en janvier 1994 et paupérisation de la population.
\end{enumerate}
\textbf{Point de vigilance :} le taux de variation se calcule toujours par rapport à la valeur de \textbf{départ} (diviseur = valeur initiale), jamais par rapport à la valeur d'arrivée.
\end{corrige}

\begin{corrige}[Exercice 5 — Chronologie : classer et dater]
Ordre chronologique correct :
\begin{enumerate}
\item Indépendance du Cameroun oriental : 1\textsuperscript{er} janvier 1960 ;
\item Plébiscites sous l'égide de l'ONU : 11 et 12 février 1961 ;
\item Réunification effective des deux Cameroun : 1\textsuperscript{er} octobre 1961 ;
\item Création de l'Union nationale camerounaise (parti unique) : 1\textsuperscript{er} septembre 1966 ;
\item Proclamation de la République unie du Cameroun : référendum du 20 mai 1972 (Constitution du 2 juin 1972) ;
\item Démission d'Ahmadou Ahidjo : 4 novembre 1982 ;
\item Tentative de coup d'État contre Paul Biya : 6-7 avril 1984 ;
\item Retour au multipartisme : lois de décembre 1990 ;
\item Promulgation de la Constitution : 18 janvier 1996 ;
\item Premières élections régionales : 6 décembre 2020.
\end{enumerate}
\begin{popfigure}
\begin{center}
\begin{tikzpicture}[scale=1]
\draw[->,thick,popdark] (0,0) -- (15.5,0);
\foreach \x/\y/\t in {0.3/1960/{Indép. Cam. oriental}, 2.2/1961/{Plébiscites ONU}, 4.1/1961/{Réunification}, 5.9/1966/{Parti unique UNC}, 7.9/1972/{Rép. unie}, 9.6/1982/{Démission Ahidjo}, 11.0/1984/{Coup d'État manqué}, 12.4/1990/{Multipartisme}, 13.7/1996/{Constitution}, 15.0/2020/{Élections régionales}}{
\fill[poporange] (\x,0) circle (0.07);
\draw[popmuted] (\x,0) -- (\x,0.35);
\node[above,align=center,font=\scriptsize,text=popdark] at (\x,0.35) {\y\\\t};
}
\end{tikzpicture}
\end{center}
\legende{Frise chronologique corrigée : les grandes dates du Cameroun indépendant, de 1960 à 2020.}
\end{popfigure}
\textbf{Point de vigilance :} pour les plébiscites, écrire « 11 et 12 février 1961 » (deux jours selon les zones) ; pour la République unie, on accepte 20 mai 1972 (référendum) ou 2 juin 1972 (Constitution), mais il faut préciser lequel on choisit.
\end{corrige}

\begin{corrige}[Exercice 6 — Question à réponse construite : la succession de 1982]
\textbf{Corrigé rédigé (modèle) :}

Le 4 novembre 1982, le président Ahmadou Ahidjo, au pouvoir depuis 1960, surprend le pays en annonçant sa démission, officiellement pour « fatigue ». Conformément à la Constitution de 1972, qui prévoit qu'en cas de vacance de la présidence le Premier ministre assure l'intérim puis la succession, Paul Biya, Premier ministre depuis 1975, devient président de la République le 6 novembre 1982. Chrétien du Centre-Sud, formé en France, longtemps fidèle collaborateur d'Ahidjo, il incarne alors la continuité. Mais la passation de pouvoir débouche vite sur des tensions : Ahidjo, resté président du parti unique UNC, entend conserver la réalité du pouvoir et tente de subordonner le chef de l'État au parti. En juin 1983, Biya écarte les fidèles d'Ahidjo ; celui-ci démissionne de la présidence de l'UNC en août 1983 et s'exile. Les 6 et 7 avril 1984, des officiers nordistes de la garde présidentielle tentent un coup d'État, qui échoue et est durement réprimé. Biya en sort consolidé : il fait élaborer une nouvelle Constitution (janvier 1984), rebaptise le pays « République du Cameroun » et transforme l'UNC en RDPC en 1985.

\textbf{Barème indicatif (sur 5 points) :} Constitution et légalité de la succession (1 pt) ; personnalité de Biya (1 pt) ; tensions avec Ahidjo 1983 (1 pt) ; coup d'État d'avril 1984 et consolidation (1 pt) ; clarté et dates exactes (1 pt).

\textbf{Points de vigilance :} insister sur la \emph{légalité constitutionnelle} de la succession ; ne pas écrire que Biya a été « élu » en 1982 ; dater précisément la démission (4 novembre 1982) et l'investiture (6 novembre 1982).
\end{corrige}

\begin{corrige}[Exercice 7 — Du fédéralisme à l'unité : compléter un tableau]
\begin{center}
\begin{tabularx}{\linewidth}{|l|X|X|}\hline
\rowcolor{popblueL} Critère & République fédérale (1961-1972) & République unie (1972-1984) \\ \hline
Nombre d'États fédérés & Deux : Cameroun oriental et Cameroun occidental & Aucun : État unitaire, provinces (puis régions) \\ \hline
Institutions propres aux États & Assemblées et gouvernements fédérés propres à chaque État & Institutions uniques : une Assemblée nationale, un gouvernement central \\ \hline
Symbole de l'unité & Drapeau à deux étoiles (les deux États) & Drapeau à une étoile, hymne national unique \\ \hline
Acte fondateur & Constitution fédérale du 1\textsuperscript{er} septembre 1961 (conférence de Foumban, juillet 1961) & Référendum du 20 mai 1972, Constitution du 2 juin 1972 \\ \hline
\end{tabularx}
\end{center}
\textbf{Conclusion :} le référendum du 20 mai 1972 marque la fin du fédéralisme et la naissance d'un État unitaire, présenté par Ahidjo comme plus économique et plus efficace. Il parachève l'unification nationale commencée en 1961 en effaçant la dualité institutionnelle héritée des colonisations française et britannique.

\textbf{Point de vigilance :} les deux étoiles du drapeau fédéral symbolisent les deux États fédérés ; le retour à une seule étoile en 1972 illustre l'unité.
\end{corrige}

\begin{corrige}[Exercice 8 — Commentaire de données chiffrées : la crise économique]
\begin{enumerate}
\item \textbf{Description :} entre 1985 et 1990, la croissance du PIB s'effondre, passant de $+7,0\,\%$ à $-6,0\,\%$, et reste négative en 1994 ($-3,0\,\%$) : l'économie est en récession. Les recettes pétrolières chutent fortement : l'indice passe de 100 à 55 en 1990 ($-45\,\%$) puis à 45 en 1994 ($-55\,\%$ par rapport à 1985). Les salaires de la fonction publique, maintenus jusqu'en 1990 (indice 100), sont brutalement réduits en 1994 (indice 30), soit une baisse de $70\,\%$.
\item \textbf{Causes de la crise (deux attendues) :} la chute des cours mondiaux des matières premières à partir de 1986 (contre-choc pétrolier, baisse des prix du cacao et du café), qui réduit les recettes d'exportation ; la mauvaise gestion et l'endettement de l'État, aggravés par la dépendance de l'économie aux exportations de produits de base.
\item \textbf{Conséquences sociales :} forte baisse des salaires (jusqu'à $70\,\%$ pour les fonctionnaires en 1993-1994), licenciements et chômage croissant, paupérisation des ménages, réduction des dépenses de santé et d'éducation dans le cadre des programmes d'ajustement structurel, dévaluation du franc CFA en janvier 1994 qui renchérit les importations.
\end{enumerate}
\textbf{Point de vigilance :} distinguer clairement \emph{causes} (extérieures : cours des matières premières ; internes : gestion, endettement) et \emph{conséquences} (sociales) ; appuyer chaque affirmation sur un chiffre du tableau.
\end{corrige}

\begin{corrige}[Exercice 9 — Dissertation type Baccalauréat]
\textbf{Corrigé détaillé (plan et éléments de rédaction).}

\textbf{Introduction (modèle).} \emph{Accroche :} le 1\textsuperscript{er} janvier 1960, le Cameroun oriental accède à l'indépendance ; un an plus tard, la réunification réalise le rêve des nationalistes. \emph{Présentation :} de 1960 à 1982, sous la présidence d'Ahmadou Ahidjo, le Cameroun construit un État à partir de deux territoires d'administrations différentes et connaît un développement économique remarquable. \emph{Problématique :} comment le Cameroun indépendant a-t-il réussi, entre 1960 et 1982, à édifier un État unifié et stable, et dans quelle mesure son développement économique a-t-il été réel et équilibré ? \emph{Annonce du plan :} nous étudierons d'abord la construction de l'État camerounais (I), puis le développement économique et ses limites (II).

\textbf{I. La construction de l'État camerounais (1960-1972).}
\begin{itemize}
\item Indépendances : 1\textsuperscript{er} janvier 1960 (Cameroun oriental, ex-français) ; contexte de la répression de l'UPC et de la guerre civile larvée.
\item Plébiscites des 11-12 février 1961 sous l'égide de l'ONU : le Northern Cameroons choisit le Nigeria, le Southern Cameroons choisit le Cameroun.
\item Réunification le 1\textsuperscript{er} octobre 1961 : naissance de la République fédérale du Cameroun (conférence de Foumban, juillet 1961 ; Constitution fédérale du 1\textsuperscript{er} septembre 1961).
\item Renforcement du pouvoir central : parti unique UNC (1\textsuperscript{er} septembre 1966), référendum du 20 mai 1972 et République unie du Cameroun (Constitution du 2 juin 1972) : État unitaire, présidentialiste et centralisé.
\end{itemize}

\textbf{II. Le développement économique et ses limites.}
\begin{itemize}
\item Réussites : planification (plans quinquennaux), essor des cultures d'exportation (cacao, café, coton, bananes), début de l'exploitation pétrolière (à partir de 1977-1978), industrialisation (ALUCAM, SONEL, sucreries, cimenteries), infrastructures (Transcamerounais, port de Douala), croissance forte (environ $5$ à $7\,\%$ par an) et stabilité monétaire dans la zone franc.
\item Limites : économie extravertie et dépendante des cours mondiaux et des capitaux étrangers ; déséquilibres régionaux (concentration sur Douala et le Sud, enclavement du Nord) ; faible transformation locale des matières premières ; inégalités sociales et corruption naissante.
\end{itemize}

\textbf{Conclusion (modèle).} Entre 1960 et 1982, le Cameroun est passé du statut de territoires sous tutelle à celui d'État unifié, stable et relativement prospère : le pari de la construction nationale est réussi, mais le développement reste fragile car dépendant de l'extérieur. La démission d'Ahidjo le 4 novembre 1982 ouvre une nouvelle ère, celle de Paul Biya, qui devra affronter dès 1986 une crise économique majeure.

\textbf{Barème indicatif (sur 20) :} introduction (accroche, problématique, plan) : 4 pts ; construction de l'État (faits et dates exacts) : 6 pts ; développement économique et limites : 6 pts ; conclusion et ouverture : 2 pts ; qualité de l'expression : 2 pts.

\textbf{Points de vigilance :} dater avec exactitude (1960, 1961, 1966, 1972) ; ne pas réduire le développement au pétrole (l'agriculture reste essentielle) ; toujours nuancer (« réussites \emph{mais} limites ») pour répondre à « discutez ».
\end{corrige}

\begin{corrige}[Exercice 10 — Étude critique de documents : du fédéralisme à l'unité]
\begin{enumerate}
\item \textbf{Présentation des documents.} Document 1 : extrait d'un discours politique prononcé par Ahmadou Ahidjo, président de la République du Cameroun, à la conférence de Foumban en juillet 1961, réunie pour préparer la Constitution fédérale avant la réunification du 1\textsuperscript{er} octobre 1961. Document 2 : extrait de la Constitution du 2 juin 1972, texte fondamental adopté après le référendum du 20 mai 1972, qui institue la République unie du Cameroun.
\item \textbf{Pourquoi le fédéralisme en 1961 ?} Selon le document 1, les deux territoires ont vécu « quarante ans d'administrations différentes » (française et britannique) : ils diffèrent par la langue, le droit, l'administration et les habitudes politiques. La fédération permet de « respecter les particularismes de chacun » tout en « construisant la nation camerounaise » : c'est un compromis rassurant pour la minorité anglophone.
\item \textbf{La rupture de 1972.} Le document 2 affirme que la République est « une et indivisible » et adopte le nom de « République unie » : les États fédérés disparaissent, ainsi que leurs assemblées et gouvernements propres. Le dualisme institutionnel de 1961 est aboli au profit d'un État unitaire centralisé.
\item \textbf{Les justifications d'Ahidjo.} À l'aide du cours : le coût élevé des institutions fédérales (doubles administrations jugées trop lourdes pour un pays en développement) ; la volonté d'achever l'unité nationale et de dépasser le clivage anglophone-francophone ; la recherche de l'efficacité de l'action de l'État (décisions plus rapides, développement harmonisé).
\item \textbf{Conclusion rédigée (modèle, une dizaine de lignes).} Le passage du fédéralisme à l'unité apparaît à la fois comme une rupture et comme un aboutissement. Rupture, car la Constitution de 1972 supprime les États fédérés, leurs institutions et le compromis de Foumban qui protégeait les particularismes ; le nom même du pays change. Mais c'est aussi l'aboutissement d'un processus engagé dès 1961 : la création du parti unique en 1966 avait déjà unifié la vie politique, et le pouvoir fédéral dominait largement les États fédérés. Le référendum du 20 mai 1972 parachève donc une centralisation progressive, au nom de l'unité nationale et de l'efficacité, tout en laissant ouverte, pour l'avenir, la question de la place de la minorité anglophone dans l'État camerounais.
\end{enumerate}
\textbf{Barème indicatif (sur 20) :} présentation des documents : 4 pts ; exploitation du document 1 : 3 pts ; exploitation du document 2 : 3 pts ; mobilisation des connaissances : 4 pts ; conclusion argumentée : 4 pts ; expression : 2 pts.

\textbf{Points de vigilance :} citer explicitement des passages des documents entre guillemets ; ne pas confondre la conférence de Foumban (juillet 1961) avec le référendum de 1972 ; la conclusion doit trancher la question posée (rupture \emph{ou} aboutissement, ou les deux nuancés).
\end{corrige}

\begin{corrige}[Exercice 11 — Travail pratique noté : croquis cartographique]
\textbf{Modèle de croquis attendu :}

\begin{popfigure}
\begin{center}
\begin{tikzpicture}[scale=1.05]
% Contour simplifié du Cameroun
\draw[thick,popdark,fill=popcream] (2.2,0) -- (4.6,0.2) -- (5.2,1.2) -- (4.8,2.4) -- (5.6,3.2) -- (5.0,4.4) -- (3.6,5.4) -- (2.6,6.6) -- (1.8,7.6) -- (1.2,6.4) -- (0.6,5.0) -- (1.0,3.6) -- (0.4,2.4) -- (1.2,1.2) -- cycle;
% Nord
\draw[->,thick,popdark] (6.3,6.8) -- (6.3,7.8);
\node[above,font=\small\bfseries] at (6.3,7.8) {N};
% Villes
\fill[popblue] (1.0,0.9) circle (0.09); \node[left,font=\scriptsize] at (0.95,0.9) {Douala};
\fill[popblue] (2.4,1.6) circle (0.09); \node[right,font=\scriptsize] at (2.45,1.6) {Yaoundé};
\fill[popblue] (1.9,7.0) circle (0.09); \node[right,font=\scriptsize] at (1.95,7.0) {Maroua};
\fill[popblue] (0.9,3.1) circle (0.09); \node[left,font=\scriptsize] at (0.85,3.1) {Buéa};
\fill[popblue] (1.6,3.9) circle (0.09); \node[right,font=\scriptsize] at (1.65,3.9) {Bamenda};
% Zone Boko Haram
\fill[poppink,opacity=0.35] (1.2,6.4) -- (1.8,7.6) -- (2.6,6.6) -- (2.4,5.8) -- (1.4,5.8) -- cycle;
% Zone anglophone
\fill[poporange,opacity=0.4] (0.4,2.4) -- (1.0,3.6) -- (1.9,4.2) -- (2.0,2.8) -- (1.2,2.2) -- cycle;
% Zone réfugiés Est
\fill[popgreen,opacity=0.35] (4.8,2.4) -- (5.6,3.2) -- (5.0,4.4) -- (4.2,4.0) -- (4.4,2.8) -- cycle;
% Camp de Minawao
\fill[poppurple] (2.3,6.5) circle (0.07); \node[right,font=\scriptsize] at (2.35,6.5) {Minawao};
% Corridor Douala-N'Djamena
\draw[->,very thick,popgold] (1.1,1.0) -- (2.4,2.2) -- (3.2,4.0) -- (2.6,6.0) -- (2.2,7.4);
\node[font=\scriptsize,popgold,rotate=62] at (3.5,4.9) {Corridor Douala-N'Djamena};
% Transcamerounais
\draw[very thick,poppurple,dashed] (1.1,1.0) -- (2.0,2.4) -- (2.8,3.6) -- (3.4,4.8);
\node[font=\scriptsize,poppurple,rotate=55] at (2.0,3.4) {Transcamerounais};
% Échelle
\draw[thick] (0.4,-0.5) -- (1.9,-0.5);
\draw[thick] (0.4,-0.6) -- (0.4,-0.4); \draw[thick] (1.9,-0.6) -- (1.9,-0.4);
\node[below,font=\scriptsize] at (1.15,-0.6) {0 \quad 200 km (indicatif)};
\end{tikzpicture}
\end{center}
\legende{1. Zone d'insécurité Boko Haram (rose) et camp de Minawao ; 2. Régions en crise anglophone (orange) ; 3. Zone d'accueil des réfugiés centrafricains (vert) ; 4. Villes principales (points bleus) ; 5. Corridor Douala-N'Djamena (flèche jaune) ; 6. Chemin de fer Transcamerounais (pointillés violets).}
\end{popfigure}

\textbf{Légende organisée attendue sur la copie :}
\begin{enumerate}
\item Zone d'insécurité liée à \cle{Boko Haram} (Extrême-Nord, depuis 2014) : figuré surfacique rose ; camp de déplacés de Minawao (point violet) ;
\item Régions du Nord-Ouest et du Sud-Ouest en \cle{crise anglophone} depuis 2016-2017 : figuré surfacique orange ;
\item Région de l'Est accueillant les \cle{réfugiés centrafricains} : figuré surfacique vert ;
\item Villes et pôles : points bleus (Douala, port et capitale économique ; Yaoundé, capitale politique ; Maroua, Buéa, Bamenda) ;
\item Axes de \cle{désenclavement} : corridor Douala-N'Djamena (flèche jaune), Transcamerounais (pointillés violets).
\end{enumerate}

\textbf{Conclusion rédigée (modèle, cinq lignes) :} le croquis montre que le Cameroun contemporain fait face à trois types de défis. Sur le plan sécuritaire, l'Extrême-Nord subit les attaques de Boko Haram depuis 2014 et l'Est accueille des réfugiés centrafricains. Sur le plan politique, la crise anglophone qui secoue le Nord-Ouest et le Sud-Ouest depuis 2016-2017 interroge l'unité nationale. Sur le plan économique enfin, la concentration des activités sur le port de Douala et l'enclavement des régions septentrionales imposent de grands axes de désenclavement comme le corridor Douala-N'Djamena et le Transcamerounais.

\textbf{Barème indicatif (sur 20) :} contour, titre, nord, échelle : 4 pts ; exactitude du report des informations : 6 pts ; qualité et organisation de la légende numérotée : 4 pts ; propreté et lisibilité : 2 pts ; conclusion rédigée : 4 pts.

\textbf{Points de vigilance :} le titre et la légende sont obligatoires (un croquis sans légende est fortement pénalisé) ; placer correctement les villes (Douala sur la côte, Maroua à l'Extrême-Nord, Bamenda et Buéa à l'ouest) ; ne pas confondre crise anglophone (ouest) et insécurité Boko Haram (Extrême-Nord).
\end{corrige}

\newpage

\coursec{Fiche bilan}

\begin{popfigure}
\begin{tikzpicture}[
  mindmap,
  level 1 concept/.append style={level distance=4.5cm,sibling angle=60},
  level 2 concept/.append style={level distance=3cm,sibling angle=45},
  concept color=popblue!80,
  text=popdark,
  font=\small\bfseries,
  every node/.style={scale=0.85},
  grow cyclic
]
\node[concept, minimum size=2.5cm, align=center]{Le Cameroun\\ indépendant\\(1960--nos jours)}
  child[concept color=poporange] {
    node[concept, minimum size=2cm, align=center]{1. Indépendances\\ \& Réunification}
    child { node[concept, minimum size=1.6cm] {1er janv. 1960 : Indép. Cameroun français} }
    child { node[concept, minimum size=1.6cm] {1er oct. 1961 : Réunification (Foumbam)} }
    child { node[concept, minimum size=1.6cm] {Ahmadou Ahidjo 1er Président} }
  }
  child[concept color=popgreen] {
    node[concept, minimum size=2cm, align=center]{2. Ère Ahidjo\\ (1961--1972)}
    child { node[concept, minimum size=1.6cm] {Fédéralisme (2 États fédérés)} }
    child { node[concept, minimum size=1.6cm] {Référendum 1972 : État unitaire} }
    child { node[concept, minimum size=1.6cm] {Parti unique : UNC $\to$ CNU} }
  }
  child[concept color=poppurple] {
    node[concept, minimum size=2cm, align=center]{3. Bilan Ahidjo\\ (1960--1982)}
    child { node[concept, minimum size=1.6cm] {Stabilité politique \& unité} }
    child { node[concept, minimum size=1.6cm] {Croissance éco. (pétrole, agro)} }
    child { node[concept, minimum size=1.6cm] {Éducation, santé, infrastructures} }
    child { node[concept, minimum size=1.6cm] {Limites : autoritarisme, dette} }
  }
  child[concept color=poppink] {
    node[concept, minimum size=2cm] {4. Transition 1982}
    child { node[concept, minimum size=1.6cm] {Démission Ahidjo (4 nov. 1982)} }
    child { node[concept, minimum size=1.6cm] {Paul Biya : succession constitutionnelle} }
    child { node[concept, minimum size=1.6cm] {Conflit Ahidjo/Biya (1983--84)} }
    child { node[concept, minimum size=1.6cm] {Tentative coup d'État 1984} }
  }
  child[concept color=popgold] {
    node[concept, minimum size=2cm, align=center]{5. Mutations\\ politiques (1982--...)}
    child { node[concept, minimum size=1.6cm] {Libéralisation 1990 (multipartisme)} }
    child { node[concept, minimum size=1.6cm] {Constitution 1996 (décentralisation)} }
    child { node[concept, minimum size=1.6cm] {Élections présidentielles régulières} }
    child { node[concept, minimum size=1.6cm] {Crise anglophone (depuis 2016)} }
  }
  child[concept color=popblue] {
    node[concept, minimum size=2cm, align=center]{6. Défis\\ contemporains}
    child { node[concept, minimum size=1.6cm] {Émergence 2035 (SND30)} }
    child { node[concept, minimum size=1.6cm] {Gouvernance \& corruption} }
    child { node[concept, minimum size=1.6cm] {Sécurité (Boko Haram, NOSO)} }
    child { node[concept, minimum size=1.6cm] {Jeunesse, emploi, formation} }
  };
\end{tikzpicture}
\legende{Carte mentale de synthèse : Le Cameroun indépendant (1960--nos jours).}
\end{popfigure}

\begin{bilan}
\courssub{Repères chronologiques essentiels}
\begin{center}
\begin{tabularx}{\linewidth}{|c|X|}\hline
\textbf{Date} & \textbf{Événement majeur} \\ \hline
1er janvier 1960 & Indépendance du Cameroun français (République du Cameroun) \\ \hline
1er octobre 1961 & Réunification : naissance de la République fédérale du Cameroun (Foumbam) \\ \hline
20 mai 1972 & Référendum : passage à l'État unitaire (République unie du Cameroun) \\ \hline
4 novembre 1982 & Démission d'Ahmadou Ahidjo ; Paul Biya devient Président \\ \hline
1983--1984 & Conflit Ahidjo/Biya ; tentative de coup d'État (avril 1984) \\ \hline
1990 & Loi sur les partis politiques : retour au multipartisme \\ \hline
18 janvier 1996 & Promulgation de la Constitution (décentralisation, Sénat, Cour suprême) \\ \hline
2016--... & Crise anglophone (revendications Nord-Ouest/Sud-Ouest) \\ \hline
\end{tabularx}
\end{center}

\courssub{Acteurs clés}
\begin{itemize}
\item \textbf{Ahmadou Ahidjo} (1924--1989) : 1er Président (1960--1982), artisan de l'unité nationale et de l'État unitaire.
\item \textbf{Paul Biya} (né 1933) : Premier ministre (1975--1982), Président depuis 1982 ; libéralisation politique (1990), Constitution 1996.
\item \textbf{John Ngu Foncha} : Vice-président fédéral (Cameroon Occidental), figure de la réunification.
\item \textbf{André-Marie Mbida}, \textbf{Ernest Ouandié}, \textbf{Félix-Roland Moumié} : figures du nationalisme et de l'opposition (UPC).
\end{itemize}

\courssub{Concepts \& Notions à maîtriser}
\begin{definition}[Réunification]
Processus d'union entre la République du Cameroun (ex-français) et le Southern Cameroons (ex-britannique) aboutissant à la Fédération (1er oct. 1961), consacrée par la Conférence de Foumbam (juillet 1961).
\end{definition}
\begin{definition}[État unitaire]
Régime issu du référendum du 20 mai 1972 supprimant le fédéralisme : un seul gouvernement, une seule assemblée, une seule constitution pour tout le territoire.
\end{definition}
\begin{definition}[Parti unique (CNU)]
Congrès National Uni (1966) : fusion des partis politiques en un parti d'État, instrument de mobilisation et de contrôle social jusqu'en 1990.
\end{definition}
\begin{definition}[Libéralisation politique (1990)]
Adoption des lois sur les partis politiques, les associations, la presse ; fin du parti unique ; ouverture au multipartisme.
\end{definition}
\begin{definition}[Décentralisation (Constitution 1996)]
Transfert de compétences et de ressources de l'État aux Collectivités Territoriales Décentralisées (CTD : régions, communes) pour une gestion de proximité.
\end{definition}

\courssub{Méthode express : Commentaire de document historique (sujet Bac)}
\begin{methode}[Analyse d'un discours / texte officiel]
\begin{enumerate}
\item \textbf{Identifier} : Nature (discours, loi, constitution, article de presse), Auteur, Date, Contexte précis.
\item \textbf{Analyser le contenu} : Idées forces, vocabulaire clé (unité, nation, démocratie, développement), arguments avancés.
\item \textbf{Expliquer la portée} : Pourquoi ce document à ce moment ? Rupture ou continuité ? Conséquences immédiates et à long terme.
\item \textbf{Conclure} : Apport du document à la compréhension de l'histoire politique du Cameroun (construction étatique, enjeux de pouvoir, défis actuels).
\end{enumerate}
\end{methode}

\courssub{Bilan synthétique par période}
\begin{center}
\begin{tabularx}{\linewidth}{|>{\bfseries}l|X|X|}\hline
\textbf{Période} & \textbf{Acquis majeurs} & \textbf{Limitations / Enjeux} \\ \hline
1960--1972 (Fédéralisme) & Indépendance, réunification, stabilité relative, institutions biculturelles. & Tensions fédéralistes, rébellion UPC, centralisation rampante. \\ \hline
1972--1982 (État unitaire Ahidjo) & Unité administrative, croissance économique (pétrole, « vert »), scolarisation. & Autoritarisme, parti unique, début de la dette, succession non préparée. \\ \hline
1982--1990 (Début Biya) & Continuité, rigueur morale, gestion crise 1984, début libéralisation éco. & Héritage autoritaire, conflit Ahidjo/Biya, crise économique fin années 80. \\ \hline
1990--nos jours (Multipartisme) & Libertés publiques, Constitution 1996, décentralisation, élections régulières. & Crise anglophone, insécurité (Extrême-Nord, NOSO), gouvernance, chômage jeunes, émergence 2035. \\ \hline
\end{tabularx}
\end{center}
\end{bilan}

\begin{aretenir}[Checklist avant le Bac]
\begin{itemize}
\item $\square$ Je sais dater et expliquer l'indépendance (1er janv. 1960) et la réunification (1er oct. 1961, Conférence de Foumbam).
\item $\square$ Je distingue fédéralisme (1961--1972) et État unitaire (depuis 1972) : institutions, avantages/inconvénients.
\item $\square$ Je connais le rôle d'Ahmadou Ahidjo (construction étatique, parti unique CNU, bilan éco/social).
\item $\square$ J'explique la transition de 1982 (démission Ahidjo, succession Biya, conflit 1983--84, coup d'État 1984).
\item $\square$ Je maîtrise la libéralisation de 1990 (lois, multipartisme) et la Constitution de 1996 (décentralisation, Sénat).
\item $\square$ Je suis capable d'analyser un document (discours d'Ahidjo ou Biya, extrait de Constitution, article de presse) avec la méthode en 4 étapes.
\item $\square$ J'identifie les défis actuels : crise anglophone (NOSO), Boko Haram, gouvernance, émergence 2035 (SND30), jeunesse/emploi.
\item $\square$ Je sais situer les grands axes de la politique économique : planification (plans quinquennaux), pétrole, agriculture, SND30.
\item $\square$ Je connais les symboles de l'unité : drapeau, hymne, fête nationale (20 mai), bilinguisme officiel.
\item $\square$ Je peux rédiger une dissertation structurée (intro, 2--3 parties équilibrées, conclusion) sur un thème du programme.
\item $\square$ Je révise les cartes schématiques : évolution territoriale (fédéralisme $\to$ unitaire), zones de crise (NOSO, Extrême-Nord).
\end{itemize}
\end{aretenir}

\findecours

\end{document}
